% Généralisation et dosages d'oxydoréduction — Chimie 1ère D — Cours signé OnBuch+
% Programme officiel MINESEC, Première C, D, E, TI. Compiler avec : tectonic N08-dosages-oxydoreduction.tex
\newcommand{\DOCMATIERE}{Chimie}
\newcommand{\DOCNIVEAU}{1ère D}
\newcommand{\DOCMODULE}{Module 2 : Oxydoréduction}
\newcommand{\DOCLECON}{08}
\newcommand{\DOCTITRE}{Généralisation et dosages d'oxydoréduction}
% OnBuch+ — préambule « cours signés » ENRICHI (compilé avec tectonic / XeLaTeX)
% Système visuel « Pop » d'OnBuch+ : fond crème (jamais blanc), bordures encre
% épaisses, ombres « dures » décalées, titres Archivo Black en capitales.
% Version enrichie pour les cours de chimie : chimie (mhchem, chemfig),
% graphes (pgfplots), boîtes pédagogiques supplémentaires (définition,
% propriété, méthode, attention, le savais-tu, exercice résolu, exercice,
% TP, bilan), page de garde refaite et signature OnBuch+ ancrée en bas.
%
% Macros attendues dans main.tex AVANT \input{preamble} :
%   \DOCMATIERE  \DOCNIVEAU  \DOCMODULE  \DOCLECON (numéro)  \DOCTITRE
\documentclass[11pt]{article}

\usepackage{fontspec}
\usepackage[french]{babel}
\usepackage[a4paper,margin=2cm,top=3.4cm,bottom=2.8cm,headheight=1.4cm,footskip=1.3cm]{geometry}
\usepackage{amsmath,amssymb,amsfonts}
\usepackage[table]{xcolor} % [table] : fournit aussi \rowcolors (alternance de lignes)
\usepackage{pagecolor}
\usepackage{tikz}
\usetikzlibrary{arrows.meta,positioning,calc,shapes.geometric,shapes.misc,shapes.arrows,decorations.pathmorphing,decorations.markings,patterns,shadows,mindmap,trees,fit,backgrounds,matrix,intersections}
\usepackage{pgfplots}
\pgfplotsset{compat=1.18}
\usepackage[version=4]{mhchem}
\usepackage{chemfig}
\usepackage{enumitem}
\usepackage{fancyhdr}
% [most] charge toutes les bibliothèques tcolorbox (listings, minted, raster,
% documentation...) et gonfle inutilement la mémoire TeX sur un document long
% (~300 boîtes) : on ne charge que ce qui est réellement utilisé.
\usepackage[skins,breakable]{tcolorbox}
\usepackage{graphicx}
\usepackage{colortbl}
\usepackage{booktabs}
\usepackage{tabularx}
\usepackage{array}
\usepackage{multicol}
\usepackage{siunitx}
% parse-numbers=false : accepte aussi \SI{2,5 \times 10^{-3}}{...}
\sisetup{locale=FR,output-decimal-marker={,},group-minimum-digits=4,parse-numbers=false}
\DeclareSIUnit\ohm{\ensuremath{\symup{\Omega}}} % U+2126 absent de la police
\usepackage{qrcode}
% \degree : commande fréquemment utilisée par les modèles pour un angle ou une
% température en dehors de \SI{}{} (non fournie par les paquets chargés).
\providecommand{\degree}{^{\circ}}
\usepackage{lastpage}
\usepackage{hyperref}
\hypersetup{hidelinks}

% ============================================================
% Palette « Pop » OnBuch+ — mêmes hex que la classe P (plus_ui.dart)
% ============================================================
\definecolor{poporange}{HTML}{F59321}
\definecolor{popdark}{HTML}{E07A0C}
\definecolor{popcream}{HTML}{FFF6E8}
\definecolor{popink}{HTML}{1C1714}
\definecolor{poppurple}{HTML}{7A5AE0}
\definecolor{popgreen}{HTML}{2E9E6A}
\definecolor{poppink}{HTML}{E0457B}
\definecolor{popblue}{HTML}{2F7DE1}
\definecolor{popgold}{HTML}{C98A12}
\definecolor{popmuted}{HTML}{8A7F76}
\definecolor{popline}{HTML}{EADFD2}
\definecolor{popgray}{HTML}{8A7F76}
\colorlet{popgrey}{popgray}
% Alias tolérants (le modèle nomme parfois les couleurs différemment)
\colorlet{popwhite}{white}
\colorlet{popblack}{popink}
\colorlet{popred}{poppink}
\colorlet{popyellow}{popgold}
% Teintes claires pour fonds de tableaux / zones
\colorlet{popblueL}{popblue!10!white}
\colorlet{poporangeL}{poporange!14!white}
\colorlet{popgreenL}{popgreen!12!white}
\colorlet{poppurpleL}{poppurple!12!white}
\colorlet{poppinkL}{poppink!12!white}
\colorlet{popgoldL}{popgold!14!white}

\pagecolor{popcream}

% ============================================================
% Typographie
% ============================================================
\defaultfontfeatures{Ligatures=TeX}
\setmainfont{PlusJakartaSans-Medium.ttf}[Path=../../../fonts/,BoldFont=PlusJakartaSans-Bold.ttf]
% Maths en sans-serif (Fira Math) avec chiffres et lettres droites de Plus
% Jakarta, pour que les formules se fondent dans le texte.
\usepackage[math-style=ISO,bold-style=ISO]{unicode-math}
\setmathfont{FiraMath-Regular.otf}[Path=../../../fonts/]
\setmathfont{PlusJakartaSans-Medium.ttf}[Path=../../../fonts/,range={up/num,up/{Latin,latin},\mathup}]
\usepackage{icomma} % virgule décimale sans espace en mode math
% Caractères absents de la police de texte (sinon carré vide) : rendus par
% leur équivalent mathématique, en texte comme en formule.
\usepackage{newunicodechar}
\newunicodechar{α}{\ensuremath{\alpha}}
\newunicodechar{Α}{\ensuremath{\alpha}}
\newunicodechar{β}{\ensuremath{\beta}}
\newunicodechar{δ}{\ensuremath{\delta}}
\newunicodechar{✓}{\popcheck{popgreen}}
\newfontfamily\popdisplay{ArchivoBlack-Regular.ttf}[Path=../../../fonts/]
\newfontfamily\poplogo{SpaceGrotesk-Bold.ttf}[Path=../../../fonts/]
\newfontfamily\popsemi{PlusJakartaSans-SemiBold.ttf}[Path=../../../fonts/]
\newfontfamily\popxbold{PlusJakartaSans-ExtraBold.ttf}[Path=../../../fonts/]
\newfontfamily\popxboldls{PlusJakartaSans-ExtraBold.ttf}[Path=../../../fonts/,LetterSpace=3.0]

\setlist[enumerate]{leftmargin=1.7em,itemsep=5pt,topsep=4pt}
\setlist[itemize]{leftmargin=1.7em,itemsep=4pt,topsep=4pt,label={\color{poporange}\textbullet}}
\setlength{\parindent}{0pt}
\setlength{\parskip}{5pt}
\renewcommand{\arraystretch}{1.25}

% chemfig : liaisons un peu plus épaisses, encre Pop
\setchemfig{atom sep=2em,bond style={line width=0.9pt,popink},cram width=3pt}

% pgfplots : style Pop par défaut
\pgfplotsset{
  every axis/.append style={axis line style={popink,line width=1pt},tick style={popink},
    grid style={popline,line width=0.5pt},label style={font=\small},tick label style={font=\footnotesize}},
  popcurve/.style={poporange,line width=1.8pt,no markers,smooth},
}

% ============================================================
% Pill / squiggle
% ============================================================
\newcommand{\poppill}[3]{%
  \tikz[baseline=-0.55ex]\node[draw=popink,line width=1.2pt,rounded corners=2.6pt,%
    fill=#2,inner xsep=6.5pt,inner ysep=3pt]{%
    \popxboldls\fontsize{7}{7}\selectfont\color{#3}\MakeUppercase{#1}};%
}
\newcommand{\popsquiggle}[1]{%
  \tikz[baseline=-0.4ex]{
    \draw[#1,line width=2.6pt,line cap=round]
      plot [smooth] coordinates {(0,0.09) (0.34,0.01) (0.68,0.21) (1.02,0.01) (1.36,0.10)};
  }%
}
% Mot-clé surligné (marqueur orange sous le texte)
\newcommand{\cle}[1]{{\popxbold\color{popdark}#1}}

% ============================================================
% En-tête / pied
% ============================================================
\pagestyle{fancy}
\fancyhf{}
\renewcommand{\headrulewidth}{0pt}
\renewcommand{\footrulewidth}{0pt}
\lhead{\raisebox{-0.15ex}{\poplogo\fontsize{13}{13}\selectfont\color{popink}OnBuch\color{poporange}+}}
\rhead{\poppill{\DOCMATIERE\ \textbullet\ \DOCNIVEAU\ \textbullet\ Leçon \DOCLECON}{white}{popink}}
\lfoot{\popxbold\fontsize{7.6}{7.6}\selectfont\color{popmuted}\MakeUppercase{Cours signé OnBuch+}\ \textbullet\ {\popsemi\DOCTITRE}}
\rfoot{\tikz[baseline=-0.6ex]\node[rounded corners=8.5pt,fill=popink,minimum height=17pt,minimum width=17pt,inner xsep=5pt,inner sep=0]{\popxbold\fontsize{8}{8}\selectfont\color{popcream}\thepage\,/\,\pageref*{LastPage}};}

% ============================================================
% Titres
% ============================================================
\newcommand{\chaptitre}[2]{%
  \begin{tcolorbox}[enhanced,colback=poporange,colframe=popink,boxrule=2.6pt,arc=11pt,
      drop shadow=popink,left=15pt,right=15pt,top=12pt,bottom=11pt,before skip=3pt,after skip=18pt]
    \poppill{#2}{popink}{popcream}\par\vspace{9pt}
    {\raggedright\hyphenpenalty=10000\exhyphenpenalty=10000\popdisplay\fontsize{22}{24}\selectfont\color{popcream}\MakeUppercase{#1}\par}
  \end{tcolorbox}
}
% Section numérotée : pastille orange avec le numéro + titre Archivo
\newcounter{popsec}
\newcommand{\coursec}[1]{%
  \stepcounter{popsec}\setcounter{popsub}{0}%
  \par\vspace{16pt}\noindent
  \begin{minipage}{\linewidth}
  \tikz[baseline=-0.7ex]\node[draw=popink,line width=1.6pt,fill=poporange,rounded corners=4pt,minimum size=22pt,inner sep=0]{\popdisplay\fontsize{12}{12}\selectfont\color{popcream}\thepopsec};%
  \hspace{8pt}{\popdisplay\fontsize{15}{17}\selectfont\color{popink}\MakeUppercase{#1}}\par
  \vspace{4pt}\noindent{\color{poporange}\rule{36pt}{3.6pt}}
  \end{minipage}\par\nopagebreak\vspace{8pt}
}
\newcounter{popsub}
\newcommand{\courssub}[1]{%
  \stepcounter{popsub}%
  \par\vspace{9pt}\noindent{\popxbold\fontsize{11.5}{13}\selectfont\color{popink}{\color{poporange}\thepopsec.\thepopsub}\hspace{6pt}#1}\par\nopagebreak\vspace{4pt}
}

% ============================================================
% Boîtes pédagogiques — même recette : fond blanc, bordure épaisse colorée,
% ombre dure encre, badge plein. Titre optionnel [#1].
% ============================================================
\tcbset{popbase/.style={breakable,enhanced,colback=white,boxrule=2.3pt,arc=9pt,
  shadow={3pt}{-3pt}{0pt}{popink},left=14pt,right=14pt,top=11pt,bottom=11pt,before skip=11pt,after skip=15pt,
  fontupper=\popsemi\fontsize{10.6}{14.5}\selectfont\color{popink},
  fontlower=\popsemi\fontsize{10.6}{14.5}\selectfont\color{popink}}}
% Coche dessinée (le glyphe ✓ n'existe pas dans les polices du document)
\newcommand{\popcheck}[1]{\tikz[baseline=-0.5ex]\draw[#1,line width=1.6pt,line cap=round,line join=round](-0.13,0.0)--(-0.04,-0.09)--(0.13,0.1);}
\newcommand{\popboxhead}[3]{\poppill{#1}{#2}{white}\ifx\relax#3\relax\else\hspace{6pt}{\popxbold\fontsize{10.6}{12}\selectfont\color{popink}#3}\fi\par\vspace{7pt}}

\newtcolorbox{aretenir}[1][]{popbase,colframe=popgreen,before upper={\popboxhead{À retenir}{popgreen}{#1}}}
\newtcolorbox{exemplebox}[1][]{popbase,colframe=poporange,before upper={\popboxhead{Exemple}{poporange}{#1}}}
\newtcolorbox{definition}[1][]{popbase,colframe=popblue,before upper={\popboxhead{Définition}{popblue}{#1}}}
\newtcolorbox{propriete}[1][]{popbase,colframe=poppurple,before upper={\popboxhead{Propriété}{poppurple}{#1}}}
\newtcolorbox{methode}[1][]{popbase,colframe=popgold,before upper={\popboxhead{Méthode}{popgold}{#1}}}
\newtcolorbox{attention}[1][]{popbase,colframe=poppink,before upper={\popboxhead{Attention}{poppink}{#1}}}
\newtcolorbox{experience}[1][]{popbase,colframe=popblue,colback=popblueL,before upper={\popboxhead{Expérience}{popblue}{#1}}}
% « Le savais-tu ? » : carte violette pleine (contexte, culture, Cameroun)
\newtcolorbox{savaistu}[1][]{popbase,colback=poppurple,colframe=popink,
  fontupper=\popsemi\fontsize{10.4}{14.2}\selectfont\color{white},
  before upper={\poppill{Le savais-tu\,?}{white}{poppurple}\ifx\relax#1\relax\else\hspace{6pt}{\popxbold\color{white}#1}\fi\par\vspace{7pt}}}
% Exercice résolu : énoncé en haut, \tcblower puis la solution sur fond crème
\newcounter{popexo}\newcounter{popexr}
\newtcolorbox{exoresolu}[1][]{popbase,colframe=poporange,colbacklower=popcream,
  segmentation style={popink,line width=1.2pt,dashed},
  before upper={\stepcounter{popexr}\popboxhead{Exercice résolu \thepopexr}{poporange}{#1}},
  before lower={\poppill{Solution}{popink}{popcream}\par\vspace{6pt}}}
% Exercice (non résolu en place) — la correction est dans la partie Corrigés
\newtcolorbox{exercice}[1][]{popbase,colframe=popink,
  before upper={\stepcounter{popexo}\popboxhead{Exercice \thepopexo}{popink}{#1}}}
% Corrigé
\newtcolorbox{corrige}[1][]{popbase,colframe=popgreen,colback=popgreenL,
  before upper={\popboxhead{Corrigé}{popgreen}{#1}}}
% Objectifs / prérequis / bilan
\newtcolorbox{objectifs}{popbase,colframe=popink,colback=poporangeL,
  before upper={\popboxhead{Objectifs de la leçon}{popink}{}}}
\newtcolorbox{prerequis}{popbase,colframe=popink,colback=popblueL,
  before upper={\popboxhead{Prérequis}{popblue}{}}}
\newtcolorbox{bilan}{popbase,colframe=popink,colback=white,boxrule=2.8pt,
  before upper={\popboxhead{Fiche bilan}{poporange}{L'essentiel à connaître pour l'examen}}}
% Figure encadrée avec légende
\newtcolorbox{popfigure}[1][]{enhanced,breakable=false,colback=white,colframe=popline,boxrule=1.4pt,arc=8pt,
  left=10pt,right=10pt,top=10pt,bottom=8pt,before skip=10pt,after skip=12pt,halign=center,
  lower separated=false,halign lower=center,
  fontlower=\popsemi\fontsize{9}{11.5}\selectfont\color{popmuted},
  #1}
\newcommand{\legende}[1]{\par\vspace{6pt}{\popsemi\fontsize{9}{11.5}\selectfont\color{popmuted}#1}}

% ============================================================
% Signature OnBuch+
% ============================================================
\newcommand{\poplogoinline}[1]{{\poplogo\fontsize{#1}{#1}\selectfont\color{popink}OnBuch\color{poporange}+}}
\newcommand{\signatureonbuch}{%
  \begin{tcolorbox}[enhanced,colback=popink,colframe=popink,boxrule=2.2pt,arc=10pt,
      drop shadow=poporange,left=14pt,right=14pt,top=10pt,bottom=10pt,before skip=0pt,after skip=0pt]
    \tikz[baseline=-0.7ex]\node[circle,fill=popgreen,draw=popcream,line width=1.3pt,minimum size=22pt,inner sep=0]{\popcheck{white}};%
    \hspace{9pt}\begin{minipage}[c]{0.62\linewidth}
      {\popxbold\fontsize{12}{13}\selectfont\color{popcream}Signé\ \textbullet\ {\poplogo OnBuch\color{poporange}+}}\par\vspace{2pt}
      {\raggedright\popsemi\fontsize{8}{10.5}\selectfont\color{popline}\MakeUppercase{Équipe pédagogique OnBuch+}\\\MakeUppercase{Conforme au programme officiel MINESEC}\par}
    \end{minipage}\hfill
    {\poplogo\fontsize{22}{22}\selectfont\color{popcream}OnBuch\color{poporange}+}
  \end{tcolorbox}%
}

% ============================================================
% Page de garde
% #1 titre · #2 matière · #3 niveau · #4 module · #5 n° leçon · #6 durée ·
% #7 description / objectifs (texte ou itemize)
% ============================================================
\newcommand{\pagegarde}[7]{%
  \thispagestyle{empty}
  \newgeometry{a4paper,margin=1.7cm,top=1.6cm,bottom=1.4cm}%
  \begin{tikzpicture}[remember picture,overlay]
    \fill[poporange!18] ([xshift=-3.2cm,yshift=-3.6cm]current page.north east) circle (4.6cm);
    \fill[poppurple!14] ([xshift=2.4cm,yshift=7.6cm]current page.south west) circle (3.8cm);
  \end{tikzpicture}%
  {\poplogo\fontsize{30}{30}\selectfont\color{popink}OnBuch\color{poporange}+}\hfill
  \raisebox{0.6ex}{\poppill{Cours signé \textbullet\ Premium}{popink}{popcream}}\par
  \vspace{1pt}\popsquiggle{poppurple}\par
  \vspace{8pt}
  \begin{tcolorbox}[enhanced,colback=poporange,colframe=popink,boxrule=2.8pt,arc=13pt,
      drop shadow=popink,left=18pt,right=18pt,top=12pt,bottom=13pt,after skip=10pt]
    \poppill{#2 \textbullet\ #3}{popink}{popcream}\hspace{5pt}\poppill{#4}{white}{popink}\par\vspace{8pt}
    {\popdisplay\fontsize{14}{14}\selectfont\color{popink}LEÇON #5}\par\vspace{4pt}
    {\raggedright\hyphenpenalty=10000\exhyphenpenalty=10000\popdisplay\fontsize{25}{27}\selectfont\color{popcream}\MakeUppercase{#1}\par}
    \vspace{8pt}
    {\popxbold\fontsize{9.5}{9.5}\selectfont\color{popink}\MakeUppercase{Durée conseillée : #6}}
  \end{tcolorbox}
  \begin{tcolorbox}[enhanced,colback=white,colframe=popink,boxrule=2.2pt,arc=10pt,
      drop shadow=popink,left=16pt,right=16pt,top=10pt,bottom=10pt,after skip=8pt]
    \poppill{Dans cette leçon}{poppurple}{white}\par\vspace{5pt}
    \popsemi\fontsize{9.3}{12.6}\selectfont\color{popink}#7
  \end{tcolorbox}
  \noindent\begin{minipage}[t]{0.487\linewidth}
    \begin{tcolorbox}[enhanced,colback=popgreen,colframe=popink,boxrule=2.2pt,arc=10pt,
        drop shadow=popink,left=11pt,right=11pt,top=10pt,bottom=10pt,height=3.1cm,valign=center]
      \tcbox[colback=white,colframe=popink,boxrule=1.3pt,arc=5pt,boxsep=0pt,nobeforeafter,
        left=3pt,right=3pt,top=3pt,bottom=3pt,valign=center]{\qrcode[height=1.9cm]{https://play.google.com/store/apps/details?id=com.luvvix.onbuch&hl=fr}}%
      \hspace{8pt}\begin{minipage}[c]{0.5\linewidth}
        {\popdisplay\fontsize{11}{12.5}\selectfont\color{white}\MakeUppercase{Télécharge OnBuch}}\par\vspace{4pt}
        {\raggedright\popsemi\fontsize{8.4}{11}\selectfont\color{white}Gratuit \textbullet\ Google Play\\Tuteur IA, annales, concours, résultats\par}
      \end{minipage}
    \end{tcolorbox}
  \end{minipage}\hfill
  \begin{minipage}[t]{0.487\linewidth}
    \begin{tcolorbox}[enhanced,colback=poppurple,colframe=popink,boxrule=2.2pt,arc=10pt,
        drop shadow=popink,left=11pt,right=11pt,top=10pt,bottom=10pt,height=3.1cm,valign=center]
      \tikz[baseline=-0.7ex]\node[circle,fill=white,draw=popink,line width=1.3pt,minimum size=1.6cm,inner sep=0]{\popdisplay\fontsize{24}{24}\selectfont\color{poppurple}+};%
      \hspace{8pt}\begin{minipage}[c]{0.6\linewidth}
        {\popdisplay\fontsize{11}{12.5}\selectfont\color{white}\MakeUppercase{Passe à OnBuch+}}\par\vspace{4pt}
        {\raggedright\popsemi\fontsize{8.4}{11}\selectfont\color{white}Tous les cours signés, Tuteur IA illimité, fiches PDF — onglet OnBuch+ de l'app\par}
      \end{minipage}
    \end{tcolorbox}
  \end{minipage}
  \vspace{8pt}
  \popcontenu
  \vfill
  \signatureonbuch
  \restoregeometry
  \newpage
}

% Rangée « Ce cours contient » (page de garde)
\newcommand{\poptile}[3]{% #1 couleur #2 gros texte #3 légende
  \begin{tcolorbox}[enhanced,colback=#1,colframe=popink,boxrule=2pt,arc=9pt,shadow={2.5pt}{-2.5pt}{0pt}{popink},
      left=6pt,right=6pt,top=9pt,bottom=9pt,halign=center,valign=center,height=1.75cm,width=\linewidth,nobeforeafter]
    {\popdisplay\fontsize{12.5}{13}\selectfont\color{white}\MakeUppercase{#2}}\par\vspace{3pt}
    {\centering\popsemi\fontsize{7.8}{9.5}\selectfont\color{white}#3\par}
  \end{tcolorbox}}
\newcommand{\popcontenu}{%
  \par\noindent{\popxboldls\fontsize{7.5}{7.5}\selectfont\color{popmuted}CE COURS SIGNÉ CONTIENT}\par\vspace{7pt}
  \noindent\begin{minipage}[t]{0.235\linewidth}\poptile{popblue}{Cours}{complet, illustré, pas à pas}\end{minipage}\hfill
  \begin{minipage}[t]{0.235\linewidth}\poptile{poporange}{Méthodes}{et exercices résolus}\end{minipage}\hfill
  \begin{minipage}[t]{0.235\linewidth}\poptile{poppink}{Exercices}{type examen + corrigés}\end{minipage}\hfill
  \begin{minipage}[t]{0.235\linewidth}\poptile{popgreen}{Fiche bilan}{l'essentiel à retenir}\end{minipage}\par}

% Sommaire
\newtcolorbox{sommaire}{enhanced,colback=white,colframe=popink,boxrule=2.2pt,arc=10pt,
  drop shadow=popink,left=18pt,right=18pt,top=13pt,bottom=13pt,before skip=6pt,after skip=20pt,
  before upper={\poppill{Sommaire}{poppurple}{white}\par\vspace{9pt}\popsemi\fontsize{10.6}{14.5}\selectfont\color{popink}}}

% Signature de fin de cours — poussée tout en bas de la dernière page
\newcommand{\findecours}{%
  \par\vfill\nopagebreak
  \begin{center}\popsquiggle{poporange}\end{center}\vspace{4pt}
  \signatureonbuch
}
% Glyphes absents de Fira Math : redéfinis à partir de glyphes disponibles
\AtBeginDocument{%
  \def\setminus{\mathbin{\backslash}}\let\smallsetminus\setminus
  \def\vdots{\mathord{\vbox{\baselineskip3.2pt\lineskiplimit0pt\kern4pt\hbox{.}\hbox{.}\hbox{.}}}}%
  \def\ddots{\mathinner{\mkern1mu\raise6.5pt\hbox{.}\mkern2mu\raise3.6pt\hbox{.}\mkern2mu\raise0.7pt\hbox{.}\mkern1mu}}%
  \def\triangle{\mathord{\scalebox{0.9}{$\symup{\Delta}$}}}%
  \def\nabla{\mathord{\raisebox{\depth}{\scalebox{1}[-1]{$\symup{\Delta}$}}}}%
  \def\checkmark{\popcheck{popdark}}%
  \def\star{\mathbin{\ast}}\def\bigstar{\mathord{\ast}}%
  \def\diamond{\mathbin{\scalebox{0.6}{$\Diamond$}}}%
  \def\bigcup{\mathop{\mathchoice{\raisebox{-0.3em}{\scalebox{1.7}{$\cup$}}}{\scalebox{1.25}{$\cup$}}{\cup}{\cup}}\displaylimits}%
  \def\bigcap{\mathop{\mathchoice{\raisebox{-0.3em}{\scalebox{1.7}{$\cap$}}}{\scalebox{1.25}{$\cap$}}{\cap}{\cap}}\displaylimits}%
  \def\lesseqqgtr{\mathrel{\lessgtr}}\def\bigoplus{\mathop{\oplus}\displaylimits}%
}
\providecommand{\ssi}{si et seulement si} % abréviation fréquente

%% FIGFIX-BEGIN v4 — correctifs globaux des figures (docs/onbuch-generation/tools/figfix)
\usepackage{adjustbox}\usepackage{etoolbox}
% toute figure TikZ/pgfplots plus large que la ligne est réduite à la largeur disponible
\BeforeBeginEnvironment{tikzpicture}{\begin{adjustbox}{max width=\linewidth}}
\AfterEndEnvironment{tikzpicture}{\end{adjustbox}}
% légendes pgfplots lisibles par-dessus les courbes
\pgfplotsset{every axis/.append style={legend style={fill=white,fill opacity=0.92,text opacity=1,draw=black!15,inner sep=3pt}}}
% étiquettes d'arêtes (syntaxe edge["..."]) avec fond blanc
\tikzset{every edge quotes/.append style={fill=white,inner sep=1.2pt}}
%% FIGFIX-END
\begin{document}

\pagegarde{Généralisation et dosages d'oxydoréduction}{Chimie}{1ère D}{Module 2 : Oxydoréduction}{08}{4 h}{Cette leçon complète la classification électrochimique avec les couples oxydant/réducteur usuels en solution aqueuse (permanganate, dichromate, diiode, thiosulfate...) et introduit le principe des dosages d'oxydoréduction. Elle prépare aux épreuves de chimie du Probatoire et aux travaux pratiques de Terminale.
\vspace{4pt}
\begin{multicols}{2}\small
\begin{itemize}
\item À la fin de cette leçon, je sais écrire les demi-équations électroniques des couples MnO4-/Mn2+, Cr2O7 2-/Cr3+, NO3-/NO, SO4 2-/SO2, Fe3+/Fe2+, Cl2/Cl-, S4O6 2-/S2O3 2- et I2/I-
\item À la fin de cette leçon, je sais établir l'équation-bilan d'une réaction d'oxydoréduction à partir des demi-équations électroniques
\item À la fin de cette leçon, je sais prévoir le sens d'une réaction d'oxydoréduction à l'aide de la classification électrochimique (règle du gamma)
\item À la fin de cette leçon, je sais définir un dosage d'oxydoréduction et citer les conditions qu'une réaction de dosage doit remplir
\item À la fin de cette leçon, je sais schématiser et décrire le dispositif expérimental d'un dosage d'oxydoréduction
\item À la fin de cette leçon, je sais repérer le point d'équivalence d'un dosage colorimétrique (changement de teinte persistant)
\end{itemize}\end{multicols}}

\begin{sommaire}
\begin{multicols}{2}
\begin{enumerate}
\item Rappels et généralisation de la classification électrochimique
\item Équations-bilan d'oxydoréduction avec les couples au programme
\item Principe général d'un dosage d'oxydoréduction
\item Dosage des ions Fe2+ par les ions MnO4- (dosage manganimétrique)
\item Dosage du diiode par les ions thiosulfate (iodométrie)
\item Méthodologie, exploitation des résultats et entraînement type Probatoire
\item Travaux pratiques
\item Méthodes et exercices résolus
\item Exercices
\item Corrigés des exercices
\item Fiche bilan
\end{enumerate}
\end{multicols}
\end{sommaire}

\begin{objectifs}
\begin{itemize}
\item À la fin de cette leçon, je sais écrire les demi-équations électroniques des couples MnO4-/Mn2+, Cr2O7 2-/Cr3+, NO3-/NO, SO4 2-/SO2, Fe3+/Fe2+, Cl2/Cl-, S4O6 2-/S2O3 2- et I2/I-
\item À la fin de cette leçon, je sais établir l'équation-bilan d'une réaction d'oxydoréduction à partir des demi-équations électroniques
\item À la fin de cette leçon, je sais prévoir le sens d'une réaction d'oxydoréduction à l'aide de la classification électrochimique (règle du gamma)
\item À la fin de cette leçon, je sais définir un dosage d'oxydoréduction et citer les conditions qu'une réaction de dosage doit remplir
\item À la fin de cette leçon, je sais schématiser et décrire le dispositif expérimental d'un dosage d'oxydoréduction
\item À la fin de cette leçon, je sais repérer le point d'équivalence d'un dosage colorimétrique (changement de teinte persistant)
\item À la fin de cette leçon, je sais exploiter la relation à l'équivalence pour calculer une concentration inconnue (dosage Fe2+/MnO4- et I2/S2O3 2-)
\item À la fin de cette leçon, je sais rédiger un compte rendu de dosage avec schéma légendé, tableau d'avancement simplifié et calcul de concentration
\end{itemize}
\end{objectifs}

\begin{prerequis}
\begin{itemize}
\item Notions d'oxydant, de réducteur, d'oxydation et de réduction (leçon précédente)
\item Écriture et équilibrage des demi-équations électroniques en milieu acide
\item Classification électrochimique qualitative et règle du gamma
\item Réaction totale, rapide et unique ; notion d'avancement et de réactif limitant
\item Calculs de quantités de matière : n = C × V, dilutions
\end{itemize}
\end{prerequis}

\newpage

\coursec{Rappels et généralisation de la classification électrochimique}

\courssub{Situation d'accroche : de l'éthylotest au dosage du fer}

Il est 22 h à l'entrée de Yaoundé. Un conducteur est invité à souffler dans le ballon d'un éthylotest. Quelques secondes plus tard, les cristaux contenus dans le tube virent du jaune-orangé au vert : l'air expiré contient de l'éthanol, qui a \emph{réduit} les ions dichromate $\ce{Cr2O7^2-}$ (orange) en ions chrome(III) $\ce{Cr^3+}$ (verts). Une réaction d'oxydoréduction vient de servir de détecteur.

Quelques kilomètres plus loin, au laboratoire municipal, un technicien détermine la teneur en fer d'un médicament contre l'anémie. Il verse goutte à goutte une solution violette de permanganate de potassium dans la solution contenant les ions $\ce{Fe^2+}$. Tant qu'il reste du fer(II), la teinte violette disparaît instantanément ; dès que tout le fer(II) a été oxydé, une teinte rose persistante apparaît : c'est le signal d'arrêt. Une réaction d'oxydoréduction vient de servir d'instrument de mesure.

Deux questions se posent alors :
\begin{itemize}
\item Comment \cle{prévoir à l'avance} si une réaction d'oxydoréduction aura lieu entre deux espèces données ?
\item Comment \cle{exploiter} une telle réaction pour mesurer précisément la concentration d'une solution ?
\end{itemize}

C'est tout l'objet de cette leçon : généraliser la classification électrochimique aux couples du programme, puis maîtriser les dosages d'oxydoréduction. Commençons par consolider les fondations.

\courssub{Rappels : oxydant, réducteur, couple Ox/Red}

Tout commence par un constat simple : certains atomes ou ions \emph{donnent} des électrons, d'autres \emph{captent} des électrons. Le fer métal d'une pointe rouille parce qu'il cède des électrons au dioxygène de l'air ; le permanganate décolore une solution de fer(II) parce qu'il lui arrache des électrons.

\begin{definition}[Oxydant, réducteur, couple Ox/Red]
\begin{itemize}
\item Un \cle{oxydant} est une espèce chimique (atome, molécule ou ion) capable de \textbf{capter} un ou plusieurs électrons.
\item Un \cle{réducteur} est une espèce chimique capable de \textbf{céder} un ou plusieurs électrons.
\item Un \cle{couple oxydant/réducteur}, noté Ox/Red, est constitué de deux espèces conjuguées qui se transforment l'une en l'autre par transfert d'électrons, selon la \cle{demi-équation électronique} :
\[ \text{Ox} + n\,\ce{e-} \ce{<=>} \text{Red} \]
\end{itemize}
La transformation d'un oxydant en son réducteur conjugué est une \cle{réduction} ; celle d'un réducteur en son oxydant conjugué est une \cle{oxydation}.
\end{definition}

\begin{attention}[Ne pas inverser les rôles]
C'est l'\textbf{oxydant} qui se \textbf{réduit} (il gagne des électrons), et le \textbf{réducteur} qui s'\textbf{oxyde} (il perd des électrons). Moyen mnémotechnique : « \emph{Le réducteur réduit l'autre, donc il s'oxyde lui-même} ». Une erreur d'inversion à ce niveau fausse toute la suite du raisonnement.
\end{attention}

Une réaction d'oxydoréduction met donc toujours en jeu \textbf{deux couples} : le réducteur de l'un cède des électrons à l'oxydant de l'autre. Les électrons ne circulent jamais librement en solution : ils sont échangés \emph{directement} entre les espèces, ce qui impose que le nombre d'électrons cédés soit exactement égal au nombre d'électrons captés.

\courssub{La règle du gamma : qui réagit avec qui ?}

Tous les oxydants ne se valent pas : le permanganate $\ce{MnO4-}$ est un oxydant redoutable, alors que l'ion $\ce{Fe^3+}$ est un oxydant modéré. De même, le fer métal est un réducteur bien plus fort que l'ion iodure $\ce{I-}$. Pour comparer ces forces, on utilise une grandeur mesurable : le \cle{potentiel normal} (ou potentiel standard) $E^\circ$, exprimé en volts (V), mesuré dans les conditions standard ($\SI{25}{\celsius}$, concentrations de $\SI{1}{mol.L^{-1}}$, pression de $\SI{1}{bar}$) par rapport à l'électrode normale à hydrogène prise comme référence ($E^\circ = \SI{0}{V}$).

\begin{propriete}[Classement des couples par potentiel normal]
Plus le potentiel normal $E^\circ$ d'un couple est \textbf{élevé} (algébriquement), plus son \textbf{oxydant} est fort et plus son \textbf{réducteur} est faible. Inversement, plus $E^\circ$ est faible (ou négatif), plus le réducteur du couple est fort.
\end{propriete}

\begin{propriete}[Règle du gamma]
Une réaction d'oxydoréduction \textbf{spontanée} a lieu entre l'oxydant le plus fort (couple de $E^\circ$ le plus élevé) et le réducteur le plus fort (couple de $E^\circ$ le plus faible). Sur une échelle verticale de potentiels où l'on place l'oxydant à gauche et le réducteur à droite de chaque couple, le chemin réactionnel dessine la lettre grecque $\gamma$ (gamma) : on descend de l'oxydant fort vers le réducteur fort, puis on remonte vers les produits.
\end{propriete}

Autrement dit : si $E^\circ_1 > E^\circ_2$, alors l'oxydant du couple 1 réagit spontanément avec le réducteur du couple 2 :
\[ \text{Ox}_1 + \text{Red}_2 \ce{->} \text{Red}_1 + \text{Ox}_2 \]

\courssub{Les nouveaux couples du programme en solution aqueuse}

Jusqu'ici, tu as surtout rencontré des couples métalliques du type $\ce{M^{n+}}/\ce{M}$. Le programme de Première élargit la famille à des couples dont les deux membres sont des espèces \emph{dissoutes} en solution aqueuse, souvent colorées, et dont les demi-équations font intervenir des molécules d'eau et des ions $\ce{H+}$. Les voici, classés par potentiel normal décroissant :

\begin{center}
\begin{tabularx}{\linewidth}{|l|X|l|}
\hline
\rowcolor{popblueL} \textbf{Couple Ox/Red} & \textbf{Remarques} & $E^\circ$ (V) \\
\hline
$\ce{MnO4-}/\ce{Mn^2+}$ & $\ce{MnO4-}$ violet, $\ce{Mn^2+}$ quasi incolore ; milieu acide & $+1,51$ \\
\hline
$\ce{Cl2}/\ce{Cl-}$ & dichlore, gaz jaune-vert, désinfectant & $+1,36$ \\
\hline
$\ce{Cr2O7^2-}/\ce{Cr^3+}$ & dichromate orange, chrome(III) vert ; milieu acide & $+1,33$ \\
\hline
$\ce{NO3-}/\ce{NO}$ & nitrate (engrais) / monoxyde d'azote gaz & $+0,96$ \\
\hline
$\ce{Fe^3+}/\ce{Fe^2+}$ & fer(III) jaune-rouille, fer(II) verdâtre & $+0,77$ \\
\hline
$\ce{I2}/\ce{I-}$ & diiode brun, iodure incolore & $+0,54$ \\
\hline
$\ce{SO4^2-}/\ce{SO2}$ & sulfate / dioxyde de soufre (conservateur) & $+0,17$ \\
\hline
$\ce{S4O6^2-}/\ce{S2O3^2-}$ & tétrathionate / thiosulfate, incolores & $+0,09$ \\
\hline
\end{tabularx}
\end{center}

\begin{attention}[Pièges classiques à éviter absolument]
\begin{itemize}
\item Dans le couple $\ce{Cr2O7^2-}/\ce{Cr^3+}$, l'ion dichromate contient \textbf{deux} atomes de chrome : il faudra écrire $2\,\ce{Cr^3+}$ à droite. Oublier ce facteur 2 est l'erreur la plus fréquente au Probatoire.
\item Ne pas confondre le \cle{thiosulfate} $\ce{S2O3^2-}$ (réducteur, utilisé en iodométrie) avec le \cle{tétrathionate} $\ce{S4O6^2-}$ (son oxydant conjugué) : ce sont les deux membres du \emph{même} couple.
\item Le permanganate $\ce{MnO4-}$ n'est réduit en $\ce{Mn^2+}$ (quasi incolore) qu'en \textbf{milieu acide} ; en milieu neutre, il donne $\ce{MnO2}$, précipité brun. On acidifie donc toujours avec l'acide sulfurique.
\end{itemize}
\end{attention}

\courssub{Écrire les demi-équations en milieu acide}

Pour ces couples « oxygénés », on ne peut plus se contenter d'ajouter des électrons : il faut équilibrer les atomes d'oxygène et d'hydrogène. En solution aqueuse acide, on dispose pour cela de deux outils : les molécules d'eau $\ce{H2O}$ et les ions hydrogène $\ce{H+}$.

\begin{methode}[Équilibrer une demi-équation en milieu acide : 4 étapes]
\begin{enumerate}
\item \textbf{Équilibrer les atomes autres que O et H} (par exemple Cr, Mn, S, N) en ajustant les coefficients stœchiométriques.
\item \textbf{Équilibrer les atomes d'oxygène} en ajoutant des molécules $\ce{H2O}$ du côté qui manque d'oxygène.
\item \textbf{Équilibrer les atomes d'hydrogène} en ajoutant des ions $\ce{H+}$ du côté qui manque d'hydrogène.
\item \textbf{Équilibrer les charges électriques} en ajoutant des électrons $\ce{e-}$ du côté de l'oxydant (à gauche pour une réduction).
\end{enumerate}
Vérification finale : les atomes \emph{et} les charges doivent être conservés des deux côtés.
\end{methode}

\begin{exemplebox}[Exemple guidé 1 : le couple $\ce{MnO4-}/\ce{Mn^2+}$]
Appliquons la méthode pas à pas.
\begin{enumerate}
\item Le manganèse est déjà équilibré : $\ce{MnO4-} \ce{<=>} \ce{Mn^2+}$.
\item Quatre oxygènes à gauche : on ajoute $4\,\ce{H2O}$ à droite : $\ce{MnO4-} \ce{<=>} \ce{Mn^2+} + 4 H2O$.
\item Huit hydrogènes à droite : on ajoute $8\,\ce{H+}$ à gauche : $\ce{MnO4-} + 8 H+ \ce{<=>} \ce{Mn^2+} + 4 H2O$.
\item Charges : à gauche $-1 + 8 = +7$ ; à droite $+2$. Il faut donc $5$ électrons à gauche :
\[ \ce{MnO4- + 8 H+ + 5 e- <=> Mn^2+ + 4 H2O} \]
\end{enumerate}
Commentaire : le permanganate ne peut capter ses 5 électrons qu'en présence d'ions $\ce{H+}$ en quantité suffisante, d'où l'obligation pratique d'acidifier le milieu (avec $\ce{H2SO4}$, jamais avec $\ce{HCl}$ dont les ions $\ce{Cl-}$ seraient oxydés en dichlore).
\end{exemplebox}

\begin{exemplebox}[Exemple guidé 2 : le couple $\ce{Cr2O7^2-}/\ce{Cr^3+}$]
\begin{enumerate}
\item Deux chrome à gauche : $\ce{Cr2O7^2-} \ce{<=>} 2 Cr^3+$.
\item Sept oxygènes : $\ce{Cr2O7^2-} \ce{<=>} 2 Cr^3+ + 7 H2O$.
\item Quatorze hydrogènes : $\ce{Cr2O7^2-} + 14 H+ \ce{<=>} 2 Cr^3+ + 7 H2O$.
\item Charges : à gauche $-2 + 14 = +12$ ; à droite $+6$. Il faut 6 électrons à gauche :
\[ \ce{Cr2O7^2- + 14 H+ + 6 e- <=> 2 Cr^3+ + 7 H2O} \]
\end{enumerate}
C'est exactement cette demi-équation qui fonctionne dans l'éthylotest : le dichromate orange est réduit en ions $\ce{Cr^3+}$ verts.
\end{exemplebox}

\begin{exemplebox}[Exemple guidé 3 : couples sans oxygène à équilibrer]
Certains couples ne demandent que l'étape des charges :
\begin{itemize}
\item Diiode/iodure : $\ce{I2 + 2 e- <=> 2 I-}$ (attention au coefficient 2 devant $\ce{I-}$).
\item Tétrathionate/thiosulfate : $\ce{S4O6^2- + 2 e- <=> 2 S2O3^2-}$ (vérification des charges : à gauche $-2 - 2 = -4$ ; à droite $2 \times (-2) = -4$).
\item Dichlore/chlorure, écrit dans le sens de la réduction : $\ce{Cl2 + 2 e- <=> 2 Cl-}$.
\end{itemize}
Pour les couples $\ce{NO3-}/\ce{NO}$ et $\ce{SO4^2-}/\ce{SO2}$, la méthode en 4 étapes donne :
\[ \ce{NO3- + 4 H+ + 3 e- <=> NO + 2 H2O} \]
\[ \ce{SO4^2- + 4 H+ + 2 e- <=> SO2 + 2 H2O} \]
Vérifie chacune de ces demi-équations toi-même, étape par étape : c'est le meilleur entraînement.
\end{exemplebox}

\courssub{L'échelle des potentiels normaux : la carte des forces}

Rassemblons tous ces couples sur une même échelle verticale, les potentiels croissant vers le haut. À gauche de chaque couple, l'oxydant ; à droite, le réducteur conjugué.

\begin{popfigure}
\begin{center}
\begin{tikzpicture}[scale=1.0]
% Echelle : y = 0.5 + 3*E
% S4O6/S2O3 : 0.09 -> 0.77 ; SO4/SO2 : 0.17 -> 1.01 ; I2/I- : 0.54 -> 2.12 ; Fe3+/Fe2+ : 0.77 -> 2.81 ; NO3/NO : 0.96 -> 3.38 ; Cr2O7/Cr3+ : 1.33 -> 4.49 ; Cl2/Cl- : 1.36 -> 4.58 ; MnO4/Mn2+ : 1.51 -> 5.03
% Axe vertical
\draw[->, thick, popdark] (0,-0.4) -- (0,6.2) node[above] {$E^\circ$ (V)};
% Graduations
\foreach \y/\lab in {0.5/{0,0}, 2.0/{0,5}, 3.5/{1,0}, 5.0/{1,5}} {
  \draw[popdark] (-0.12,\y) -- (0.12,\y) node[left=2pt] {\small \lab};
}
\draw[popline, dashed] (-3.6,0.77) -- (3.6,0.77);
\node[popgreen] at (-2.4,1.05) {\small $\ce{S4O6^2-}$};
\node[popgreen] at (2.4,1.05) {\small $\ce{S2O3^2-}$};
\node[popmuted, right] at (3.7,0.77) {\scriptsize $+0,09$};

\draw[popline, dashed] (-3.6,1.01) -- (3.6,1.01);
\node[popgold] at (-2.4,0.72) {\small $\ce{SO4^2-}$};
\node[popgold] at (2.4,0.72) {\small $\ce{SO2}$};
\node[popmuted, right] at (3.7,1.01) {\scriptsize $+0,17$};

\draw[popline, dashed] (-3.6,2.12) -- (3.6,2.12);
\node[poporange] at (-2.4,2.12) {\small $\ce{I2}$};
\node[poporange] at (2.4,2.12) {\small $\ce{I-}$};
\node[popmuted, right] at (3.7,2.12) {\scriptsize $+0,54$};

\draw[popline, dashed] (-3.6,2.81) -- (3.6,2.81);
\node[popblue] at (-2.4,2.81) {\small $\ce{Fe^3+}$};
\node[popblue] at (2.4,2.81) {\small $\ce{Fe^2+}$};
\node[popmuted, right] at (3.7,2.81) {\scriptsize $+0,77$};

\draw[popline, dashed] (-3.6,3.38) -- (3.6,3.38);
\node[poppurple] at (-2.4,3.38) {\small $\ce{NO3-}$};
\node[poppurple] at (2.4,3.38) {\small $\ce{NO}$};
\node[popmuted, right] at (3.7,3.38) {\scriptsize $+0,96$};

\draw[popline, dashed] (-3.6,4.49) -- (3.6,4.49);
\node[popink] at (-2.4,4.28) {\small $\ce{Cr2O7^2-}$};
\node[popink] at (2.4,4.28) {\small $\ce{Cr^3+}$};
\node[popmuted, right] at (3.7,4.49) {\scriptsize $+1,33$};

\draw[popline, dashed] (-3.6,4.58) -- (3.6,4.58);
\node[popgreen!60!black] at (-2.4,4.82) {\small $\ce{Cl2}$};
\node[popgreen!60!black] at (2.4,4.82) {\small $\ce{Cl-}$};
\node[popmuted, right] at (3.7,4.58) {\scriptsize $+1,36$};

\draw[popline, dashed] (-3.6,5.03) -- (3.6,5.03);
\node[poppink] at (-2.4,5.32) {\small $\ce{MnO4-}$};
\node[poppink] at (2.4,5.32) {\small $\ce{Mn^2+}$};
\node[popmuted, right] at (3.7,5.03) {\scriptsize $+1,51$};

% Flèches de force
\draw[->, very thick, popink] (-4.6,5.6) -- (-4.6,0.7) node[midway, left, align=center] {\small oxydants\\ \small de plus\\ \small en plus forts};
\draw[->, very thick, popblue] (4.9,0.7) -- (4.9,5.6) node[midway, right, align=center] {\small réducteurs\\ \small de plus\\ \small en plus forts};
\end{tikzpicture}
\end{center}
\legende{Échelle des potentiels normaux des couples du programme (valeurs indicatives à \SI{25}{\celsius}). En haut à gauche : les oxydants les plus forts ($\ce{MnO4-}$, $\ce{Cl2}$, $\ce{Cr2O7^2-}$). En bas à droite : les réducteurs les plus forts ($\ce{S2O3^2-}$, $\ce{SO2}$, $\ce{I-}$).}
\end{popfigure}

Cette échelle est une véritable \emph{carte des forces} : elle permet de prévoir, sans aucune expérience, le sens d'une réaction d'oxydoréduction. Retenons les positions clés : $\ce{MnO4-}$ est l'oxydant le plus puissant de notre liste, $\ce{S2O3^2-}$ le réducteur le plus puissant.

\courssub{Prévoir une réaction spontanée : la règle du gamma en action}

\begin{exemplebox}[Exemple chiffré : le permanganate oxyde-t-il le fer(II) ?]
Données : $E^\circ(\ce{MnO4-}/\ce{Mn^2+}) = \SI{+1,51}{V}$ et $E^\circ(\ce{Fe^3+}/\ce{Fe^2+}) = \SI{+0,77}{V}$.
\begin{itemize}
\item Le couple du permanganate a le potentiel le plus élevé : $\ce{MnO4-}$ est l'oxydant le plus fort.
\item Le couple du fer a le potentiel le plus faible : $\ce{Fe^2+}$ est le réducteur le plus fort.
\item Conclusion par la règle du gamma : $\ce{MnO4-}$ oxyde \textbf{spontanément} $\ce{Fe^2+}$ selon :
\[ \ce{MnO4- + 5 Fe^2+ + 8 H+ -> Mn^2+ + 5 Fe^3+ + 4 H2O} \]
\end{itemize}
Commentaire : cette réaction est rapide, totale et facile à repérer (disparition de la teinte violette) : c'est précisément pour cela qu'elle sert de \emph{réaction support} au dosage manganimétrique du fer, que nous étudierons dans la suite de la leçon.
\end{exemplebox}

\begin{popfigure}
\begin{center}
\begin{tikzpicture}[scale=1.0]
% Echelle verticale simplifiée : y = 2*E ; Fe : 0.77 -> 1.54 ; Mn : 1.51 -> 3.02
\draw[->, thick, popdark] (0,0) -- (0,6.4) node[above] {$E^\circ$ (V)};
\draw[popline, dashed] (-3.2,1.54) -- (3.2,1.54);
\node[popblue] at (-2.0,1.54) {\small $\ce{Fe^3+}$};
\node[popblue] at (2.0,1.54) {\small $\ce{Fe^2+}$};
\node[popmuted, left] at (-3.3,1.54) {\scriptsize $+0,77$ V};
\draw[popline, dashed] (-3.2,3.02) -- (3.2,3.02);
\node[poppink] at (-2.0,3.02) {\small $\ce{MnO4-}$};
\node[poppink] at (2.0,3.02) {\small $\ce{Mn^2+}$};
\node[popmuted, left] at (-3.3,3.02) {\scriptsize $+1,51$ V};
% Gamma : de MnO4- (haut gauche) vers Fe2+ (bas droite)
\draw[->, very thick, poporange, rounded corners=6pt] (-1.6,3.02) -- (-1.6,4.6) -- (1.6,4.6) -- (1.6,1.9);
\node[poporange, align=center] at (0,4.95) {\small $\ce{MnO4-}$ oxyde $\ce{Fe^2+}$};
\draw[->, thick, poppink] (-2.6,3.02) -- (-2.6,2.2) node[left, midway] {\scriptsize réduction};
\draw[->, thick, popblue] (2.6,1.54) -- (2.6,2.4) node[right, midway] {\scriptsize oxydation};
\node[popdark, align=center] at (0,0.6) {\small Produits : $\ce{Mn^2+}$ et $\ce{Fe^3+}$};
\end{tikzpicture}
\end{center}
\legende{Règle du gamma appliquée aux couples $\ce{MnO4-}/\ce{Mn^2+}$ et $\ce{Fe^3+}/\ce{Fe^2+}$ : la réaction spontanée relie l'oxydant placé le plus haut ($\ce{MnO4-}$) au réducteur placé le plus bas ($\ce{Fe^2+}$), dessinant un $\gamma$.}
\end{popfigure}

\begin{exoresolu}[Le thiosulfate réduit-il le diiode ?]
\textbf{Énoncé.} On donne $E^\circ(\ce{I2}/\ce{I-}) = \SI{+0,54}{V}$ et $E^\circ(\ce{S4O6^2-}/\ce{S2O3^2-}) = \SI{+0,09}{V}$. Une solution brune de diiode est versée dans une solution incolore de thiosulfate de sodium. Prévoir s'il se produit une réaction et, le cas échéant, écrire son équation-bilan.
\tcblower
\textbf{Solution détaillée.}
\begin{enumerate}
\item \textbf{Identification des forces.} Le couple $\ce{I2}/\ce{I-}$ a le potentiel le plus élevé : $\ce{I2}$ est l'oxydant le plus fort. Le couple $\ce{S4O6^2-}/\ce{S2O3^2-}$ a le potentiel le plus faible : $\ce{S2O3^2-}$ est le réducteur le plus fort.
\item \textbf{Règle du gamma.} La réaction spontanée a donc lieu entre $\ce{I2}$ et $\ce{S2O3^2-}$.
\item \textbf{Demi-équations.}
\[ \ce{I2 + 2 e- -> 2 I-} \qquad \ce{2 S2O3^2- -> S4O6^2- + 2 e-} \]
\item \textbf{Équation-bilan} (les 2 électrons échangés se simplifient) :
\[ \ce{I2 + 2 S2O3^2- -> 2 I- + S4O6^2-} \]
\end{enumerate}
Vérification : atomes conservés (2 I, 4 S, 6 O de chaque côté) et charges conservées ($-4$ de chaque côté). Observation : la teinte brune du diiode disparaît. Cette réaction, rapide et totale, est la base du dosage iodométrique que nous détaillerons plus loin.
\end{exoresolu}

\begin{attention}[Deux réflexes à acquérir dès maintenant]
\begin{itemize}
\item Avant de conclure, vérifie toujours \textbf{quelle espèce est présente} dans le milieu : l'échelle des potentiels ne prévoit une réaction qu'entre un oxydant et un réducteur \emph{effectivement mis en contact}.
\item L'équation-bilan ne doit \textbf{jamais faire apparaître d'électrons} : multiplie les demi-équations par les coefficients qui égalisent les électrons échangés avant de les additionner.
\end{itemize}
\end{attention}

\begin{savaistu}[L'éthylotest : la chimie au service de la sécurité routière]
Le tube de l'éthylotest classique contient du dichromate de potassium $\ce{K2Cr2O7}$ en milieu acide, de couleur jaune-orangée. Si l'air expiré contient de l'éthanol $\ce{C2H5OH}$, celui-ci est oxydé en acide éthanoïque tandis que le dichromate est réduit en ions chrome(III) $\ce{Cr^3+}$, verts :
\[ \ce{2 Cr2O7^2- + 3 C2H5OH + 16 H+ -> 4 Cr^3+ + 3 CH3COOH + 11 H2O} \]
Le virage de l'orange au vert trahit donc la présence d'alcool. Aux contrôles de gendarmerie de Yaoundé ou de Douala, c'est la règle du gamma qui monte la garde : $E^\circ(\ce{Cr2O7^2-}/\ce{Cr^3+}) = \SI{+1,33}{V}$, très supérieur au potentiel du couple de l'éthanol, la réaction est spontanée. Aujourd'hui, les éthylotests électroniques utilisent des piles à combustible, mais le principe reste électrochimique.
\end{savaistu}

\begin{aretenir}[L'essentiel de cette section]
\begin{itemize}
\item Un couple Ox/Red échange des électrons selon $\text{Ox} + n\,\ce{e-} \ce{<=>} \text{Red}$.
\item En milieu acide, on équilibre une demi-équation en 4 étapes : atomes autres que O et H, puis O avec $\ce{H2O}$, puis H avec $\ce{H+}$, puis charges avec $\ce{e-}$.
\item Plus $E^\circ$ est élevé, plus l'oxydant du couple est fort.
\item Règle du gamma : l'oxydant du couple de plus haut potentiel réagit spontanément avec le réducteur du couple de plus bas potentiel.
\item Couples clés à connaître : $\ce{MnO4-}/\ce{Mn^2+}$ ($+1,51$ V), $\ce{Cr2O7^2-}/\ce{Cr^3+}$ ($+1,33$ V), $\ce{Fe^3+}/\ce{Fe^2+}$ ($+0,77$ V), $\ce{I2}/\ce{I-}$ ($+0,54$ V), $\ce{S4O6^2-}/\ce{S2O3^2-}$ ($+0,09$ V).
\end{itemize}
\end{aretenir}

Maintenant que la carte des forces est établie et que la règle du gamma n'a plus de secret pour toi, passons à l'étape suivante : construire rigoureusement les équations-bilan d'oxydoréduction avec tous ces couples.

\coursec{Équations-bilan d'oxydoréduction avec les couples au programme}

Dans la section précédente, nous avons classé les couples oxydant/réducteur selon leur potentiel normal d'électrode $E^\circ$ et nous avons appris à prévoir le sens d'une réaction spontanée grâce à la règle du gamma. Mais prévoir qu'une réaction \emph{a lieu} ne suffit pas : pour exploiter cette réaction (doser une solution, interpréter un éthylotest, comprendre la décoloration du permanganate), il faut écrire son \cle{équation-bilan}, c'est-à-dire l'équation chimique complète, équilibrée, sans électrons apparents. C'est l'objet de cette section : tu vas apprendre une méthode infaillible pour combiner deux demi-équations électroniques, puis tu l'appliqueras aux grands couples du programme.

\courssub{Principe : combiner deux demi-équations électroniques}

\begin{definition}[Équation-bilan d'oxydoréduction]
L'\cle{équation-bilan} d'une réaction d'oxydoréduction est l'équation chimique globale obtenue en combinant la demi-équation d'oxydation du réducteur et la demi-équation de réduction de l'oxydant, de telle sorte que \textbf{les électrons n'apparaissent pas} dans le bilan final.
\end{definition}

L'idée physique est simple et belle : les électrons ne se baladent jamais librement en solution aqueuse. Chaque électron cédé par le réducteur est \emph{immédiatement} capté par l'oxydant. Il doit donc y avoir \textbf{autant d'électrons cédés que d'électrons captés}. Si le réducteur cède 1 électron et que l'oxydant en capte 5, il faudra 5 réducteurs pour 1 oxydant : c'est tout le secret des coefficients stœchiométriques.

\begin{methode}[Écrire une équation-bilan d'oxydoréduction]
\begin{enumerate}
\item \textbf{Identifier} l'oxydant qui se réduit et le réducteur qui s'oxyde (à l'aide des potentiels normaux ou des observations expérimentales).
\item \textbf{Écrire} les deux demi-équations électroniques dans le sens où elles se produisent réellement (réduction pour l'oxydant, oxydation pour le réducteur).
\item \textbf{Multiplier} chaque demi-équation par le plus petit coefficient entier tel que le nombre d'électrons cédés égale le nombre d'électrons captés (PPCM des nombres d'électrons).
\item \textbf{Additionner} membre à membre les deux demi-équations : les électrons s'éliminent.
\item \textbf{Simplifier} éventuellement les espèces \ce{H+} ou \ce{H2O} présentes des deux côtés.
\item \textbf{Vérifier} la conservation des éléments et la conservation de la charge électrique totale.
\end{enumerate}
\end{methode}

\begin{attention}[Erreurs fréquentes dans le bilan final]
Deux fautes reviennent sans cesse au Probatoire :
\begin{itemize}
\item laisser apparaître des \textbf{électrons} dans l'équation-bilan finale : c'est rédhibitoire, les électrons doivent avoir disparu ;
\item laisser des ions \ce{H+} ou des molécules \ce{H2O} \textbf{des deux côtés} de l'équation : il faut les simplifier (par exemple $16\,\ce{H+}$ à gauche et $10\,\ce{H+}$ à droite se simplifient en $6\,\ce{H+}$ à gauche).
\end{itemize}
\end{attention}

\courssub{Exemple type 1 : oxydation des ions \ce{Fe^2+} par les ions permanganate \ce{MnO4-}}

C'est \emph{la} réaction reine de cette leçon : elle servira de support au dosage manganimétrique (section 4). Observons d'abord : quand on verse une solution violette de permanganate de potassium dans une solution d'ions fer(II) acidifiée, la solution se décolore. Le violet de \ce{MnO4-} disparaît car il est réduit en \ce{Mn^2+} (incolore), tandis que \ce{Fe^2+} est oxydé en \ce{Fe^3+} (jaune pâle).

\textbf{Étape 1 et 2 : les deux demi-équations.}

Réduction du permanganate (couple \ce{MnO4-}/\ce{Mn^2+}, $E^\circ = \SI{1,51}{V}$) :
\[ \ce{MnO4- + 8H+ + 5e- -> Mn^2+ + 4H2O} \]

Oxydation du fer(II) (couple \ce{Fe^3+}/\ce{Fe^2+}, $E^\circ = \SI{0,77}{V}$) :
\[ \ce{Fe^2+ -> Fe^3+ + e-} \]

\textbf{Étape 3 : égaliser les électrons.} \ce{MnO4-} capte 5 électrons, \ce{Fe^2+} n'en cède qu'un. Il faut donc multiplier la demi-équation du fer par \textbf{5} et celle du permanganate par \textbf{1} :

\begin{popfigure}
\begin{tikzpicture}[
  bloc/.style={draw=popblue, thick, rounded corners=4pt, fill=popblueL, text width=5.6cm, align=center, minimum height=2.6cm},
  bloc2/.style={draw=poporange, thick, rounded corners=4pt, fill=poporangeL, text width=5.6cm, align=center, minimum height=2.6cm},
  res/.style={draw=popgreen, thick, rounded corners=4pt, fill=popgreenL, text width=12.6cm, align=center},
  fleche/.style={-{Stealth[length=3mm]}, very thick, popdark}]
\node[bloc] (ox) at (0,0) {\textbf{Oxydation ($\times 5$)}\\[2pt] $\ce{Fe^2+ -> Fe^3+ + e-}$ \quad $\times 5$\\[3pt] $\ce{5Fe^2+ -> 5Fe^3+ + 5e-}$};
\node[bloc2] (red) at (8.6,0) {\textbf{Réduction ($\times 1$)}\\[2pt] $\ce{MnO4- + 8H+ + 5e- -> Mn^2+ + 4H2O}$ \quad $\times 1$\\[3pt] $\ce{MnO4- + 8H+ + 5e- -> Mn^2+ + 4H2O}$};
\draw[fleche, poppurple] (ox.south) -- ++(0,-0.9) -| node[pos=0.25, below, poppurple]{\textbf{les 5 électrons cédés = les 5 électrons captés : ils s'éliminent}} (red.south);
\node[res] at (4.3,-3.4) {\textbf{Équation-bilan :}\quad $\ce{MnO4- + 5Fe^2+ + 8H+ -> Mn^2+ + 5Fe^3+ + 4H2O}$};
\draw[fleche] (4.3,-1.6) -- (4.3,-2.6);
\end{tikzpicture}
\legende{Combinaison des demi-équations : le facteur 5 égalise les électrons échangés, qui disparaissent du bilan.}
\end{popfigure}

\textbf{Étape 6 : vérification systématique.} Cette étape, trop souvent négligée, est ton filet de sécurité.

\begin{exemplebox}[Vérification de l'équation-bilan]
Vérifions $\ce{MnO4- + 5Fe^2+ + 8H+ -> Mn^2+ + 5Fe^3+ + 4H2O}$.

\emph{Conservation des éléments :}
\begin{center}
\begin{tabularx}{\linewidth}{|l|X|X|X|X|}
\hline
\rowcolor{popblueL} Élément & Mn & O & Fe & H \\
\hline
À gauche & 1 & 4 & 5 & 8 \\
\hline
À droite & 1 & 4 (dans $4\,\ce{H2O}$) & 5 & 8 \\
\hline
\end{tabularx}
\end{center}

\emph{Conservation de la charge :}
\begin{align*}
\text{gauche : } & (-1) + 5\times(+2) + 8\times(+1) = -1 + 10 + 8 = +17\\
\text{droite : } & (+2) + 5\times(+3) + 4\times 0 = +17
\end{align*}
Les éléments et la charge ($+17$ de chaque côté) sont conservés : l'équation est correcte.
\end{exemplebox}

\begin{aretenir}[La réaction manganimétrique à connaître par cœur]
\[ \ce{MnO4- + 5Fe^2+ + 8H+ -> Mn^2+ + 5Fe^3+ + 4H2O} \]
Réaction \textbf{totale}, \textbf{rapide} et \textbf{unique} : elle possède les trois qualités d'une réaction de dosage. Remarque le coefficient \textbf{5} devant \ce{Fe^2+} : c'est lui qui commande tous les calculs du dosage manganimétrique.
\end{aretenir}

\courssub{Exemple type 2 : réaction du diiode avec les ions thiosulfate}

Cette réaction est le cœur de l'\cle{iodométrie}, la technique de dosage du diiode (section 5). Le diiode \ce{I2} (solution brune) est un oxydant modéré ($E^\circ = \SI{0,54}{V}$) ; l'ion thiosulfate \ce{S2O3^2-} est un réducteur qui s'oxyde en ion tétrathionate \ce{S4O6^2-} ($E^\circ = \SI{0,08}{V}$).

\textbf{Demi-équations :}
\[ \ce{I2 + 2e- -> 2I-} \qquad \text{(réduction, } \times 1\text{)} \]
\[ \ce{2S2O3^2- -> S4O6^2- + 2e-} \qquad \text{(oxydation, } \times 1\text{)} \]

Les deux demi-équations échangent déjà le même nombre d'électrons (2) : les coefficients multiplicateurs sont tous égaux à 1. En additionnant :

\[ \ce{I2 + 2S2O3^2- -> 2I- + S4O6^2-} \]

\textbf{Vérification des charges :} gauche : $0 + 2\times(-2) = -4$ ; droite : $2\times(-1) + (-2) = -4$. Parfait.

Expérimentalement, la solution brune de diiode se décolore progressivement : quand tout le diiode a réagi, la solution devient incolore (\ce{I-} et \ce{S4O6^2-} sont incolores). C'est ce changement de teinte qui permettra de repérer l'équivalence.

\begin{savaistu}[Pourquoi le diiode est-il dissous dans l'iodure de potassium ?]
Le diiode \ce{I2} est très peu soluble dans l'eau pure (environ \SI{0,3}{g.L^{-1}} à \SI{25}{\celsius}). Pour préparer des solutions de diiode, on le dissout dans une solution d'\textbf{iodure de potassium} \ce{K+ + I-} : le diiode s'y combine à l'ion iodure pour former l'ion \cle{triiodure} \ce{I3-}, très soluble et de couleur brune :
\[ \ce{I2 + I- <=> I3-} \]
L'ion \ce{I3-} se comporte exactement comme \ce{I2} vis-à-vis des réducteurs : il libère le diiode au fur et à mesure de sa consommation. C'est pourquoi, dans les équations-bilan, on écrit simplement \ce{I2}.
\end{savaistu}

\courssub{Exemple type 3 : action des ions dichromate — l'éthylotest}

Le couple \ce{Cr2O7^2-}/\ce{Cr^3+} ($E^\circ = \SI{1,33}{V}$) est un oxydant puissant en milieu acide. Sa demi-équation de réduction est :
\[ \ce{Cr2O7^2- + 14H+ + 6e- -> 2Cr^3+ + 7H2O} \]

\textbf{Application à l'éthylotest.} L'éthanol \ce{CH3CH2OH} expiré par un conducteur est oxydé en acide éthanoïque \ce{CH3COOH} (couple \ce{CH3COOH}/\ce{CH3CH2OH}) :
\[ \ce{CH3CH2OH + H2O -> CH3COOH + 4H+ + 4e-} \]

Pour combiner : le dichromate capte 6 électrons, l'éthanol en cède 4. Le plus petit commun multiple est 12 : on multiplie la première par \textbf{2} et la seconde par \textbf{3} :
\begin{align*}
\ce{2Cr2O7^2- + 28H+ + 12e- &-> 4Cr^3+ + 14H2O}\\
\ce{3CH3CH2OH + 3H2O &-> 3CH3COOH + 12H+ + 12e-}
\end{align*}
En additionnant puis en simplifiant les \ce{H+} ($28-12=16$ à gauche) et les \ce{H2O} ($14-3=11$ à droite) :
\[ \ce{2Cr2O7^2- + 3CH3CH2OH + 16H+ -> 4Cr^3+ + 3CH3COOH + 11H2O} \]

C'est le virage de l'\textbf{orange} (\ce{Cr2O7^2-}) au \textbf{vert} (\ce{Cr^3+}) qui trahit la présence d'alcool dans l'air expiré. Vérifie la charge : gauche $2\times(-2)+16 = +12$ ; droite $4\times(+3) = +12$. 

\textbf{Variante : le dichromate oxyde aussi les ions \ce{Fe^2+}.} Même méthode : 6 électrons captés par \ce{Cr2O7^2-}, 1 électron cédé par \ce{Fe^2+}, donc on multiplie le fer par 6 :
\[ \ce{Cr2O7^2- + 6Fe^2+ + 14H+ -> 2Cr^3+ + 6Fe^3+ + 7H2O} \]

\courssub{Le rôle du milieu : pourquoi acidifier, et avec quel acide ?}

Tu as remarqué que les demi-équations de \ce{MnO4-} et de \ce{Cr2O7^2-} font intervenir des ions \ce{H+} : ces oxydants \textbf{n'agissent qu'en milieu acide}. En pratique, on ajoute toujours quelques millilitres d'\textbf{acide sulfurique dilué} (\ce{2H+ + SO4^2-}) à la solution avant le dosage. Sans acidification, le permanganate est réduit non pas en \ce{Mn^2+} incolore mais en dioxyde de manganèse \ce{MnO2}, solide brun qui trouble la solution et fausse le repérage de l'équivalence.

\begin{attention}[Jamais d'acide chlorhydrique avec le permanganate !]
Pour acidifier un milieu contenant \ce{MnO4-}, il est \textbf{interdit} d'utiliser l'acide chlorhydrique (\ce{H+ + Cl-}). En effet, $E^\circ(\ce{MnO4-}/\ce{Mn^2+}) = \SI{1,51}{V} > E^\circ(\ce{Cl2}/\ce{Cl-}) = \SI{1,36}{V}$ : le permanganate oxyderait les ions chlorure en dichlore, gaz toxique :
\[ \ce{2MnO4- + 10Cl- + 16H+ -> 2Mn^2+ + 5Cl2 + 8H2O} \]
Le permanganate serait consommé par une réaction parasite et le dosage serait faux. On utilise donc l'\textbf{acide sulfurique dilué}, dont l'anion \ce{SO4^2-} n'est pas oxydé dans ces conditions.
\end{attention}

\courssub{Tableau récapitulatif des couples au programme}

\begin{center}
\begin{tabularx}{\linewidth}{|l|X|l|}
\hline
\rowcolor{popblueL} \textbf{Couple Ox/Red} & \textbf{Demi-équation électronique (réduction)} & \textbf{Couleurs en solution} \\
\hline
\ce{MnO4-}/\ce{Mn^2+} & \ce{MnO4- + 8H+ + 5e- <=> Mn^2+ + 4H2O} & violet / incolore \\
\hline
\ce{Cr2O7^2-}/\ce{Cr^3+} & \ce{Cr2O7^2- + 14H+ + 6e- <=> 2Cr^3+ + 7H2O} & orange / vert \\
\hline
\ce{NO3-}/\ce{NO} & \ce{NO3- + 4H+ + 3e- <=> NO + 2H2O} & incolore / gaz incolore \\
\hline
\ce{SO4^2-}/\ce{SO2} & \ce{SO4^2- + 4H+ + 2e- <=> SO2 + 2H2O} & incolore / gaz \\
\hline
\ce{Fe^3+}/\ce{Fe^2+} & \ce{Fe^3+ + e- <=> Fe^2+} & jaune pâle / verdâtre \\
\hline
\ce{Cl2}/\ce{Cl-} & \ce{Cl2 + 2e- <=> 2Cl-} & jaune verdâtre / incolore \\
\hline
\ce{S4O6^2-}/\ce{S2O3^2-} & \ce{S4O6^2- + 2e- <=> 2S2O3^2-} & incolore / incolore \\
\hline
\ce{I2}/\ce{I-} & \ce{I2 + 2e- <=> 2I-} & brun / incolore \\
\hline
\end{tabularx}
\end{center}

\courssub{Entraînement guidé : deux bilans à construire toi-même}

\begin{exoresolu}[Action de l'acide nitrique sur le cuivre]
Le cuivre métallique (couple \ce{Cu^2+}/\ce{Cu}) est attaqué par l'acide nitrique : les ions nitrate \ce{NO3-} sont réduits en monoxyde d'azote \ce{NO}. Écris l'équation-bilan de la réaction.
\tcblower
\textbf{Demi-équations :}
\begin{align*}
\ce{NO3- + 4H+ + 3e- &-> NO + 2H2O} && (\times 2)\\
\ce{Cu &-> Cu^2+ + 2e-} && (\times 3)
\end{align*}
Le ppcm de 3 et 2 est 6 : on multiplie la réduction par 2 et l'oxydation par 3, puis on additionne :
\[ \ce{2NO3- + 3Cu + 8H+ -> 2NO + 3Cu^2+ + 4H2O} \]
\textbf{Vérification des charges :} gauche : $2\times(-1) + 8 = +6$ ; droite : $3\times(+2) = +6$. Correct. C'est cette réaction qui produit les vapeurs rousses de \ce{NO2} (formé par oxydation de \ce{NO} à l'air) observées quand on attaque un fil de cuivre avec l'acide nitrique.
\end{exoresolu}

\begin{exoresolu}[Réduction des ions sulfate par le cuivre]
L'acide sulfurique concentré et chaud oxyde le cuivre : le couple mis en jeu est \ce{SO4^2-}/\ce{SO2}. Écris l'équation-bilan.
\tcblower
\textbf{Demi-équations :}
\begin{align*}
\ce{SO4^2- + 4H+ + 2e- &-> SO2 + 2H2O} && (\times 1)\\
\ce{Cu &-> Cu^2+ + 2e-} && (\times 1)
\end{align*}
Les deux demi-équations échangent 2 électrons : addition directe :
\[ \ce{SO4^2- + Cu + 4H+ -> SO2 + Cu^2+ + 2H2O} \]
\textbf{Vérification des charges :} gauche : $-2 + 4 = +2$ ; droite : $+2$. Correct. Remarque : contrairement à l'acide sulfurique \emph{dilué} (qui n'attaque pas le cuivre et sert à acidifier les milieux), l'acide sulfurique \emph{concentré et chaud} est un oxydant. Le milieu change tout !
\end{exoresolu}

\begin{aretenir}[L'essentiel de la section]
\begin{itemize}
\item Une équation-bilan s'obtient en multipliant chaque demi-équation par un coefficient tel que les électrons cédés égalent les électrons captés, puis en additionnant : \textbf{aucun électron ne subsiste}.
\item Les trois bilans à maîtriser absolument :
\[ \ce{MnO4- + 5Fe^2+ + 8H+ -> Mn^2+ + 5Fe^3+ + 4H2O} \]
\[ \ce{I2 + 2S2O3^2- -> 2I- + S4O6^2-} \]
\[ \ce{Cr2O7^2- + 6Fe^2+ + 14H+ -> 2Cr^3+ + 6Fe^3+ + 7H2O} \]
\item Toujours vérifier : conservation des éléments \emph{et} conservation de la charge.
\item \ce{MnO4-} et \ce{Cr2O7^2-} n'oxydent qu'en milieu acide ; on acidifie avec l'acide sulfurique dilué, \textbf{jamais} avec l'acide chlorhydrique.
\end{itemize}
\end{aretenir}

Maintenant que tu sais écrire et vérifier n'importe quelle équation-bilan d'oxydoréduction, il ne te reste qu'un pas pour devenir un véritable analyste : utiliser ces réactions pour \emph{mesurer} des concentrations inconnues. C'est précisément l'objet de la section suivante, consacrée au principe général d'un dosage d'oxydoréduction.

\coursec{Principe général d'un dosage d'oxydoréduction}

Dans la section précédente, nous avons appris à écrire et à équilibrer les équations-bilan d'oxydoréduction mettant en jeu les couples au programme : \ce{MnO4-}/\ce{Mn^2+}, \ce{Cr2O7^2-}/\ce{Cr^3+}, \ce{Fe^3+}/\ce{Fe^2+}, \ce{I2}/\ce{I-}, \ce{S4O6^2-}/\ce{S2O3^2-}... À quoi cela sert-il concrètement ? Imagine un contrôle de qualité à la SONARA ou dans une brasserie de Douala : comment vérifier la concentration en fer(II) d'un médicament contre l'anémie, ou la teneur en chlore actif d'une eau de Javel commerciale ? La réponse est une technique expérimentale fondamentale : le \cle{dosage} (ou \cle{titrage}) d'oxydoréduction. C'est l'une des épreuves classiques du Probatoire, alors maîtrisons-la pas à pas.

\courssub{Qu'est-ce qu'un dosage ?}

\textbf{Intuition.} Tu veux connaître la quantité de sucre contenue dans un verre de jus de bissap, mais tu ne peux pas « peser » le sucre dissous. En revanche, si tu disposes d'un réactif qui réagit \emph{uniquement} avec le sucre, goutte après goutte, et qui te signale (par un changement de couleur, par exemple) l'instant précis où tout le sucre a réagi, alors le volume de réactif versé te renseigne sur la quantité de sucre initiale. C'est exactement l'idée d'un dosage.

\begin{definition}[Dosage ou titrage]
Un \cle{dosage} (ou \cle{titrage}) est une méthode d'analyse quantitative qui consiste à déterminer la \cle{concentration inconnue} d'une espèce chimique en solution, en la faisant réagir avec une solution de concentration \emph{exactement connue}, appelée \cle{solution titrante}. La réaction mise en jeu est appelée \cle{réaction de dosage}.
\end{definition}

Le vocabulaire est précis et il faut l'employer correctement :

\begin{itemize}
\item La \cle{solution titrée} est la solution dont on cherche la concentration. On en prélève un volume $V$ \emph{précisément mesuré} à la \textbf{pipette jaugée}, que l'on place dans l'\textbf{erlenmeyer} (préféré au bécher pour éviter les projections lors de l'agitation).
\item La \cle{solution titrante} est la solution de concentration connue $C_{\text{titrante}}$. Elle est placée dans la \textbf{burette graduée}, qui permet de mesurer le volume versé.
\item Le \cle{volume équivalent} $V_E$ est le volume de solution titrante versé lorsque l'équivalence est atteinte.
\end{itemize}

\courssub{Conditions d'une réaction de dosage}

Toutes les réactions d'oxydoréduction ne conviennent pas pour un dosage. Pour être utilisable, la réaction de dosage doit satisfaire \textbf{trois conditions} impératives :

\begin{enumerate}
\item Elle doit être \cle{totale} : les réactifs doivent être entièrement consommés (constante d'équilibre très grande, $\Delta E^\circ$ élevé). Sinon, une partie de l'espèce à doser resterait en solution et le résultat serait faussé.
\item Elle doit être \cle{rapide} : le changement doit se produire dès l'ajout du réactif, sans temps d'attente notable. Une réaction lente rend le repérage de l'équivalence impossible en temps réel.
\item Elle doit être \cle{unique} : le réactif titrant ne doit réagir qu'avec l'espèce dosée, sans \textbf{réaction parasite} (avec le solvant, l'air, une autre espèce présente, ou le récipient).
\end{enumerate}

\begin{attention}[Une réaction de dosage doit être totale, rapide et unique]
Une erreur classique consiste à croire que n'importe quelle réaction d'oxydoréduction peut servir à doser. C'est faux ! Une réaction \textbf{lente} (comme certaines oxydations par \ce{NO3-} à froid) ou \textbf{partielle} (équilibre non déplacé) ne permet pas de repérer nettement l'équivalence. Au Probatoire, si l'on te demande de justifier le choix d'une réaction de dosage, cite toujours ces trois adjectifs : \emph{totale, rapide et unique}.
\end{attention}

\begin{exemplebox}[La réaction de dosage de référence : dosage manganimétrique]
L'oxydation des ions fer(II) par les ions permanganate en milieu acide sulfurique :
\[ \ce{MnO4- + 5Fe^2+ + 8H+ -> Mn^2+ + 5Fe^3+ + 4H2O} \]
Cette réaction est \textbf{totale} (les potentiels normaux $E^\circ(\ce{MnO4-}/\ce{Mn^2+}) = \SI{1,51}{V}$ et $E^\circ(\ce{Fe^3+}/\ce{Fe^2+}) = \SI{0,77}{V}$ sont très éloignés, $\Delta E^\circ = \SI{0,74}{V}$), \textbf{rapide} à température ambiante et \textbf{unique} si le milieu est correctement acidifié à l'acide sulfurique \ce{H2SO4} (l'acide chlorhydrique serait oxydé en \ce{Cl2} par \ce{MnO4-}, l'acide nitrique est lui-même oxydant). Elle remplit donc parfaitement les trois conditions d'un dosage.
\end{exemplebox}

\courssub{Le point d'équivalence}

\textbf{Intuition.} Versons progressivement la solution titrante dans la solution titrée. Au début, chaque goutte de titrant est immédiatement consommée : il reste encore de l'espèce à doser. Puis vient un instant magique où la quantité d'espèce à doser est \emph{exactement} épuisée : c'est l'\cle{équivalence}. La goutte suivante met le titrant en excès, et c'est cet excès qui va se trahir par une couleur.

\begin{definition}[Équivalence d'un dosage]
L'\cle{équivalence} d'un dosage est atteinte lorsque les réactifs ont été introduits dans les \cle{proportions stœchiométriques} de l'équation de dosage. Les deux réactifs sont alors \textbf{entièrement consommés} : ni l'un ni l'autre n'est en excès. Le volume de solution titrante versé à cet instant est le \cle{volume équivalent}, noté $V_E$.
\end{definition}

\begin{attention}[Ne pas confondre $V$ et $V_E$]
Erreur extrêmement fréquente au Probatoire : confondre le \textbf{volume $V$ de la prise d'essai} (volume de solution titrée prélevé à la pipette jaugée et placé dans l'erlenmeyer) avec le \textbf{volume équivalent $V_E$} (volume de solution titrante lu sur la burette à l'équivalence). Ce sont deux volumes différents, mesurés avec deux verreries différentes, qui apparaissent à deux endroits différents de la relation de dosage. Souligne-les de couleurs différentes dans tes calculs !
\end{attention}

\courssub{Repérage de l'équivalence : le dosage colorimétrique}

Comment savoir expérimentalement que l'équivalence est atteinte ? En oxydoréduction, la méthode la plus simple est le \cle{dosage colorimétrique} : on observe un \textbf{changement de couleur} de la solution dans l'erlenmeyer.

\textbf{Premier cas : l'auto-indicateur.} Certains réactifs titrants sont eux-mêmes colorés, alors que leur forme réduite est incolore. C'est le cas de l'ion permanganate \ce{MnO4-}, violet intense, qui est réduit en ion \ce{Mn^2+} quasi incolore (rose très pâle, masqué par le milieu). Tant qu'il reste des ions \ce{Fe^2+} dans l'erlenmeyer, chaque goutte de \ce{MnO4-} versée est décolorée. À l'équivalence, la première goutte en excès n'est plus consommée : la solution vire au \textbf{rose persistant}. Le permanganate joue donc à la fois le rôle de réactif titrant et d'indicateur de fin de réaction : on dit qu'il est \cle{auto-indicateur}.

\textbf{Deuxième cas : l'indicateur coloré ajouté.} Quand le réactif n'est pas assez coloré, on ajoute un indicateur qui signale sa présence. L'exemple au programme est l'\cle{empois d'amidon}, qui forme avec le diiode \ce{I2} un complexe \textbf{bleu nuit} très intense. Lors du dosage de \ce{I2} par les ions thiosulfate \ce{S2O3^2-} (incolores), la solution est bleue tant qu'il reste du diiode ; à l'équivalence, tout le \ce{I2} a été consommé et la solution se \textbf{décolore} brutalement.

\begin{methode}[Repérer l'équivalence par changement de teinte]
\begin{enumerate}
\item Verser la solution titrante \textbf{rapidement} au début, tant que la décoloration (ou l'absence de coloration) est immédiate.
\item Ralentir dès que la décoloration devient plus lente : on approche de l'équivalence.
\item Verser alors \textbf{goutte à goutte}, en agitant \emph{en permanence} (agitateur magnétique ou agitation manuelle vigoureuse) pour homogénéiser le mélange.
\item S'arrêter dès que la teinte (rose pour \ce{MnO4-}, bleue pour l'empois d'amidon) devient \textbf{persistante} (ou disparaît durablement, selon le cas) après agitation : c'est l'équivalence.
\item Lire le volume équivalent $V_E$ sur la burette, au bas du ménisque, l'œil à hauteur du ménisque pour éviter l'erreur de parallaxe.
\end{enumerate}
\end{methode}

\courssub{Le dispositif expérimental}

\begin{popfigure}
\begin{tikzpicture}[scale=1.0, every node/.style={font=\small}]
% Support
\draw[popdark, line width=1.2pt] (-0.2,0) -- (3.2,0);
\draw[popdark, line width=1.2pt] (0.4,0) -- (0.4,7.2);
% Pince
\draw[popdark, line width=1pt] (0.4,5.6) -- (1.35,5.6);
\draw[popdark, line width=1pt] (1.35,5.45) -- (1.35,5.75);
% Burette
\draw[popblue, line width=1pt] (1.55,7.0) -- (1.55,2.6) -- (1.65,2.3) -- (1.65,1.9);
\draw[popblue, line width=1pt] (2.15,7.0) -- (2.15,2.6) -- (2.05,2.3) -- (2.05,1.9);
% Robinet
\draw[popdark, line width=1pt] (1.55,2.35) -- (2.15,2.35);
\fill[popdark] (1.85,2.35) circle (0.05);
% Graduations
\foreach \y in {3.0,3.6,4.2,4.8,5.4,6.0,6.6} \draw[popblue] (2.15,\y) -- (2.0,\y);
% Liquide titrant
\fill[poppinkL] (1.57,4.6) rectangle (2.13,6.95);
\draw[popink] (1.57,4.6) -- (2.13,4.6);
% Goutte
\fill[popink] (1.85,1.55) circle (0.045);
% Erlenmeyer (conique)
\draw[popdark, line width=1pt] (1.3,0.15) -- (1.0,1.5) -- (2.7,1.5) -- (2.4,0.15);
\fill[popgoldL] (1.05,0.17) -- (1.3,0.95) -- (2.45,0.95) -- (2.2,0.17) -- cycle;
% Barreau magnétique
\fill[popdark, rounded corners=1pt] (1.6,0.2) rectangle (2.1,0.32);
% Agitateur magnétique
\draw[popmuted, line width=1.2pt, rounded corners=2pt] (0.6,-0.05) rectangle (3.1,-0.55);
\node[popmuted] at (1.85,-0.3) {\footnotesize agitateur magnétique};
% Légendes
\node[popink, anchor=west] at (3.4,5.8) {\footnotesize burette graduée :};
\node[popink, anchor=west] at (3.4,5.45) {\footnotesize solution titrante};
\node[popink, anchor=west] at (3.4,5.1) {\footnotesize (concentration $C_{\text{titrante}}$ connue)};
\node[popink, anchor=west] at (3.4,1.0) {\footnotesize erlenmeyer : solution titrée,};
\node[popink, anchor=west] at (3.4,0.65) {\footnotesize volume $V$ prélevé à la pipette jaugée};
\node[popink, anchor=west] at (3.4,2.1) {\footnotesize volume versé à l'équivalence : $V_E$};
\draw[popink, ->] (3.35,2.0) -- (2.2,1.75);
\draw[popink, ->] (3.35,0.9) -- (2.5,0.7);
\draw[popink, ->] (3.35,5.6) -- (2.2,5.8);
\end{tikzpicture}
\legende{Dispositif d'un dosage d'oxydoréduction : la solution titrante est dans la burette graduée ; la solution titrée (volume $V$ exact) est dans l'erlenmeyer, maintenu en agitation permanente.}
\end{popfigure}

\begin{attention}[Précision du dosage]
La qualité d'un dosage dépend de la verrerie et des gestes : prélèvement du volume $V$ à la \textbf{pipette jaugée} (jamais à l'éprouvette ni au bécher !), lecture de $V_E$ sur la \textbf{burette graduée} au dixième de millilitre (\SI{0,1}{mL}), versement \textbf{goutte à goutte} près de l'équivalence et \textbf{agitation permanente}. Un dosage réussi se répète : on effectue au moins deux essais concordants (écart inférieur à \SI{0,2}{mL}) et l'on retient la moyenne des $V_E$.
\end{attention}

\courssub{Relation à l'équivalence et exploitation}

C'est ici que tout se joue au Probatoire. Considérons l'équation générale d'un dosage d'oxydoréduction :
\[ \nu_{\text{titrant}}\,\ce{Esp_{\text{titrant}}} + \nu_{\text{titré}}\,\ce{Esp_{\text{titré}}} \ce{->} \text{produits} \]
où la solution titrante contient l'espèce \ce{Esp_{\text{titrant}}} (concentration $C_{\text{titrante}}$) et la prise d'essai contient l'espèce \ce{Esp_{\text{titré}}} (concentration inconnue $C$, volume $V$). Les coefficients stœchiométriques sont $\nu_{\text{titrant}}$ et $\nu_{\text{titré}}$.

À l'équivalence, les réactifs ont été mélangés dans les proportions stœchiométriques, c'est-à-dire :
\begin{aretenir}[Relation fondamentale à l'équivalence]
\[ \frac{n(\ce{Esp_{\text{titrant}}}\ \text{versé})}{\nu_{\text{titrant}}} = \frac{n(\ce{Esp_{\text{titré}}}\ \text{présent})}{\nu_{\text{titré}}} \]
soit, en introduisant concentrations et volumes :
\[ \frac{C_{\text{titrante}} \times V_E}{\nu_{\text{titrant}}} = \frac{C \times V}{\nu_{\text{titré}}} \quad\Longrightarrow\quad C = \frac{\nu_{\text{titré}}}{\nu_{\text{titrant}}}\times \frac{C_{\text{titrante}} \times V_E}{V} \]
Les volumes $V_E$ et $V$ doivent être exprimés dans la \textbf{même unité} (peu importe laquelle, car ils se simplifient dans le rapport).
\end{aretenir}

\begin{popfigure}
\begin{tikzpicture}
\begin{axis}[
width=0.85\linewidth, height=6cm,
xlabel={Volume de titrant versé $V$ (mL)},
ylabel={Quantité de matière (mmol)},
xmin=0, xmax=20, ymin=0, ymax=3.2,
legend style={at={(0.97,0.97)}, anchor=north east, font=\footnotesize},
grid=both, grid style={popline!40},
]
\addplot[popcurve, popink, domain=0:12.5, samples=60] {2.5 - 0.2*x};
\addplot[popcurve, popink, domain=12.5:20, samples=20] {0};
\addlegendentry{$n(\ce{Esp_{\text{titré}}})$ restant}
\addplot[popcurve, poporange, domain=0:12.5, samples=60] {0.16*x};
\addplot[popcurve, poporange, domain=12.5:20, samples=20] {2};
\addlegendentry{$n$(produit) formé}
\addplot[popcurve, popgreen, domain=12.5:20, samples=20] {0.2*(x-12.5)};
\addlegendentry{$n(\ce{Esp_{\text{titrant}}})$ en excès}
\draw[dashed, draw=poppurple, thick] (axis cs:12.5,0) -- (axis cs:12.5,3.2);
\node[color=poppurple, anchor=south, font=\footnotesize] at (axis cs:12.5,3.05) {équivalence $V_E$};
\fill[fill=poppurple] (axis cs:12.5,0) circle (2pt);
\end{axis}
\end{tikzpicture}
\legende{Évolution qualitative des quantités de matière en fonction du volume versé : l'espèce titrée s'épuise, le produit s'accumule ; à l'équivalence $V_E$, l'espèce titrée est entièrement consommée et le titrant commence à s'accumuler en excès.}
\end{popfigure}

\begin{exoresolu}[Dosage des ions fer(II) par le permanganate]
On dose $V = \SI{20,0}{mL}$ d'une solution d'ions \ce{Fe^2+} de concentration inconnue $C$ par une solution de permanganate de potassium de concentration $C_{\text{titrante}} = \SI{0,020}{mol.L^{-1}}$, en milieu acide sulfurique. Le virage au rose persistant est obtenu pour $V_E = \SI{12,5}{mL}$. Calculer $C$.
\tcblower
\textbf{Étape 1 : équation de dosage.}
\[ \ce{MnO4- + 5Fe^2+ + 8H+ -> Mn^2+ + 5Fe^3+ + 4H2O} \]
Les coefficients stœchiométriques sont $\nu_{\text{titrant}} = 1$ pour \ce{MnO4-} et $\nu_{\text{titré}} = 5$ pour \ce{Fe^2+}.

\textbf{Étape 2 : relation à l'équivalence.}
\[ \frac{n(\ce{MnO4-})}{1} = \frac{n(\ce{Fe^2+})}{5} \quad\Longrightarrow\quad n(\ce{Fe^2+}) = 5\,n(\ce{MnO4-}) \]

\textbf{Étape 3 : application numérique.}
\begin{align*}
n(\ce{MnO4-}) &= C_{\text{titrante}} \times V_E = \SI{0,020}{mol.L^{-1}} \times \SI{12,5e-3}{L} = \SI{2,5e-4}{mol}\\
n(\ce{Fe^2+}) &= 5 \times \SI{2,5e-4}{mol} = \SI{1,25e-3}{mol}\\
C &= \frac{n(\ce{Fe^2+})}{V} = \frac{\SI{1,25e-3}{mol}}{\SI{20,0e-3}{L}} = \SI{0,0625}{mol.L^{-1}}
\end{align*}
Vérification par la formule directe : $C = \dfrac{5 \times \num{0,020} \times 12{,}5}{20{,}0} = \SI{0,0625}{mol.L^{-1}}$. La solution titrée a une concentration de \SI{6,25e-2}{mol.L^{-1}} en ions fer(II).
\end{exoresolu}

\begin{attention}[Le piège des coefficients stœchiométriques]
L'oubli du facteur $\nu_{\text{titré}}/\nu_{\text{titrant}}$ (ici $5/1$) est l'erreur la plus fréquente : beaucoup d'élèves écrivent $n(\ce{Fe^2+}) = n(\ce{MnO4-})$, ce qui est faux car l'équation n'est pas « 1 pour 1 ». Réflexe de survie : \textbf{toujours écrire l'équation-bilan équilibrée avant tout calcul}, puis encadrer les coefficients $\nu_{\text{titrant}}$ et $\nu_{\text{titré}}$.
\end{attention}

\begin{savaistu}[L'éthylotest : un dosage d'oxydoréduction dans la vie quotidienne]
Lors d'un contrôle routier à Yaoundé ou à Douala, l'alcootest jetable (tube) contient du dichromate de potassium \ce{K2Cr2O7} (orange) acidifié, déposé sur un support solide (silice). Si l'air expiré contient de l'éthanol \ce{C2H5OH}, celui-ci est oxydé en acide éthanoïque tandis que les ions dichromate sont réduits en ions \ce{Cr^3+} (verts) :
\[ \ce{2Cr2O7^2- + 3C2H5OH + 16H+ -> 4Cr^3+ + 3CH3COOH + 11H2O} \]
La longueur de la zone verdie est proportionnelle à la quantité d'alcool : c'est bien un dosage d'oxydoréduction, fondé sur les mêmes principes que ceux étudiés ici (réaction totale, rapide, unique, repérage colorimétrique) !
\end{savaistu}

\begin{aretenir}[L'essentiel du principe d'un dosage]
\begin{itemize}
\item Un dosage détermine une concentration inconnue grâce à une réaction \textbf{totale, rapide et unique} avec une solution titrante de concentration connue.
\item Solution \textbf{titrée} : volume $V$ prélevé à la pipette jaugée, dans l'erlenmeyer. Solution \textbf{titrante} : dans la burette graduée.
\item À l'\textbf{équivalence}, les réactifs sont dans les proportions stœchiométriques : $\dfrac{n_{\text{titrant}}}{\nu_{\text{titrant}}} = \dfrac{n_{\text{titré}}}{\nu_{\text{titré}}}$.
\item Repérage colorimétrique : rose persistant avec \ce{MnO4-} (auto-indicateur), décoloration du bleu avec l'empois d'amidon pour \ce{I2}.
\item Exploitation : $C = \dfrac{\nu_{\text{titré}}}{\nu_{\text{titrant}}}\times\dfrac{C_{\text{titrante}}\,V_E}{V}$, sans jamais confondre $V$ et $V_E$.
\end{itemize}
\end{aretenir}

Maintenant que le principe général est solidement acquis, passons à la pratique avec les deux dosages incontournables du programme : d'abord le dosage manganimétrique des ions \ce{Fe^2+}, puis l'iodométrie.

\coursec{Dosage des ions Fe2+ par les ions MnO4- (dosage manganimétrique)}

Dans la section précédente, nous avons posé le principe général d'un dosage d'oxydoréduction : faire réagir, goutte à goutte, un réactif titrant de concentration connue sur l'espèce à doser, jusqu'à l'\cle{équivalence}, moment où les deux réactifs ont été mélangés dans les proportions stœchiométriques de l'équation de dosage. Restait une question pratique : comment \emph{voir} l'équivalence sans instrument coûteux ? Le dosage manganimétrique apporte une réponse d'une élégance rare : le réactif titrant, l'ion permanganate \ce{MnO4-}, est si coloré qu'il signale lui-même la fin du dosage. Étudions ce dosage emblématique, grand classique du Probatoire.

\courssub{Contexte : pourquoi doser les ions fer(II) ?}

Le fer est un oligo-élément indispensable : il entre dans la composition de l'hémoglobine qui transporte le dioxygène dans le sang. Une carence en fer provoque l'\cle{anémie}, maladie fréquente que l'on combat notamment par des médicaments à base de \cle{sulfate de fer(II)} \ce{FeSO4}. Le sulfate de fer(II) est aussi utilisé comme amendement dans certains engrais et pour traiter la chlorose (jaunissement) des plantes.

Pour vérifier la teneur en fer d'un comprimé ou d'une solution commerciale, le chimiste doit doser les ions \ce{Fe^2+}. Problème : les ions \ce{Fe^2+} sont quasiment incolores en solution diluée (très légère teinte verdâtre), impossible donc de les « compter » à l'œil. L'astuce consiste à les faire réagir avec un oxydant puissant et intensément coloré : l'ion \cle{permanganate} \ce{MnO4-}, d'un violet profond, apporté par une solution de permanganate de potassium \ce{KMnO4}.

\courssub{Les couples mis en jeu et la spontanéité de la réaction}

\textbf{Les deux couples.} Les demi-équations électroniques, en milieu acide, s'écrivent :

\begin{align*}
\ce{MnO4- + 8 H+ + 5 e- &<=> Mn^2+ + 4 H2O} && E^\circ = \SI{1,51}{V} \\
\ce{Fe^3+ + e- &<=> Fe^2+} && E^\circ = \SI{0,77}{V}
\end{align*}

\textbf{Prévision par la règle du gamma.} Plaçons les deux couples sur l'axe des potentiels : le couple \ce{MnO4-}/\ce{Mn^2+} ($E^\circ = \SI{1,51}{V}$) est \emph{au-dessus} du couple \ce{Fe^3+}/\ce{Fe^2+} ($E^\circ = \SI{0,77}{V}$). La règle du gamma prédit une réaction naturelle et quasi totale entre l'oxydant le plus fort, \ce{MnO4-}, et le réducteur le plus fort, \ce{Fe^2+} :

\begin{popfigure}
\begin{center}
\begin{tikzpicture}[scale=1]
\draw[-{Stealth}, thick, popdark] (0,0) -- (0,4.6) node[above, font=\small] {$E^\circ$ (V)};
\draw[thick, popdark] (-0.12,3.8) -- (0.12,3.8) node[right=4pt, font=\small] {\SI{1,51}{V}};
\draw[thick, popdark] (-0.12,1.6) -- (0.12,1.6) node[right=4pt, font=\small] {\SI{0,77}{V}};
\node[font=\small\bfseries, poppurple, anchor=west] at (1.6,3.8) {\ce{MnO4-}};
\node[font=\small, popmuted, anchor=west] at (5.2,3.8) {\ce{Mn^2+}};
\node[font=\small, popmuted, anchor=west] at (1.6,1.6) {\ce{Fe^3+}};
\node[font=\small\bfseries, popgreen, anchor=west] at (5.2,1.6) {\ce{Fe^2+}};
\draw[-{Stealth}, very thick, poppurple] (2.9,3.8) -- (4.9,1.75);
\draw[-{Stealth}, very thick, popgreen] (6.6,1.6) -- (4.4,3.65);
\node[font=\footnotesize, popdark, align=center] at (4.6,0.6) {r\'eaction naturelle : $\gamma$};
\end{tikzpicture}
\end{center}
\legende{Règle du gamma : \ce{MnO4-} (oxydant placé en haut à gauche) réagit spontanément avec \ce{Fe^2+} (réducteur placé en bas à droite).}
\end{popfigure}

L'écart de potentiel $\Delta E^\circ = 1{,}51 - 0{,}77 = \SI{0,74}{V}$ est grand : la réaction est \cle{totale} et \cle{rapide}, deux qualités indispensables pour une réaction de dosage.

\textbf{Équation-bilan du dosage.} L'ion \ce{MnO4-} capte 5 électrons ; chaque ion \ce{Fe^2+} n'en cède qu'un seul. Il faut donc multiplier la demi-équation du fer par 5 pour que les électrons s'éliminent :

\[ \ce{MnO4- + 5 Fe^2+ + 8 H+ -> Mn^2+ + 5 Fe^3+ + 4 H2O} \]

\begin{aretenir}[L'équation de dosage à connaître par cœur]
\[ \ce{MnO4- + 5 Fe^2+ + 8 H+ -> Mn^2+ + 5 Fe^3+ + 4 H2O} \]
Réaction \textbf{totale}, \textbf{rapide} et \textbf{unique} en milieu acide : c'est une excellente réaction de dosage. Retiens surtout le coefficient \cle{5} devant \ce{Fe^2+} : un ion permanganate « consomme » cinq ions fer(II).
\end{aretenir}

\courssub{Pourquoi acidifier, et pourquoi avec l'acide sulfurique ?}

L'équation-bilan fait apparaître 8 ions \ce{H+} parmi les réactifs : la réaction \emph{consomme} de l'acidité. Avant de commencer le dosage, on ajoute donc à la prise d'essai quelques millilitres d'\cle{acide sulfurique dilué} \ce{H2SO4}. Cette acidification a deux rôles :

\begin{itemize}
\item elle fournit les ions \ce{H+} nécessaires à la réaction et garantit que le permanganate est réduit jusqu'à l'ion \ce{Mn^2+} (incolore) ;
\item en milieu insuffisamment acide, \ce{MnO4-} ne serait réduit qu'en dioxyde de manganèse \ce{MnO2}, solide \textbf{brun} qui trouble la solution et fausse le repérage de l'équivalence.
\end{itemize}

\begin{attention}[Jamais d'acide chlorhydrique !]
On acidifie avec \ce{H2SO4} dilué, \textbf{jamais} avec \ce{HCl}. En effet, l'ion permanganate est un oxydant assez puissant pour oxyder les ions chlorure \ce{Cl-} en dichlore \ce{Cl2} (couple \ce{Cl2}/\ce{Cl-}, $E^\circ = \SI{1,36}{V} < \SI{1,51}{V}$) :
\[ \ce{2 MnO4- + 10 Cl- + 16 H+ -> 2 Mn^2+ + 5 Cl2 + 8 H2O} \]
Une partie du permanganate versé serait alors consommée par les ions chlorure au lieu de réagir avec \ce{Fe^2+} : le volume équivalent serait trop grand et le dosage \textbf{faux}. On évite aussi l'acide nitrique \ce{HNO3}, qui est lui-même un oxydant (couple \ce{NO3-}/\ce{NO}) capable d'oxyder \ce{Fe^2+}.
\end{attention}

\courssub{Dispositif expérimental et déroulement du dosage}

Le montage est celui d'un dosage colorimétrique classique : la solution titrante (permanganate de potassium de concentration connue) est placée dans la burette ; la solution à doser (ions \ce{Fe^2+} acidifiés) est placée dans le bécher, agité en permanence.

\begin{popfigure}
\begin{center}
\begin{tikzpicture}[scale=0.95, every node/.style={font=\small}]
% Support
\draw[very thick, popdark] (-2.6,-0.2) -- (-2.6,7.2);
\draw[very thick, popdark] (-3.4,-0.2) -- (-1.8,-0.2);
\draw[thick, popdark] (-2.6,5.6) -- (-1.1,5.6);
\draw[thick, popdark] (-1.1,5.45) rectangle (-0.75,5.75);
% Burette
\draw[thick, popdark] (-1.15,2.6) -- (-1.15,6.9);
\draw[thick, popdark] (-0.65,2.6) -- (-0.65,6.9);
\draw[thick, popdark] (-1.15,2.6) -- (-1.0,2.1) -- (-1.0,1.75);
\draw[thick, popdark] (-0.65,2.6) -- (-0.8,2.1) -- (-0.8,1.75);
% robinet
\draw[thick, popdark, fill=popcream] (-1.25,2.25) rectangle (-0.55,2.05);
% liquide violet dans burette
\fill[poppurple!55] (-1.13,3.4) rectangle (-0.67,6.85);
\draw[poppurple, thick] (-1.13,3.4) -- (-0.67,3.4);
% graduations
\foreach \y in {3.8,4.3,4.8,5.3,5.8,6.3} \draw[popdark] (-1.15,\y) -- (-1.0,\y);
% goutte
\fill[poppurple!70] (-0.9,1.45) circle (0.07);
% Bécher
\draw[thick, popdark] (-2.1,-0.05) -- (-2.1,1.5) -- (0.3,1.5) -- (0.3,-0.05);
\draw[thick, popdark] (-2.1,-0.05) -- (0.3,-0.05);
% liquide bécher
\fill[popgreen!18] (-2.05,0.0) rectangle (0.25,1.0);
\draw[popgreen!60!black, thick] (-2.05,1.0) -- (0.25,1.0);
% agitateur
\draw[thick, popmuted, fill=popmuted!40] (-0.9,0.18) ellipse (0.35 and 0.1);
% Légendes
\node[align=left, anchor=west, text=poppurple!70!black] at (1.2,5.4) {Burette : solution de \ce{KMnO4}\\ de concentration $C_0$ connue\\ (violette)};
\draw[-{Stealth}, poppurple!70!black] (1.2,5.1) -- (-0.55,4.6);
\node[align=left, anchor=west, text=popgreen!50!black] at (1.2,0.7) {Bécher : volume $V$ de solution\\ de \ce{Fe^2+} acidifiée (\ce{H2SO4}),\\ concentration $C$ inconnue};
\draw[-{Stealth}, popgreen!50!black] (1.2,0.55) -- (0.35,0.6);
\node[align=left, anchor=west, text=popmuted] at (-4.9,0.2) {agitateur\\ magnétique};
\draw[-{Stealth}, popmuted] (-3.0,0.25) -- (-1.35,0.22);
\end{tikzpicture}
\end{center}
\legende{Montage du dosage manganimétrique : le permanganate de potassium (solution titrante, violette) est versé goutte à goutte dans la solution de \ce{Fe^2+} acidifiée.}
\end{popfigure}

\begin{experience}[Protocole du dosage manganimétrique]
\begin{enumerate}
\item Rincer puis remplir la burette avec la solution de permanganate de potassium de concentration $C_0$ connue ; ajuster le zéro. (La solution étant très colorée, on lit le \emph{haut} du ménisque.)
\item Prélever à la pipette jaugée un volume $V = \SI{10,0}{mL}$ de la solution de \ce{Fe^2+} ; l'introduire dans un bécher.
\item Ajouter environ \SI{5}{mL} d'acide sulfurique dilué.
\item Verser \ce{KMnO4} goutte à goutte en agitant : la teinte violette disparaît immédiatement au début.
\item Ralentir l'addition près de l'équivalence ; stopper dès que la teinte rose-violette \textbf{persiste au moins 30 secondes}. Noter le volume équivalent $V_E$.
\end{enumerate}
\end{experience}

\courssub{Repérage de l'équivalence : le permanganate, auto-indicateur}

Voici le point le plus joli de ce dosage. Tant qu'il reste des ions \ce{Fe^2+} dans le bécher, chaque goutte violette de \ce{MnO4-} est immédiatement \cle{décolorée} : le permanganate est consommé et transformé en ions \ce{Mn^2+} incolores. La solution reste donc incolore (ou très légèrement jaunâtre à cause des ions \ce{Fe^3+} formés).

À l'équivalence, il n'y a \textbf{plus aucun ion \ce{Fe^2+}} pour réduire le permanganate : la première goutte de \ce{MnO4-} versée en excès n'est plus détruite et communique à toute la solution une teinte \cle{rose-violette persistante}. On dit que l'ion permanganate joue le rôle d'\cle{auto-indicateur} (ou indicateur coloré interne) : nul besoin d'ajouter un indicateur coloré supplémentaire.

\begin{popfigure}
\begin{center}
\begin{tikzpicture}[scale=1]
% Frise des volumes
\draw[-{Stealth}, thick, popdark] (0,0) -- (12.4,0) node[below left=2pt and -40pt, font=\small] {};
\foreach \x/\lab in {0.2/{$V=0$}, 3.2/{$V<V_E$}, 6.2/{$V \approx V_E$}, 9.2/{$V = V_E$}, 11.6/{$V > V_E$}}
  \draw[thick, popdark] (\x,0.1) -- (\x,-0.1) node[below, font=\footnotesize] {\lab};
% Bécher 1 : incolore
\draw[thick, popdark] (0.6,0.5) -- (0.6,1.7) -- (2.0,1.7) -- (2.0,0.5);
\fill[popgreen!12] (0.64,0.54) rectangle (1.96,1.2);
\node[font=\footnotesize, align=center] at (1.3,2.15) {incolore /\\ verdâtre};
% Bécher 2 : incolore (décoloration)
\draw[thick, popdark] (3.6,0.5) -- (3.6,1.7) -- (5.0,1.7) -- (5.0,0.5);
\fill[popgreen!12] (3.64,0.54) rectangle (4.96,1.2);
\node[font=\footnotesize, align=center] at (4.3,2.15) {chaque goutte\\ est décolorée};
% Bécher 3 : équivalence, rose persistant
\draw[thick, popdark] (8.0,0.5) -- (8.0,1.7) -- (9.4,1.7) -- (9.4,0.5);
\fill[poppink!45] (8.04,0.54) rectangle (9.36,1.2);
\node[font=\footnotesize, align=center, text=poppink!60!black] at (8.7,2.15) {teinte rose\\ persistante};
% Bécher 4 : excès violet
\draw[thick, popdark] (10.6,0.5) -- (10.6,1.7) -- (12.0,1.7) -- (12.0,0.5);
\fill[poppurple!45] (10.64,0.54) rectangle (11.96,1.2);
\node[font=\footnotesize, align=center, text=poppurple!70!black] at (11.3,2.15) {violet\\ (excès)};
% flèche équivalence
\draw[-{Stealth}, very thick, poporange] (6.6,1.1) -- (7.8,1.1);
\node[font=\footnotesize\bfseries, text=poporange, align=center] at (7.2,1.55) {équivalence};
\end{tikzpicture}
\end{center}
\legende{Évolution de la couleur du bécher au cours du dosage : tant que \ce{Fe^2+} est présent, le violet disparaît ; à l'équivalence, la teinte rose persistante apparaît ; au-delà, la solution devient violet foncé.}
\end{popfigure}

\begin{attention}[La teinte rose doit être persistante]
Une coloration rose qui s'atténue et disparaît en quelques secondes signifie qu'il reste encore des ions \ce{Fe^2+} : l'équivalence n'est \textbf{pas} atteinte, il faut continuer à verser. On ne valide le volume équivalent que si la teinte rose-violette persiste \textbf{au moins 30 secondes} en agitant. À l'inverse, si l'on verse trop vite, on dépasse l'équivalence (solution franchement violette) et $V_E$ est surestimé : il faut refaire le dosage.
\end{attention}

\courssub{Exploitation : la relation à l'équivalence}

\begin{methode}[Écrire la relation à l'équivalence — la démarche à suivre]
\begin{enumerate}
\item Écrire l'équation-bilan du dosage et souligner les coefficients stœchiométriques : ici 1 pour \ce{MnO4-} et 5 pour \ce{Fe^2+}.
\item Traduire l'équivalence : les réactifs ont été introduits dans les proportions stœchiométriques, donc
\[ \frac{n(\ce{Fe^2+})_{\text{présent}}}{5} = \frac{n(\ce{MnO4-})_{\text{versé}}}{1} \]
\item En déduire la relation pratique : $n(\ce{Fe^2+}) = 5 \times n(\ce{MnO4-})_E$, soit
\[ C(\ce{Fe^2+}) \times V = 5 \times C(\ce{MnO4-}) \times V_E \quad\Longrightarrow\quad C(\ce{Fe^2+}) = \frac{5 \, C(\ce{MnO4-}) \, V_E}{V} \]
\item Seulement \emph{ensuite}, remplacer par les valeurs numériques (volumes dans la \textbf{même} unité).
\end{enumerate}
\end{methode}

\begin{attention}[L'erreur classique du Probatoire : oublier le facteur 5]
Beaucoup d'élèves écrivent machinalement $C_1 V_1 = C_2 V_2$ comme en dosage acide-base. C'est \textbf{faux} ici : les coefficients stœchiométriques ne sont pas égaux ! Un ion \ce{MnO4-} réagit avec \textbf{cinq} ions \ce{Fe^2+}, donc $n(\ce{Fe^2+}) = 5\, n(\ce{MnO4-})_E$. Oublier ce facteur 5 donne une concentration cinq fois trop petite — et fait perdre la quasi-totalité des points de la question.
\end{attention}

\begin{exoresolu}[Dosage d'une solution de sulfate de fer(II)]
On dose un volume $V = \SI{10,0}{mL}$ d'une solution de sulfate de fer(II) acidifiée par une solution de permanganate de potassium de concentration $C_0 = \SI{0,010}{mol.L^{-1}}$. L'équivalence est repérée à $V_E = \SI{15,8}{mL}$.
\begin{enumerate}
\item Écrire l'équation-bilan du dosage.
\item Calculer la concentration molaire $C$ de la solution en ions \ce{Fe^2+}.
\item En déduire la concentration massique en élément fer de la solution. On donne $M(\ce{Fe}) = \SI{55,8}{g.mol^{-1}}$.
\end{enumerate}
\tcblower
\textbf{1. Équation-bilan.} Couples \ce{MnO4-}/\ce{Mn^2+} et \ce{Fe^3+}/\ce{Fe^2+}, en milieu acide :
\[ \ce{MnO4- + 5 Fe^2+ + 8 H+ -> Mn^2+ + 5 Fe^3+ + 4 H2O} \]

\textbf{2. Concentration molaire.} À l'équivalence, les réactifs sont dans les proportions stœchiométriques :
\[ \frac{n(\ce{Fe^2+})}{5} = n(\ce{MnO4-})_E \quad\Longrightarrow\quad C \times V = 5 \, C_0 \, V_E \]
\[ C = \frac{5 \times C_0 \times V_E}{V} = \frac{5 \times \num{0,010} \times \num{15,8}}{\num{10,0}} = \SI{0,079}{mol.L^{-1}} \]
La solution contient donc $C(\ce{Fe^2+}) = \SI{0,079}{mol.L^{-1}}$.

\textbf{3. Concentration massique en fer.} Chaque ion \ce{Fe^2+} apporte un atome de fer, donc :
\[ t(\ce{Fe}) = C \times M(\ce{Fe}) = 0{,}079 \times 55{,}8 \approx \SI{4,4}{g.L^{-1}} \]
La solution contient environ \SI{4,4}{g} d'élément fer par litre.
\end{exoresolu}

\begin{attention}[Précautions pratiques supplémentaires]
\begin{itemize}
\item Les solutions de \ce{KMnO4} tachent durablement la peau et les vêtements (traces brunes de \ce{MnO2}) : manipuler avec soin, porter blouse et gants.
\item Les ions \ce{Fe^2+} s'oxydent lentement à l'air en \ce{Fe^3+} : la solution à doser doit être préparée et dosée rapidement, sans quoi la concentration mesurée sera trop faible.
\item La solution de permanganate se décompose lentement à la lumière : elle se conserve en flacon ambré et doit être étalonnée régulièrement.
\end{itemize}
\end{attention}

\begin{savaistu}[Le permanganate de potassium, un oxydant aux mille usages]
Le permanganate de potassium ne sert pas qu'aux dosages scolaires ! Grâce à son fort pouvoir oxydant, il est utilisé comme \textbf{désinfectant} (antiseptique dilué en pharmacopée) et pour le \textbf{traitement de l'eau de boisson} : il élimine le fer et le manganèse dissous, détruit de nombreux micro-organismes et oxyde les matières organiques responsables des mauvais goûts. Certaines stations de traitement d'eau au Cameroun y ont recours, en complément de la chloration. C'est aussi le réactif de l'\cle{éthylotest} chimique d'autrefois : le dichromate (puis le permanganate dans certaines versions) changeait de couleur en présence de l'éthanol expiré par un conducteur ayant bu — une application directe de l'oxydoréduction au contrôle routier !
\end{savaistu}

\begin{exercice}[Pour t'entraîner]
Un comprimé anti-anémique est dissous dans de l'eau distillée acidifiée ; le volume est ajusté à \SI{100,0}{mL}. On prélève \SI{20,0}{mL} de cette solution que l'on dose par une solution de \ce{KMnO4} à \SI{0,0050}{mol.L^{-1}}. L'équivalence est obtenue pour $V_E = \SI{12,4}{mL}$. Calculer la masse de fer contenue dans un comprimé. \emph{(Réponse attendue : environ \SI{86}{mg}.)}
\end{exercice}

En résumé, le dosage manganimétrique des ions \ce{Fe^2+} repose sur une réaction totale et rapide, un repérage de l'équivalence sans indicateur grâce à la couleur propre du permanganate, et une relation stœchiométrique dominée par le fameux facteur 5. Dans la section suivante, nous découvrirons un second grand dosage d'oxydoréduction, l'iodométrie, où le repérage de l'équivalence fait cette fois appel à un indicateur externe : l'empois d'amidon.

\coursec{Dosage du diiode par les ions thiosulfate (iodométrie)}

Dans la section précédente, tu as découvert le dosage manganimétrique : l'ion permanganate \ce{MnO4-}, magnifiquement violet, jouait à la fois le rôle de \cle{réactif titrant} et celui d'\cle{indicateur} de fin de réaction, puisqu'à l'équivalence la solution passait de l'incolore au rose persistant. Mais tous les oxydants ne sont pas assez colorés pour s'auto-signaler. Que faire, par exemple, pour déterminer la concentration d'une solution de \cle{diiode} \ce{I2}, comme la teinture d'iode que l'on achète en pharmacie pour désinfecter les plaies ? Le diiode est bien coloré (brun-jaune), mais sa couleur s'estompe progressivement et devient difficile à juger à l'œil nu. Il faut alors un \emph{auxiliaire} : un indicateur coloré ajouté. C'est toute la richesse de l'\cle{iodométrie}, que nous allons étudier pas à pas.

\courssub{Contexte et principe général de l'iodométrie}

Imagine que tu veuilles vérifier l'étiquette d'un flacon de teinture d'iode acheté dans une pharmacie de Douala : la concentration annoncée en diiode est-elle correcte ? Le diiode est un oxydant ; pour le doser, il faut un réducteur qui réagisse avec lui de façon \emph{rapide}, \emph{totale} et \emph{unique}. Le candidat idéal est l'ion \cle{thiosulfate} \ce{S2O3^2-}, fourni par une solution de thiosulfate de sodium \ce{2Na+ + S2O3^2-}, incolore.

\begin{definition}[Iodométrie]
L'\cle{iodométrie} est l'ensemble des dosages d'oxydoréduction mettant en jeu le couple \ce{I2}/\ce{I-}. Dans le cas le plus courant, le diiode \ce{I2} (oxydant, espèce titrée) est dosé par les ions \cle{thiosulfate} \ce{S2O3^2-} (réducteur, espèce titrante), qui sont oxydés en ions \cle{tétrathionate} \ce{S4O6^2-}. L'iodométrie permet aussi des dosages \emph{en retour} : un excès de diiode est introduit, puis l'excès restant est dosé par le thiosulfate.
\end{definition}

\courssub{Les couples mis en jeu et la prévision de la réaction}

Les deux couples d'oxydoréduction concernés sont :

\begin{itemize}
\item le couple \ce{I2}/\ce{I-}, de potentiel normal $E^\circ = \SI{0,54}{V}$ :
\[ \ce{I2 + 2e- <=> 2I-} \]
\item le couple \ce{S4O6^2-}/\ce{S2O3^2-}, de potentiel normal $E^\circ = \SI{0,08}{V}$ :
\[ \ce{S4O6^2- + 2e- <=> 2S2O3^2-} \]
\end{itemize}

Appliquons la \cle{règle du gamma} : on place les deux couples sur un axe vertical de potentiels croissants, l'oxydant le plus fort en haut à gauche, le réducteur le plus fort en bas à droite.

\begin{center}
\begin{tikzpicture}[scale=1]
\draw[->, thick, popdark] (0,0) -- (0,3.6) node[above] {$E^\circ$ (V)};
\draw[thick, popdark] (-0.15,2.7) -- (0.15,2.7) node[right=2pt, popink] {\SI{0,54}{V}};
\draw[thick, popdark] (-0.15,0.6) -- (0.15,0.6) node[right=2pt, popblue] {\SI{0,08}{V}};
\node[anchor=east, popink] at (-0.3,2.7) {\ce{I2}};
\node[anchor=west, popink] at (1.6,2.7) {\ce{I-}};
\node[anchor=east, popblue] at (-0.3,0.6) {\ce{S4O6^2-}};
\node[anchor=west, popblue] at (1.6,0.6) {\ce{S2O3^2-}};
\draw[->, very thick, poporange] (-0.15,2.9) .. controls (-1.4,2.2) and (-1.4,1.1) .. (-0.15,0.4);
\draw[->, very thick, poporange] (1.75,0.5) .. controls (3.2,1.2) and (3.2,2.2) .. (1.75,2.8);
\node[poporange, font=\bfseries] at (3.6,1.65) {$\gamma$};
\end{tikzpicture}
\legende{Règle du gamma : \ce{I2} (oxydant du couple le plus haut) réagit spontanément avec \ce{S2O3^2-} (réducteur du couple le plus bas).}
\end{center}

L'oxydant le plus fort, \ce{I2}, réagit donc avec le réducteur le plus fort, \ce{S2O3^2-}. La réaction est spontanée, quasi totale et rapide à température ambiante : toutes les conditions d'un bon dosage sont réunies.

\begin{aretenir}[Équation de la réaction de dosage]
En combinant la réduction du diiode et l'oxydation du thiosulfate :
\[ \ce{I2 + 2e- -> 2I-} \]
\[ \ce{2S2O3^2- -> S4O6^2- + 2e-} \]
on obtient l'équation-bilan du dosage iodométrique :
\[ \ce{I2 + 2S2O3^2- -> 2I- + S4O6^2-} \]
Réaction \emph{rapide}, \emph{totale} et \emph{unique} à température ambiante, sans catalyseur. Remarque le coefficient \textbf{2} devant \ce{S2O3^2-} : il sera décisif dans les calculs.
\end{aretenir}

\courssub{Dispositif expérimental et repérage de l'équivalence}

Le montage est celui d'un dosage classique : la solution titrante de thiosulfate de sodium, de concentration connue, est placée dans la \cle{burette} ; un volume précis $V$ de solution de diiode, prélevé à la pipette jaugée, est placé dans le bécher (ou l'erlenmeyer) agité.

\begin{popfigure}
\begin{center}
\begin{tikzpicture}[scale=0.95, every node/.style={font=\small}]
% Support
\draw[very thick, popdark] (-3.4,0) -- (-3.4,7.2);
\draw[very thick, popdark] (-4.2,0) -- (-2.6,0);
\draw[thick, popdark] (-3.4,5.6) -- (-2.2,5.6);
% Burette
\draw[thick, popdark] (-2.0,2.6) -- (-2.0,6.9);
\draw[thick, popdark] (-1.4,2.6) -- (-1.4,6.9);
\draw[thick, popdark] (-2.0,6.9) -- (-1.4,6.9);
\fill[popblueL] (-1.97,3.4) rectangle (-1.43,6.85);
\draw[thick, popdark] (-2.0,2.6) -- (-1.85,2.2) -- (-1.55,2.2) -- (-1.4,2.6);
\draw[thick, popdark] (-1.85,2.2) -- (-1.85,1.9);
\draw[thick, popdark] (-1.55,2.2) -- (-1.55,1.9);
\draw[thick, popdark] (-2.15,2.35) -- (-1.25,2.35);
% Goutte
\fill[popblue] (-1.7,1.55) circle (0.07);
\node[align=center, popblue] at (0.4,5.6) {Burette : solution de\\ thiosulfate de sodium\\ \ce{S2O3^2-}, concentration\\ $C_{\text{thio}}$ connue (incolore)};
% Bécher
\draw[very thick, popdark] (-3.1,0.15) -- (-3.1,1.6) -- (-0.3,1.6) -- (-0.3,0.15);
\fill[popgoldL] (-3.05,0.2) rectangle (-0.35,1.05);
\draw[popgold, thick] (-3.05,1.05) -- (-0.35,1.05);
% Barreau magnétique
\fill[popdark] (-2.1,0.28) rectangle (-1.3,0.42);
\node[align=center, popdark] at (1.9,0.9) {Bécher : volume $V$ de\\ solution de \ce{I2} (brune)\\ + empois d'amidon\\ (ajouté en fin de dosage)};
% Agitateur
\draw[thick, popmuted] (-3.6,-0.15) rectangle (0.2,-0.75);
\node[popmuted, font=\footnotesize] at (-1.7,-0.45) {agitateur magnétique};
\end{tikzpicture}
\end{center}
\legende{Montage du dosage iodométrique : le thiosulfate de sodium (burette) titre le diiode (bécher).}
\end{popfigure}

Au fur et à mesure de l'addition de thiosulfate, le diiode est consommé : la teinte \emph{brune} de la solution s'atténue et devient \emph{jaune pâle}. Mais juger la disparition totale de ce jaune pâle est délicat. On utilise alors un indicateur coloré remarquable : l'\cle{empois d'amidon}.

\begin{propriete}[L'empois d'amidon, indicateur du diiode]
L'empois d'amidon forme avec le diiode un complexe d'inclusion d'un \textbf{bleu foncé} (presque noir) très intense, décelable même à très faible concentration en \ce{I2}. En l'absence de diiode, il est \textbf{incolore}. À l'équivalence, la solution passe donc brutalement du bleu foncé à l'incolore : le virage est net et facile à repérer.
\end{propriete}

\begin{methode}[Quand ajouter l'empois d'amidon ?]
\begin{enumerate}
\item Verser le thiosulfate jusqu'à ce que la solution brune devienne \textbf{jaune pâle} (on est alors très proche de l'équivalence).
\item Ajouter seulement \textbf{à ce moment-là} quelques gouttes d'empois d'amidon : la solution vire au \textbf{bleu foncé}.
\item Poursuivre l'addition de thiosulfate \emph{goutte à goutte} jusqu'à la \textbf{décoloration} complète : c'est l'équivalence. Noter le volume versé $V_E$.
\end{enumerate}
\textbf{Pourquoi ne pas ajouter l'amidon dès le début ?} En présence d'une grande quantité de diiode, il se formerait un complexe iode--amidon abondant et \emph{stable}, difficile à détruire : la décoloration serait retardée et l'équivalence serait repérée trop tard, faussant le résultat.
\end{methode}

\begin{popfigure}
\begin{center}
\begin{tikzpicture}[scale=1]
% Etape 1
\fill[popgold!80!black] (0,0) circle (0.55);
\node[white, font=\footnotesize\bfseries] at (0,0) {\ce{I2}};
\node[align=center, font=\footnotesize] at (0,-1.05) {Début :\\ \textbf{brun}};
\draw[->, very thick, popmuted] (0.75,0) -- (1.85,0);
% Etape 2
\fill[popgoldL] (2.6,0) circle (0.55);
\node[align=center, font=\footnotesize] at (2.6,-1.05) {Près de l'équivalence :\\ \textbf{jaune pâle}\\ (ajout de l'amidon)};
\draw[->, very thick, popmuted] (3.35,0) -- (4.45,0);
% Etape 3
\fill[popblue!85!black] (5.2,0) circle (0.55);
\node[align=center, font=\footnotesize] at (5.2,-1.05) {Amidon + \ce{I2} :\\ \textbf{bleu foncé}};
\draw[->, very thick, popmuted] (5.95,0) -- (7.05,0);
% Etape 4
\draw[very thick, popdark] (7.8,0) circle (0.55);
\node[align=center, font=\footnotesize] at (7.8,-1.05) {Équivalence :\\ \textbf{incolore}};
\end{tikzpicture}
\end{center}
\legende{Frise des couleurs observées au cours du dosage iodométrique.}
\end{popfigure}

\courssub{Relation à l'équivalence et exploitation}

À l'\cle{équivalence}, les réactifs ont été mélangés dans les proportions stœchiométriques de l'équation $\ce{I2 + 2S2O3^2- -> 2I- + S4O6^2-}$ : il faut \textbf{2 moles} de thiosulfate pour \textbf{1 mole} de diiode. Donc :

\[ n(\ce{S2O3^2-})_{\text{versé}} = 2 \times n(\ce{I2})_{\text{présent}} \]

En notant $C_{\text{thio}}$ la concentration du thiosulfate, $V_E$ le volume versé à l'équivalence, $C(\ce{I2})$ et $V$ la concentration et le volume de la solution de diiode :

\[ C_{\text{thio}} \times V_E = 2 \times C(\ce{I2}) \times V
\qquad\Longrightarrow\qquad
\boxed{\; C(\ce{I2}) = \dfrac{C_{\text{thio}} \times V_E}{2 \times V} \;} \]

\begin{attention}[Le piège du facteur 2]
L'erreur la plus fréquente au Probatoire est d'\textbf{inverser le facteur 2}. Retiens bien : c'est $n(\ce{S2O3^2-}) = 2\, n(\ce{I2})$, et non $n(\ce{I2}) = 2\, n(\ce{S2O3^2-})$. Astuce de vérification : le thiosulfate est le réactif dont le coefficient stœchiométrique est le plus grand, donc sa quantité à l'équivalence doit être la plus grande. Si ton calcul donne une quantité de diiode supérieure à celle du thiosulfate versé, c'est que tu as inversé le facteur.
\end{attention}

\begin{exoresolu}[Dosage d'une solution de diiode]
\textbf{Énoncé.} On prélève $V = \SI{20,0}{mL}$ d'une solution de diiode (teinture d'iode diluée) que l'on place dans un erlenmeyer. On dose par une solution de thiosulfate de sodium de concentration $C_{\text{thio}} = \SI{0,050}{mol.L^{-1}}$. L'empois d'amidon est ajouté lorsque la solution est jaune pâle ; la décoloration du bleu foncé survient pour $V_E = \SI{16,4}{mL}$.
\begin{enumerate}
\item Écrire l'équation de la réaction de dosage.
\item Calculer la concentration molaire $C(\ce{I2})$ de la solution.
\item En déduire sa concentration massique $C_m$ en diiode. On donne $M(\ce{I2}) = \SI{253,8}{g.mol^{-1}}$.
\end{enumerate}
\tcblower
\textbf{Solution détaillée.}
\begin{enumerate}
\item Les couples sont \ce{I2}/\ce{I-} ($E^\circ = \SI{0,54}{V}$) et \ce{S4O6^2-}/\ce{S2O3^2-} ($E^\circ = \SI{0,08}{V}$). La règle du gamma donne :
\[ \ce{I2 + 2S2O3^2- -> 2I- + S4O6^2-} \]
\item À l'équivalence : $C_{\text{thio}} \times V_E = 2 \times C(\ce{I2}) \times V$, d'où :
\[ C(\ce{I2}) = \frac{C_{\text{thio}} \times V_E}{2 \times V} = \frac{0{,}050 \times 16{,}4}{2 \times 20{,}0} = \SI{0,0205}{mol.L^{-1}} \]
Vérification : $n(\ce{S2O3^2-}) = 0{,}050 \times 16{,}4 \times 10^{-3} = \SI{8,2e-4}{mol}$ et $n(\ce{I2}) = 0{,}0205 \times 20{,}0 \times 10^{-3} = \SI{4,1e-4}{mol}$ : on a bien $n(\ce{S2O3^2-}) = 2\,n(\ce{I2})$.
\item Concentration massique :
\[ C_m = C(\ce{I2}) \times M(\ce{I2}) = 0{,}0205 \times 253{,}8 \approx \SI{5,2}{g.L^{-1}} \]
\end{enumerate}
La solution contient donc environ \SI{5,2}{g} de diiode par litre.
\end{exoresolu}

\courssub{Comparaison des deux dosages d'oxydoréduction étudiés}

\begin{center}
\begin{tabularx}{\linewidth}{|l|X|X|}
\hline
\rowcolor{popblueL} \textbf{Critère} & \textbf{Dosage manganimétrique} & \textbf{Dosage iodométrique} \\
\hline
Espèce titrée & ions \ce{Fe^2+} (réducteur) & diiode \ce{I2} (oxydant) \\
\hline
Espèce titrante & \ce{MnO4-} (oxydant, violet) & \ce{S2O3^2-} (réducteur, incolore) \\
\hline
Milieu & acide (imposé par \ce{MnO4-}) & neutre ou légèrement acide \\
\hline
Repérage de l'équivalence & \textbf{auto-indicateur} : apparition du rose persistant de \ce{MnO4-} & \textbf{indicateur ajouté} : empois d'amidon, virage bleu foncé $\rightarrow$ incolore \\
\hline
Relation à l'équivalence & $n(\ce{Fe^2+}) = 5\, n(\ce{MnO4-})$ & $n(\ce{S2O3^2-}) = 2\, n(\ce{I2})$ \\
\hline
\end{tabularx}
\end{center}

\begin{aretenir}[L'essentiel de l'iodométrie]
\begin{itemize}
\item Équation de dosage : $\ce{I2 + 2S2O3^2- -> 2I- + S4O6^2-}$, rapide et totale.
\item Indicateur : l'empois d'amidon, ajouté \emph{uniquement} quand la solution est jaune pâle ; virage du bleu foncé à l'incolore à l'équivalence.
\item Relation à l'équivalence : $C(\ce{I2}) = \dfrac{C_{\text{thio}} \times V_E}{2 \times V}$.
\item Contrairement au permanganate, le thiosulfate n'est pas auto-indicateur : un indicateur coloré ajouté est indispensable.
\end{itemize}
\end{aretenir}

\begin{savaistu}[La teinture d'iode, un antiseptique contrôlé par iodométrie]
La teinture d'iode vendue dans les pharmacies de Douala, Yaoundé ou Garoua est une solution de diiode (souvent en présence d'iodure de potassium, qui augmente la solubilité de \ce{I2} en formant l'ion \ce{I3-}) utilisée comme \emph{antiseptique} pour désinfecter les plaies. Son efficacité dépend directement de sa teneur en diiode : les laboratoires de contrôle qualité vérifient ce titre précisément par iodométrie, avec du thiosulfate de sodium et de l'empois d'amidon. L'amidon utilisé pour fabriquer cet indicateur n'est autre que celui que l'on extrait du manioc ou du maïs, produits emblématiques de l'agriculture camerounaise !
\end{savaistu}

\begin{exercice}[À toi de jouer]
On dose $V = \SI{10,0}{mL}$ d'une solution de diiode par une solution de thiosulfate de sodium de concentration $C_{\text{thio}} = \SI{0,020}{mol.L^{-1}}$. L'équivalence est obtenue pour $V_E = \SI{12,5}{mL}$. Calcule la concentration molaire puis la concentration massique de la solution en diiode.
\end{exercice}

\begin{corrige}[Exercice]
À l'équivalence : $C(\ce{I2}) = \dfrac{C_{\text{thio}} \times V_E}{2 \times V} = \dfrac{0{,}020 \times 12{,}5}{2 \times 10{,}0} = \SI{0,0125}{mol.L^{-1}}$. Concentration massique : $C_m = 0{,}0125 \times 253{,}8 \approx \SI{3,2}{g.L^{-1}}$.
\end{corrige}

\coursec{Méthodologie, exploitation des résultats et entraînement type Probatoire}

Dans les sections précédentes, tu as appris à réaliser concrètement deux dosages d'oxydoréduction : le dosage manganimétrique des ions \ce{Fe^2+} par les ions permanganate \ce{MnO4-}, puis l'iodométrie, c'est-à-dire le dosage du diiode \ce{I2} par les ions thiosulfate \ce{S2O3^2-}. Tu sais donc manipuler la burette, repérer l'équivalence par un changement de couleur et écrire la relation fondamentale entre les quantités de matière. Reste maintenant l'étape décisive : celle qui transforme un élève qui « sait faire » en un candidat qui \emph{réussit} au Probatoire. Cette dernière section te donne une méthode infaillible, en cinq étapes, pour traiter \emph{n'importe quel} exercice de dosage, puis elle t'entraîne sur les pièges les plus fréquents et les exigences de rédaction de l'examen.

\courssub{La démarche complète en cinq étapes}

Face à un exercice de dosage, le candidat pressé saute directement au calcul. C'est la première erreur. Un dosage se résout comme on construit une maison : d'abord les fondations (les couples et les demi-équations), puis les murs (l'équation-bilan), puis le toit (la relation à l'équivalence), et seulement à la fin la décoration (l'application numérique). Voici la démarche complète, à apprendre par cœur.

\begin{methode}[Résoudre un exercice de dosage d'oxydoréduction en cinq étapes]
\begin{enumerate}
\item \textbf{Identifier les deux couples mis en jeu.} Repérer dans l'énoncé l'oxydant et le réducteur, puis écrire les deux couples sous la forme Ox/Red : par exemple \ce{MnO4-}/\ce{Mn^2+} et \ce{Fe^3+}/\ce{Fe^2+}. Vérifier leur position dans la classification électrochimique (l'oxydant le plus fort réagit avec le réducteur le plus fort, règle du gamma).
\item \textbf{Écrire et équilibrer les deux demi-équations électroniques.} Équilibrer d'abord les atomes autres que O et H, puis l'oxygène avec \ce{H2O}, puis l'hydrogène avec \ce{H+}, puis la charge avec les électrons.
\item \textbf{Établir l'équation-bilan.} Multiplier chaque demi-équation par le coefficient qui égalise les électrons échangés, additionner, puis simplifier. Vérifier la conservation des atomes \emph{et} des charges.
\item \textbf{Écrire la relation à l'équivalence.} À l'équivalence, les réactifs ont été mélangés dans les proportions stœchiométriques : si l'équation s'écrit $a\,\text{Ox}_1 + b\,\text{Red}_2 \ce{->} \text{produits}$, alors
\[ \frac{n(\text{Ox}_1)_{\text{versé}}}{a} = \frac{n(\text{Red}_2)_{\text{dosé}}}{b} \]
\item \textbf{Calculer, avec les unités.} Remplacer $n$ par $C \times V$, convertir les volumes en litres, isoler l'inconnue, calculer, puis exprimer le résultat avec trois chiffres significatifs et son unité.
\end{enumerate}
\end{methode}

\begin{attention}[Les volumes : toujours en litres !]
Dans la relation $n = C \times V$, la concentration $C$ s'exprime en \si{mol.L^{-1}} : le volume $V$ doit donc être converti en \textbf{litres} avant tout calcul. Un volume d'équivalence $V_E = \SI{12,6}{mL}$ s'écrit $V_E = \SI{0,0126}{L}$. Oublier cette conversion multiplie ton résultat par $1000$ : c'est l'erreur numéro un des copies de Probatoire. Prends l'habitude d'écrire la conversion \emph{explicitement} sur ta copie.
\end{attention}

\begin{popfigure}
\begin{tikzpicture}[
node distance=0.55cm,
etape/.style={rectangle, rounded corners=3pt, draw=popblue, fill=popblueL, text width=10.2cm, align=left, font=\small, inner sep=6pt},
fleche/.style={-{Stealth[length=2.5mm]}, thick, popdark}
]
\node[etape] (e1) {\textbf{Étape 1 — Identifier les couples} : oxydant titrant, réducteur titré ; vérifier la spontanéité par la règle du gamma.};
\node[etape, below=of e1] (e2) {\textbf{Étape 2 — Demi-équations} : équilibrer atomes, puis O avec \ce{H2O}, puis H avec \ce{H+}, puis charges avec $e^-$.};
\node[etape, below=of e2] (e3) {\textbf{Étape 3 — Équation-bilan} : égaliser les électrons, additionner, simplifier, vérifier atomes et charges.};
\node[etape, below=of e3] (e4) {\textbf{Étape 4 — Relation à l'équivalence} : $\dfrac{n_{\text{Ox}}}{a} = \dfrac{n_{\text{Red}}}{b}$ (proportions stœchiométriques).};
\node[etape, below=of e4, draw=popgreen, fill=popgreenL] (e5) {\textbf{Étape 5 — Calcul} : $n = C \times V$ avec $V$ en litres ; résultat avec 3 chiffres significatifs et unité.};
\draw[fleche] (e1) -- (e2);
\draw[fleche] (e2) -- (e3);
\draw[fleche] (e3) -- (e4);
\draw[fleche] (e4) -- (e5);
\end{tikzpicture}
\legende{Organigramme de la démarche de résolution d'un exercice de dosage d'oxydoréduction.}
\end{popfigure}

\courssub{Le schéma normalisé du dispositif expérimental}

Au Probatoire, on te demande souvent de « schématiser et légender le dispositif de dosage ». Un schéma non légendé ne rapporte qu'une fraction des points. La légende doit impérativement préciser :

\begin{itemize}
\item la \textbf{nature} de chaque solution (ex. « solution de permanganate de potassium acidifiée ») ;
\item sa \textbf{concentration} quand elle est connue (ex. $C_1 = \SI{0,020}{mol.L^{-1}}$) ;
\item le \textbf{volume} prélevé à la pipette jaugée (ex. $V_2 = \SI{20,0}{mL}$) ;
\item le rôle de chaque verrerie : burette graduée (solution titrante), erlenmeyer ou bécher (solution titrée), agitateur magnétique.
\end{itemize}

Rappelle-toi : la solution \emph{titrante} (de concentration connue) va dans la burette ; la solution \emph{titrée} (de concentration inconnue) est prélevée à la pipette jaugée et placée dans l'erlenmeyer. Pour le dosage manganimétrique, l'ion \ce{MnO4-} joue en plus le rôle d'indicateur de fin de réaction : avant l'équivalence, la solution reste incolore (le permanganate est consommé instantanément) ; à l'équivalence, la première goutte en excès donne une coloration \cle{rose persistante}.

\courssub{Détermination précise du volume équivalent : les essais concordants}

Une seule mesure de $V_E$ ne suffit jamais : une goutte de trop et le volume est faussé. La démarche expérimentale correcte comporte deux phases.

\begin{experience}[Détermination de $V_E$ par essais concordants]
\textbf{Essai grossier (essai repère).} On verse la solution titrante \emph{millilitre par millilitre} en agitant, jusqu'au changement de couleur. On note le volume obtenu, par exemple $V_E \approx \SI{13}{mL}$. Cet essai, volontairement rapide, sert uniquement à localiser la zone de l'équivalence.

\textbf{Essais précis.} On refait le dosage en versant rapidement jusqu'à environ $\SI{1}{mL}$ avant le volume repère (ici jusqu'à $\SI{12}{mL}$), puis \emph{goutte à goutte}, en rinçant les parois de l'erlenmeyer à l'eau distillée si nécessaire. On note le volume à la goutte près. On recommence jusqu'à obtenir \textbf{deux ou trois valeurs concordantes}, c'est-à-dire écartées d'au plus $\SI{0,2}{mL}$.

\textbf{Exploitation.} On retient la moyenne des essais concordants (on écarte l'essai grossier et tout essai aberrant).
\end{experience}

\begin{exemplebox}[Tableau des volumes et calcul de $V_E$]
Lors du dosage de $V_2 = \SI{20,0}{mL}$ d'une solution de \ce{Fe^2+} par une solution de \ce{KMnO4} à $C_1 = \SI{0,020}{mol.L^{-1}}$, un élève consigne :

\begin{center}
\begin{tabularx}{\linewidth}{|l|X|X|X|X|}
\hline
\rowcolor{popblueL} Essai & Grossier & Précis n\textsuperscript{o}1 & Précis n\textsuperscript{o}2 & Précis n\textsuperscript{o}3 \\
\hline
$V_E$ (\si{mL}) & 13,0 & 12,8 & 12,6 & 12,7 \\
\hline
\end{tabularx}
\end{center}

Les trois essais précis sont concordants (écart maximal $\SI{0,2}{mL}$). On retient :
\[ V_E = \frac{12,8 + 12,6 + 12,7}{3} = \SI{12,7}{mL} = \SI{0,0127}{L} \]
\end{exemplebox}

\courssub{De la concentration molaire à la concentration massique}

Certains énoncés demandent la \cle{concentration massique} $C_m$ (en \si{g.L^{-1}}), par exemple pour comparer une teneur en fer à une norme. Le passage est immédiat :

\begin{aretenir}[Concentration massique]
\[ C_m = C \times M \]
avec $C$ la concentration molaire (\si{mol.L^{-1}}), $M$ la masse molaire de l'espèce dosée (\si{g.mol^{-1}}) et $C_m$ la concentration massique (\si{g.L^{-1}}). Attention : $M$ est la masse molaire de l'\emph{espèce dosée} (par exemple $M(\text{Fe}) = \SI{55,8}{g.mol^{-1}}$ si l'on veut la teneur en fer), pas celle du sel dissous.
\end{aretenir}

\begin{exemplebox}[Teneur en fer d'une solution]
Si le dosage donne $C(\ce{Fe^2+}) = \SI{0,079}{mol.L^{-1}}$, la concentration massique en fer vaut :
\[ C_m = 0,079 \times 55,8 = \SI{4,41}{g.L^{-1}} \]
Soit environ $\SI{4,4}{g}$ de fer par litre de solution.
\end{exemplebox}

\courssub{Le piège de la dilution : remonter à la solution mère}

Très souvent, l'échantillon à doser est trop concentré pour être dosé directement : on le \textbf{dilue} avant le dosage. L'énoncé dit alors, par exemple : « on dilue 10 fois la solution commerciale ; on dose $\SI{20,0}{mL}$ de la solution diluée ». Le calcul fournit la concentration de la solution \emph{diluée} ; il faut ensuite \emph{remonter} à la solution mère en multipliant par le facteur de dilution.

\begin{attention}[Erreur fréquente au Probatoire : oublier le facteur de dilution]
Beaucoup de candidats calculent correctement la concentration de la solution diluée... et s'arrêtent là, alors que la question porte sur la solution mère (commerciale, initiale). Si la solution a été diluée $F$ fois, alors :
\[ C_{\text{mère}} = F \times C_{\text{diluée}} \]
avec $F = \dfrac{V_{\text{fiole}}}{V_{\text{prélevé}}}$. Souligne dans l'énoncé les mots « diluée », « solution mère », « solution commerciale » : ils t'indiquent qu'un facteur de dilution se cache quelque part.
\end{attention}

\begin{exemplebox}[Solution de \ce{Fe^2+} diluée dix fois]
Une solution de sulfate de fer(II) est diluée 10 fois ($\SI{10,0}{mL}$ prélevés, complétés à $\SI{100,0}{mL}$). Le dosage manganimétrique de la solution diluée donne :
\[ C_{\text{diluée}} = \SI{0,0079}{mol.L^{-1}} \]
La concentration de la solution mère est donc :
\[ C_{\text{mère}} = 10 \times 0,0079 = \SI{0,079}{mol.L^{-1}} \]
Oublier le facteur 10, c'est donner une réponse dix fois trop faible : zéro point sur la question, même si tout le reste est juste.
\end{exemplebox}

\courssub{Compléments utiles au Probatoire (hors programme strict)}

Les deux notions qui suivent ne figurent pas explicitement dans le programme officiel de Première, mais elles apparaissent régulièrement dans les sujets et te donnent une longueur d'avance.

\textbf{Le degré de pureté.} Lorsqu'on dose un échantillon solide (par exemple des cristaux de sulfate de fer(II) partiellement oxydés à l'air), on compare la masse d'espèce pure trouvée par le dosage à la masse totale de l'échantillon :
\[ \text{degré de pureté} = \frac{m_{\text{espèce pure}}}{m_{\text{échantillon}}} \times 100 \quad (\text{en } \%) \]
Par exemple, si le dosage montre qu'un échantillon de $\SI{2,50}{g}$ de « sulfate de fer(II) » contient en réalité $\SI{2,30}{g}$ de \ce{FeSO4}, le degré de pureté vaut $\frac{2,30}{2,50} \times 100 = \SI{92}{\%}$.

\textbf{Incertitude simple sur $V_E$.} Le volume équivalent est lu à la burette avec une précision d'environ une goutte, soit $\pm \SI{0,05}{mL}$ à $\pm \SI{0,1}{mL}$. On écrit alors le résultat sous la forme $V_E = (12,7 \pm 0,1)\,\si{mL}$. Cela justifie l'exigence des essais concordants : deux essais écartés de plus de $\SI{0,2}{mL}$ signalent une erreur de manipulation (goutte de trop, agitation insuffisante, lecture du ménisque biaisée).

\textbf{Chiffres significatifs.} Au Probatoire, les données comportent en général trois chiffres significatifs ($C_1 = \SI{0,0200}{mol.L^{-1}}$, $V_2 = \SI{20,0}{mL}$) : ton résultat final doit donc être donné avec \textbf{trois chiffres significatifs}. Écris $C_2 = \SI{0,0127}{mol.L^{-1}}$, et non $C_2 = \SI{0,012698}{mol.L^{-1}}$ (précision illusoire) ni $C_2 = \SI{0,01}{mol.L^{-1}}$ (perte d'information). Garde les valeurs intermédiaires non arrondies dans ta calculatrice et n'arrondis qu'à la toute fin.

\begin{savaistu}[Combien vaut un exercice de dosage au Probatoire ?]
Aux Probatoires C et D, un exercice de dosage d'oxydoréduction rapporte en moyenne \textbf{4 à 6 points} sur 20 en chimie. La répartition typique est la suivante : demi-équations et équation-bilan (2 points), relation à l'équivalence correctement écrite (1 point à elle seule, même sans calcul !), application numérique avec unités (1 à 2 points), schéma légendé ou question de réflexion (1 point). Moralité : même si tu doutes de ton calcul, écris \emph{toujours} la relation à l'équivalence et l'équation-bilan : ce sont des points garantis.
\end{savaistu}

\courssub{Exercice résolu type Probatoire}

\begin{exoresolu}[Dosage du fer(II) dans un anti-anémique]
Pour traiter l'anémie, très répandue chez les femmes enceintes au Cameroun, on prescrit des comprimés de sulfate de fer(II). On dissout un comprimé dans de l'eau distillée acidifiée et on complète à $\SI{200,0}{mL}$ (solution S). On prélève $V_2 = \SI{20,0}{mL}$ de S que l'on dose par une solution acidifiée de permanganate de potassium de concentration $C_1 = \SI{0,0100}{mol.L^{-1}}$. L'équivalence (coloration rose persistante) est obtenue pour $V_E = \SI{11,4}{mL}$.

1. Écrire les demi-équations et l'équation-bilan du dosage.
2. Écrire la relation à l'équivalence et calculer la concentration $C_2$ en ions \ce{Fe^2+} de la solution S.
3. En déduire la masse de fer contenue dans un comprimé. On donne $M(\text{Fe}) = \SI{55,8}{g.mol^{-1}}$.

\tcblower
\textbf{1. Couples et équation-bilan.} Couples mis en jeu : \ce{MnO4-}/\ce{Mn^2+} ($E^\circ = \SI{1,51}{V}$) et \ce{Fe^3+}/\ce{Fe^2+} ($E^\circ = \SI{0,77}{V}$). D'après la règle du gamma, \ce{MnO4-} (oxydant le plus fort) oxyde \ce{Fe^2+} (réducteur le plus fort) : la réaction est spontanée, totale et rapide en milieu acide.

Demi-équations :
\[ \ce{MnO4- + 8H+ + 5e- -> Mn^2+ + 4H2O} \qquad \ce{Fe^2+ -> Fe^3+ + e-} \]
On multiplie la seconde par 5 pour égaliser les électrons, puis on additionne :
\[ \ce{MnO4- + 5Fe^2+ + 8H+ -> Mn^2+ + 5Fe^3+ + 4H2O} \]
Vérification : charges $-1 + 10 + 8 = +17$ à gauche, $+2 + 15 = +17$ à droite ; atomes équilibrés. La réaction est totale, rapide et unique : elle convient parfaitement à un dosage.

\textbf{2. Relation à l'équivalence et calcul de $C_2$.} À l'équivalence, les réactifs sont dans les proportions stœchiométriques :
\[ \frac{n(\ce{MnO4-})}{1} = \frac{n(\ce{Fe^2+})}{5} \quad \Longrightarrow \quad C_1 V_E = \frac{C_2 V_2}{5} \quad \Longrightarrow \quad C_2 = \frac{5\,C_1 V_E}{V_2} \]
Application numérique (volumes en litres, mais le rapport $V_E/V_2$ permet de garder les \si{mL}) :
\[ C_2 = \frac{5 \times 0,0100 \times 11,4}{20,0} = \SI{0,0285}{mol.L^{-1}} \]

\textbf{3. Masse de fer par comprimé.} La solution S ($\SI{200,0}{mL}$) contient :
\[ n(\ce{Fe^2+}) = C_2 \times V_S = 0,0285 \times 0,2000 = \SI{5,70e-3}{mol} \]
\[ m(\text{Fe}) = n \times M = 5,70 \times 10^{-3} \times 55,8 = \SI{0,318}{g} \]
Soit environ $\SI{318}{mg}$ de fer par comprimé, valeur cohérente avec les posologies usuelles.
\end{exoresolu}

\courssub{Les erreurs classiques à corriger impérativement}

Voici la « galerie des horreurs » relevée dans les copies, avec le remède pour chacune.

\begin{center}
\begin{tabularx}{\linewidth}{|X|X|}
\hline
\rowcolor{poporangeL} \textbf{Erreur classique} & \textbf{Correction} \\
\hline
Équation-bilan non équilibrée (électrons non égalisés) & Multiplier chaque demi-équation par le bon coefficient avant d'additionner ; vérifier atomes \emph{et} charges. \\
\hline
Facteur stœchiométrique oublié : écrire $C_1V_E = C_2V_2$ alors que l'équation est $\ce{MnO4-} + 5\ce{Fe^2+}$ & Toujours écrire $\dfrac{n_1}{a} = \dfrac{n_2}{b}$ à partir de l'équation-bilan. Ici $C_2 = \dfrac{5C_1V_E}{V_2}$. \\
\hline
Volumes laissés en \si{mL} dans $n = C \times V$ & Convertir en litres : $V_E = \SI{12,7}{mL} = \SI{0,0127}{L}$. \\
\hline
Oublier de remonter à la solution mère après dilution & Multiplier par le facteur de dilution $F$ : $C_{\text{mère}} = F \times C_{\text{diluée}}$. \\
\hline
Confondre solution titrante et solution titrée dans la relation & Titrante = burette, concentration connue ; titrée = erlenmeyer, concentration inconnue. \\
\hline
Résultat sans unité ou avec 8 chiffres & Trois chiffres significatifs et une unité : $C_2 = \SI{0,0285}{mol.L^{-1}}$. \\
\hline
\end{tabularx}
\end{center}

\begin{exercice}[Entraînement type Probatoire : iodométrie]
On dose le diiode présent dans $V_2 = \SI{25,0}{mL}$ d'une solution de Lugol (antiseptique) par une solution de thiosulfate de sodium de concentration $C_1 = \SI{0,0500}{mol.L^{-1}}$, en présence d'empois d'amidon. La décoloration (disparition du bleu) survient pour $V_E = \SI{14,2}{mL}$.

1. Écrire l'équation-bilan du dosage (couples \ce{I2}/\ce{I-} et \ce{S4O6^2-}/\ce{S2O3^2-}).
2. Établir la relation à l'équivalence et calculer $C(\ce{I2})$.
3. Calculer la concentration massique en diiode ($M(\ce{I2}) = \SI{253,8}{g.mol^{-1}}$).
\end{exercice}

\begin{corrige}[Exercice : iodométrie]
\textbf{1.} Demi-équations : $\ce{I2 + 2e- -> 2I-}$ et $\ce{2S2O3^2- -> S4O6^2- + 2e-}$. Équation-bilan :
\[ \ce{I2 + 2S2O3^2- -> 2I- + S4O6^2-} \]
Réaction totale, rapide, sans chauffage ni catalyseur ; l'empois d'amidon sert d'indicateur (bleu foncé avec \ce{I2}, incolore à l'équivalence).

\textbf{2.} À l'équivalence : $\dfrac{n(\ce{S2O3^2-})}{2} = \dfrac{n(\ce{I2})}{1}$, d'où $C(\ce{I2}) = \dfrac{C_1 V_E}{2 V_2} = \dfrac{0,0500 \times 14,2}{2 \times 25,0} = \SI{0,0142}{mol.L^{-1}}$.

\textbf{3.} $C_m = C \times M = 0,0142 \times 253,8 = \SI{3,60}{g.L^{-1}}$.
\end{corrige}

\courssub{Bilan de la leçon : tout est relié}

Cette leçon t'a mené de la classification électrochimique aux dosages, en passant par la prévision des réactions. La logique est implacable et belle : les potentiels normaux $E^\circ$ permettent de \emph{prévoir} quelle réaction aura lieu (règle du gamma) ; cette réaction, totale et rapide, devient le \emph{support} d'un dosage ; et la stœchiométrie de cette réaction fournit la \emph{relation à l'équivalence}, clé de tous les calculs. La carte mentale ci-dessous résume l'ensemble : garde-la sous les yeux pour réviser.

\begin{popfigure}
\begin{tikzpicture}[
centre/.style={rectangle, rounded corners=4pt, draw=popdark, fill=popblue, text=white, font=\bfseries\small, align=center, inner sep=7pt},
branche/.style={rectangle, rounded corners=3pt, draw=popblue, fill=popblueL, font=\small, align=center, text width={4.6cm}, inner sep=5pt},
feuille/.style={rectangle, rounded corners=3pt, draw=popgreen, fill=popgreenL, font=\footnotesize, align=center, text width={4.2cm}, inner sep=4pt},
lien/.style={thick, popdark}
]
\node[centre] (c) at (0,0) {Oxydoréduction\\ et dosages};
\node[branche] (b1) at (-3.4,2.1) {Classification\\ électrochimique};
\node[branche] (b2) at (3.4,2.1) {Prévision\\ des réactions};
\node[branche] (b3) at (-3.4,-2.1) {Dosages\\ d'oxydoréduction};
\node[branche] (b4) at (3.4,-2.1) {Exploitation\\ des résultats};
\node[feuille, anchor=south] (f1) at (-3.4,3.0) {Couples au programme :\\ \ce{MnO4-}/\ce{Mn^2+}, \ce{Cr2O7^2-}/\ce{Cr^3+},\\ \ce{Fe^3+}/\ce{Fe^2+}, \ce{I2}/\ce{I-},\\ \ce{S4O6^2-}/\ce{S2O3^2-}, \ce{Cl2}/\ce{Cl-},\\ \ce{NO3-}/\ce{NO}, \ce{SO4^2-}/\ce{SO2}};
\node[feuille, anchor=south] (f2) at (3.4,3.0) {Règle du gamma :\\ l'oxydant le plus fort ($E^\circ$ élevé)\\ réagit avec le réducteur\\ le plus fort ($E^\circ$ faible)};
\node[feuille, anchor=north] (f3) at (-3.4,-3.0) {Réaction totale, rapide, unique ;\\ équivalence repérée par\\ changement de couleur\\ (\ce{MnO4-} rose, amidon bleu)};
\node[feuille, anchor=north] (f4) at (3.4,-3.0) {À l'équivalence : $\dfrac{n_1}{a}=\dfrac{n_2}{b}$ ;\\ $n = C \times V$ ($V$ en L) ;\\ $C_m = C \times M$ ;\\ $C_{\text{mère}} = F \times C_{\text{diluée}}$};
\draw[lien] (c) -- (b1);
\draw[lien] (c) -- (b2);
\draw[lien] (c) -- (b3);
\draw[lien] (c) -- (b4);
\draw[lien] (b1) -- (f1);
\draw[lien] (b2) -- (f2);
\draw[lien] (b3) -- (f3);
\draw[lien] (b4) -- (f4);
\end{tikzpicture}
\legende{Carte mentale de synthèse : de la classification électrochimique aux calculs de dosage.}
\end{popfigure}

\begin{aretenir}[L'essentiel de la section]
\begin{itemize}
\item Un exercice de dosage se résout en \textbf{cinq étapes} : couples $\rightarrow$ demi-équations $\rightarrow$ équation-bilan $\rightarrow$ relation à l'équivalence $\rightarrow$ calcul.
\item À l'équivalence : $\dfrac{n_{\text{Ox}}}{a} = \dfrac{n_{\text{Red}}}{b}$, avec $n = C \times V$ et $V$ \textbf{en litres}.
\item $V_E$ se détermine par un essai grossier puis des essais précis \textbf{concordants} (écart $\leq \SI{0,2}{mL}$) dont on fait la moyenne.
\item Concentration massique : $C_m = C \times M$ ; solution diluée : $C_{\text{mère}} = F \times C_{\text{diluée}}$.
\item Résultat final : \textbf{trois chiffres significatifs} et une unité. La relation à l'équivalence, même seule, rapporte des points : écris-la toujours.
\end{itemize}
\end{aretenir}

Tu es désormais armé pour affronter n'importe quel sujet de dosage d'oxydoréduction au Probatoire. Entraîne-toi sur les exercices de la plateforme : la méthode des cinq étapes doit devenir un réflexe, aussi automatique que le geste de lire le ménisque à la burette. Bon courage, et rendez-vous à la prochaine leçon !

\newpage

\coursec{Travaux pratiques}

\courssub{Objectifs du TP}
\begin{itemize}
    \item Mettre en œuvre une technique de \cle{dosage d'oxydoréduction} (manganimétrie) pour déterminer la concentration molaire inconnue d'une solution d'ions fer(II) \ce{Fe^2+}.
    \item Identifier le \cle{point d'équivalence} grâce au changement de couleur intrinsèque du réactif titrant (permanganate), sans indicateur ajouté.
    \item Maîtriser les manipulations de base : rinçage et remplissage de la burette, jaugeage à la pipette, agitation efficace, lecture du ménisque.
    \item Écrire les demi-équations électroniques, l'équation-bilan de la réaction et établir le tableau d'avancement pour calculer la concentration recherchée.
    \item Évaluer la précision des résultats (essais concordants) et discuter les sources d'erreurs expérimentales.
    \item Connaître l'alternative iodométrique (dosage du diiode par le thiosulfate) si le permanganate fait défaut.
\end{itemize}

\courssub{Matériel et produits}
\begin{center}
\begin{tabularx}{\linewidth}{|l|X|l|}\hline
\rowcolor{popblueL}
\textbf{Catégorie} & \textbf{Désignation} & \textbf{Quantité / Concentration} \\\hline
Verrerie & Burette graduée (classe B) & \SI{25}{mL} (une par binôme) \\\hline
& Pipette jaugée (classe B) & \SI{10,0}{mL} (une par binôme) \\\hline
& Bécher (ou erlenmeyer) & \SI{100}{mL} ou \SI{150}{mL} (deux) \\\hline
& Entonnoir en verre & 1 \\\hline
& Fiole jaugée (pour préparation \ce{Fe^2+}) & \SI{100}{mL} ou \SI{250}{mL} \\\hline
& Bagotte agitateur / Barre magnétique + agitateur & 1 \\\hline
Produits & Solution de permanganate de potassium \ce{KMnO4} & \SI{0,0200}{mol.L^{-1}} (\approx \SI{3,16}{g.L^{-1}}) \\\hline
& Sulfate de fer(II) heptahydraté \ce{FeSO4,7H2O} (solide) & \SI{3,0}{g} (pour \SI{100}{mL} \SI{0,10}{mol.L^{-1}}) \\\hline
& Acide sulfurique dilué \ce{H2SO4} & \SI{1,0}{mol.L^{-1}} (\approx \SI{55}{mL} par titration) \\\hline
& Eau distillée & En quantité suffisante \\\hline
Sécurité & Blouse, lunettes de protection, gants nitrile & Obligatoires \\\hline
\multicolumn{3}{|l|}{\textbf{Alternatives locales (contexte Cameroun) :}} \\\hline
\multicolumn{3}{|X|}{Si \ce{KMnO4} absent : Teinture d'iode (pharmacie, \approx \SI{0,1}{mol.L^{-1}} en \ce{I2}) diluée 1/10 + Solution de thiosulfate de sodium \ce{Na2S2O3} \SI{0,050}{mol.L^{-1}} + Amidon de maïs (maïzena) bouilli comme indicateur.} \\\hline
\multicolumn{3}{|X|}{Si pipette/burette manquantes : Seringue graduée \SI{10}{mL} (pipette) + Burette de Mohr en plastique ou burette en verre récupérée (vérifier l'étanchéité du robinet).} \\\hline
\end{tabularx}
\end{center}

\courssub{Sécurité}
\begin{attention}[Consignes de sécurité impératives - Pictogrammes SGH]
\textbf{Acide sulfurique \ce{H2SO4} \SI{1}{mol.L^{-1}} :} Pictogramme \textbf{GHS05 (Corrosif)}. Risque de brûlures graves oculaires et cutanées. \textbf{Gestes :} Porter gants, lunettes, blouse. En cas de projection : rincer \textbf{immédiatement} à l'eau courante pendant 15 min (œil : rincer paupières ouvertes). Ne jamais verser l'eau dans l'acide (éclaboussures violentes) : verser l'acide dans l'eau. Neutraliser les petits déversements avec du bicarbonate de sodium \ce{NaHCO3} (effervescence \ce{CO2}).

\textbf{Permanganate de potassium \ce{KMnO4} \SI{0,02}{mol.L^{-1}} :} Pictogrammes \textbf{GHS03 (Comburant)}, \textbf{GHS07 (Danger)}, \textbf{GHS09 (Dangereux pour l'environnement)}. Fort pouvoir colorant (taches brunes \ce{MnO2} tenaces sur peau, tissu, paillasse). \textbf{Gestes :} Gants obligatoires. Ne pas chauffer au contact de matières organiques (risque d'inflammation). Éliminer les résidus dans le bac « Oxydants/Metaux lourds », pas dans l'évier.

\textbf{Sulfate de fer(II) \ce{FeSO4,7H2O} :} Pictogramme \textbf{GHS07 (Danger)}. Nocif par ingestion, irritant. S'oxyde rapidement à l'air en \ce{Fe^3+} (précipité brun \ce{Fe(OH)3}) : \textbf{préparer la solution juste avant le TP} et la boucher.

\textbf{Règles générales :} Pas de manger/boire au labo. Laver les mains à la fin. Connaître l'emplacement de la douche de sécurité et du lave-œil.
\end{attention}

\courssub{Schéma du dispositif}
\begin{popfigure}
\begin{tikzpicture}[scale=0.9, >=Stealth]
% Support (pied + tige)
\draw[popdark, line width=3pt] (0,0) -- (0,-0.5) -- (6,-0.5) -- (6,0); % Pied
\draw[popdark, line width=2pt] (5.5,0) -- (5.5,8); % Tige verticale
\draw[popdark, line width=2pt] (5.5,7.5) -- (3.5,7.5); % Bras horizontal

% Burette
\draw[popblueL, line width=1.5pt] (3.5, 7.5) -- (3.5, 6.8) -- (3.0, 6.8) -- (3.0, 1.5) -- (3.5, 1.5) -- (3.5, 1.0) -- (4.0, 1.0) -- (4.0, 1.5) -- (4.5, 1.5) -- (4.5, 6.8) -- (4.0, 6.8) -- (4.0, 7.5) -- cycle;
% Graduations burette
\foreach \y in {1.5, 2.0, ..., 6.5} {
    \draw[popdark, thin] (3.0, \y) -- (3.3, \y);
    \draw[popdark, very thin] (4.5, \y) -- (4.2, \y);
}
\foreach \y/\label in {1.5/0, 2.5/10, 3.5/20, 4.5/30, 5.5/40, 6.5/50} {
    \node[popdark, left, font=\tiny] at (2.9, \y) {\label};
}
% Niveau liquide initial (violet)
\fill[poppurple!60] (3.05, 1.55) rectangle (4.45, 5.8);
% Menisque
\draw[poppurple, thick, domain=3.05:4.45, smooth, variable=\x] plot ({\x}, {5.8 + 0.05*sin((\x-3.05)*180/1.4)});
% Robinet
\draw[popdark, line width=2pt] (4.5, 1.0) -- (5.0, 1.0);
\draw[popdark, line width=2pt] (4.75, 1.2) -- (4.75, 0.6) -- (5.25, 0.6) -- (5.25, 1.2) -- cycle; % Poignée
% Gouttes qui tombent
\fill[poppurple] (4.75, 0.5) circle (2pt);
\fill[poppurple] (4.75, 0.2) circle (1.5pt);

% Erlenmeyer / Bécher sous la burette
\draw[popblueL, line width=1.5pt] (2.5, 0) -- (2.5, 3.5) -- (5.5, 3.5) -- (5.5, 0) -- cycle; % Bécher simple
% Solution incolore Fe2+ + H2SO4
\fill[popgreenL!50!white] (2.55, 0.05) rectangle (5.45, 2.8);
% Surface liquide
\draw[popdark, thick] (2.55, 2.8) -- (5.45, 2.8);

% Agitateur magnétique (barre)
\draw[poporange, line width=2pt, rotate around={15:(4.0,0.8)}] (4.0, 0.8) ellipse (0.4 and 0.1);
% Plaque chauffante/agitatrice
\draw[popdark, fill=popmuted!20, rounded corners=2pt] (2.0, -0.3) rectangle (6.0, -0.1);

% Feuille blanche pour contraste
\draw[popcream, line width=1pt] (2.0, -0.5) rectangle (6.0, -0.3);
\node[popdark, font=\tiny, align=center] at (4.0, -0.4) {Feuille blanche (repérage couleur)};

% Étiquettes
\node[poppurple, font=\small\bfseries, align=center] at (1.2, 4.5) {Burette :\\ \ce{KMnO4} \SI{0,020}{M}};
\node[popgreen, font=\small\bfseries, align=center] at (6.8, 1.5) {Bécher :\\ \SI{10,0}{mL} \ce{Fe^2+} + \ce{H2SO4} \SI{1}{M}};
\node[popdark, font=\small, align=center] at (4.0, 7.8) {Support à burette};
\end{tikzpicture}
\legende{Dispositif de titration manganimétrique. La burette contient la solution violette de \ce{MnO4^-}. Le bécher (sur fond blanc) reçoit la solution incolore de \ce{Fe^2+} acidifiée. L'équivalence est signalée par l'apparition d'une teinte rose pâle persistante (excès de \ce{MnO4^-}).}
\end{popfigure}

\courssub{Protocole}
\begin{enumerate}
    \item \textbf{Préparation de la solution titrée (\ce{Fe^2+}) :} Péser \SI{3,00}{g} de \ce{FeSO4,7H2O} (masse molaire \SI{278,0}{g.mol^{-1}}) dans un bécher, dissoudre dans \SI{50}{mL} d'eau distillée, transvaser dans une fiole jaugée \SI{100}{mL}, compléter à l'eau distillée jusqu'au trait de jauge. Homogénéiser. Concentration théorique : \SI{0,108}{mol.L^{-1}}. \emph{Nota : Préparer juste avant la manipulation pour limiter l'oxydation par l'air.}
    \item \textbf{Mise en place de la burette (titrant \ce{KMnO4}) :} Rincer la burette 2 fois avec \SI{5}{mL} d'eau distillée (ouvrir le robinet), puis 2 fois avec \SI{5}{mL} de solution \ce{KMnO4} \SI{0,0200}{M} (récupérer les rinçages dans un bécher « déchets »). Remplir la burette au-dessus du trait \SI{0}{mL} à l'entonnoir, chasser la bulle d'air au robinet. Ajuster le ménisque sur le trait \SI{0,00}{mL} (lecture bas du ménisque, œil au niveau du liquide). Noter $V_i = \SI{0,00}{mL}$.
    \item \textbf{Prélèvement de l'analyte :} Rincer la pipette jaugée \SI{10,0}{mL} 2 fois eau distillée, puis 2 fois avec la solution \ce{Fe^2+} (peu de volume). Prélever \SI{10,0}{mL} de solution \ce{Fe^2+}, laisser s'écouler dans le bécher \SI{100}{mL} (ou erlenmeyer). Toucher la paroi avec l'embout 15 s, ne pas souffler dans la pipette.
    \item \textbf{Acidification :} Ajouter \SI{20}{mL} d'eau distillée puis \SI{10,0}{mL} d'acide sulfurique \ce{H2SO4} \SI{1,0}{M} (mesuré au cylindre gradué ou à la pipette). L'acide fournit les \ce{H+} nécessaires (\ce{8 H+} par \ce{MnO4-}) et empêche l'hydrolyse/précipitation de \ce{MnO2} brun. \textbf{Ne pas utiliser \ce{HCl} (oxydé en \ce{Cl2}) ni \ce{HNO3} (oxydant).}
    \item \textbf{Titrage :} Placer le bécher sur une feuille blanche sous la burette. Agiter doucement (baguette ou barre magnétique). Verser le \ce{KMnO4} \SI{1}{mL} par \SI{1}{mL} d'abord : la couleur violette disparaît instantanément (réaction rapide). Lorsque la décoloration ralentit, passer en \textbf{goutte à goutte} (1 goutte $\approx$ \SI{0,05}{mL}).
    \item \textbf{Repérage de l'équivalence :} Arrêter dès l'apparition d'une \textbf{teinte rose pâle persistante} (\textgreater 30 s sous agitation) signe d'un excès infime de \ce{MnO4-}. Noter le volume final $V_f$. Le volume équivalent $V_E = V_f - V_i$.
    \item \textbf{Répétition :} Refaire l'opération (étapes 3 à 6) deux fois minimum. Les volumes équivalents $V_{E1}, V_{E2}, V_{E3}$ doivent être \textbf{concordants} (écart \textless \SI{0,10}{mL} pour une burette classe B). Moyenne $V_E = \frac{\sum V_{Ei}}{n}$.
\end{enumerate}

\courssub{Observations et résultats attendus}
\begin{center}
\begin{tabularx}{\linewidth}{|c|c|c|c|c|}\hline
\rowcolor{popblueL}
\textbf{Essai} & \textbf{$V_i$ (mL)} & \textbf{$V_f$ (mL)} & \textbf{$V_E$ (mL)} & \textbf{Commentaire} \\\hline
1 & 0,00 & 10,05 & 10,05 & Rose pâle stable après 30 s \\\hline
2 & 0,00 & 10,02 & 10,02 & Concordant avec essai 1 \\\hline
3 & 0,00 & 10,04 & 10,04 & Concordant \\\hline
\multicolumn{2}{|c|}{\textbf{Moyenne $V_E$}} & \multicolumn{2}{c|}{\textbf{\SI{10,04}{mL}}} & \textbf{Écart max \SI{0,03}{mL} $<$ \SI{0,10}{mL} : résultats exploitables} \\\hline
\end{tabularx}
\end{center}
\begin{aretenir}[Observations clés]
\begin{itemize}
    \item La solution \ce{Fe^2+} est incolore (ou vert très pâle). L'ajout d'\ce{H2SO4} ne change pas la couleur.
    \item Tant que $\ce{Fe^2+}$ est en excès, le \ce{MnO4-} violet est réduit en \ce{Mn^2+} incolore : la solution reste incolore.
    \item Au point d'équivalence, tout le \ce{Fe^2+} est consommé. La goutte suivante de \ce{MnO4-} n'est plus réduite : la couleur violette/rose persiste.
    \item Pas d'indicateur ajouté : le \ce{MnO4-} est \textbf{auto-indicateur} (couleur intense, $\varepsilon \approx \SI{2300}{L.mol^{-1}.cm^{-1}}$ à \SI{525}{nm}).
\end{itemize}
\end{aretenir}

\courssub{Exploitation}
\begin{exoresolu}[Question 1 : Demi-équations électroniques]
Écrire les deux demi-équations électroniques des couples \ce{MnO4^-}/\ce{Mn^2+} et \ce{Fe^3+}/\ce{Fe^2+} en milieu acide.
\tcblower
\textbf{Solution :}
\begin{itemize}
    \item Couple \ce{MnO4^-}/\ce{Mn^2+} : \ce{Mn} passe de $+7$ à $+2$, gain de 5 électrons. Bilan O par \ce{H2O}, H par \ce{H+} :
    \[ \ce{MnO4^- + 8 H+ + 5 e- -> Mn^2+ + 4 H2O} \quad (E^\circ = \SI{+1,51}{V}) \]
    \item Couple \ce{Fe^3+}/\ce{Fe^2+} : \ce{Fe} passe de $+2$ à $+3$, perte de 1 électron :
    \[ \ce{Fe^2+ -> Fe^3+ + e-} \quad (E^\circ = \SI{+0,77}{V}) \]
\end{itemize}
\end{exoresolu}

\begin{exoresolu}[Question 2 : Équation-bilan]
Déduire l'équation-bilan de la réaction de dosage.
\tcblower
On multiplie la 2ème équation par 5 pour égaler les électrons, puis on additionne :
\[ \ce{MnO4^- + 8 H+ + 5 Fe^2+ -> Mn^2+ + 5 Fe^3+ + 4 H2O} \]
Vérification : Charges : $(-1) + 8(+1) + 5(+2) = +17$ à gauche ; $(+2) + 5(+3) = +17$ à droite. Atomes : OK.
\end{exoresolu}

\begin{exoresolu}[Question 3 : Choix de l'acide]
Justifier l'utilisation de l'acide sulfurique et non de l'acide chlorhydrique ou nitrique.
\tcblower
\begin{itemize}
    \item \ce{HCl} : \ce{Cl-} est réducteur ($E^\circ_{\ce{Cl2}/\ce{Cl-}} = \SI{+1,36}{V} < \SI{+1,51}{V}$). \ce{MnO4-} oxyderait \ce{Cl-} en \ce{Cl2} (gaz toxique, vert), faussant le dosage.
    \item \ce{HNO3} : \ce{NO3-} est oxydant ($E^\circ_{\ce{NO3-}/\ce{NO}} = \SI{+0,96}{V}$). Il pourrait oxyder \ce{Fe^2+} en \ce{Fe^3+} avant le titrage, ou perturber la cinétique.
    \item \ce{H2SO4} : \ce{SO4^2-} n'est ni oxydant ni réducteur dans ces conditions ($E^\circ_{\ce{SO4^2-}/\ce{SO2}} = \SI{+0,17}{V}$). Elle fournit les \ce{H+} nécessaires sans interférence redox. De plus, \ce{Fe^2+} et \ce{Fe^3+} forment des complexes stables avec \ce{SO4^2-} limitant l'hydrolyse.
\end{itemize}
\end{exoresolu}

\begin{exoresolu}[Question 4 : Auto-indication]
Pourquoi n'a-t-on pas besoin d'ajouter d'indicateur coloré (phénolphtaléine, etc.) ?
\tcblower
Le réactif titrant \ce{KMnO4} est intensément coloré (violet, $\lambda_{\max} \approx \SI{525}{nm}$). Le produit de réduction \ce{Mn^2+} est incolore en solution diluée. Tant que $\ce{Fe^2+}$ est présent, le \ce{MnO4-} ajouté est consommé $\rightarrow$ solution incolore. À l'équivalence, le \ce{Fe^2+} est épuisé ; le \ce{MnO4-} suivant reste en solution $\rightarrow$ apparition de la couleur rose/violette. C'est un \textbf{dosage auto-indiqué}. La sensibilité visuelle est excellente (seuil $\approx \SI{10^{-6}}{M}$ en \ce{MnO4-}).
\end{exoresolu}

\begin{exoresolu}[Question 5 : Calcul de la concentration]
Calculer la concentration molaire $C_{\ce{Fe^2+}}$ de la solution titrée.
\tcblower
Données : $C_{\ce{MnO4-}} = \SI{0,0200}{mol.L^{-1}}$, $V_E = \SI{10,04}{mL} = \SI{10,04 \times 10^{-3}}{L}$, $V_{\ce{Fe^2+}} = \SI{10,00}{mL} = \SI{10,00 \times 10^{-3}}{L}$.
\underline{Méthode 1 : Tableau d'avancement (recommandée au Probatoire)}
\[
\begin{array}{c|cccccc}
\text{État} & \ce{MnO4-} & + & 5\ce{Fe^2+} & + & 8\ce{H+} & \rightarrow \ce{Mn^2+} + 5\ce{Fe^3+} + 4\ce{H2O} \\
\hline
\text{Initial} & n_{\ce{MnO4-}}^0 & & n_{\ce{Fe^2+}}^0 & & \text{(excès)} & \\
\text{Final (Éq)} & 0 & & 0 & & & \\
\hline
\text{Avancement } x_E & x_E & & 5x_E & & & \\
\end{array}
\]
À l'équivalence : $n_{\ce{MnO4-}}^0 = x_E$ et $n_{\ce{Fe^2+}}^0 = 5x_E = 5 n_{\ce{MnO4-}}^0$.
\[ n_{\ce{MnO4-}}^0 = C_{\ce{MnO4-}} \times V_E = \SI{0,0200}{mol.L^{-1}} \times \SI{10,04 \times 10^{-3}}{L} = \SI{2,008 \times 10^{-4}}{mol} \]
\[ n_{\ce{Fe^2+}}^0 = 5 \times \SI{2,008 \times 10^{-4}}{mol} = \SI{1,004 \times 10^{-3}}{mol} \]
\[ C_{\ce{Fe^2+}} = \frac{n_{\ce{Fe^2+}}^0}{V_{\ce{Fe^2+}}} = \frac{\SI{1,004 \times 10^{-3}}{mol}}{\SI{10,00 \times 10^{-3}}{L}} = \SI{0,1004}{mol.L^{-1}} \]
\underline{Méthode 2 : Relation de dosage directe}
\[ \frac{C_{\ce{Fe^2+}} \times V_{\ce{Fe^2+}}}{5} = C_{\ce{MnO4-}} \times V_E \implies C_{\ce{Fe^2+}} = 5 \times \frac{C_{\ce{MnO4-}} \times V_E}{V_{\ce{Fe^2+}}} \]
\[ C_{\ce{Fe^2+}} = 5 \times \frac{\SI{0,0200}{mol.L^{-1}} \times \SI{10,04}{mL}}{\SI{10,00}{mL}} = \SI{0,1004}{mol.L^{-1}} \]
\textbf{Résultat :} $C_{\ce{Fe^2+}} = \SI{0,100}{mol.L^{-1}}$ (arrondi à 3 chiffres significatifs, cohérent avec la préparation théorique \SI{0,108}{M} compte tenu de l'oxydation partielle à l'air).
\end{exoresolu}

\begin{exoresolu}[Question 6 : Incertitude et résultat final]
Estimer l'incertitude-type $u(C_{\ce{Fe^2+}})$ et donner le résultat final sous la forme $C \pm u(C)$.
\tcblower
Incertitudes principales (burette classe B \SI{25}{mL} : $u(V_E) \approx \SI{0,03}{mL}$ ; pipette classe B \SI{10}{mL} : $u(V_{\ce{Fe^2+}}) \approx \SI{0,02}{mL}$ ; solution titrant certifiée : $u_{rel}(C_{\ce{MnO4-}}) \approx 0,5\%$).
\[ \frac{u(C_{\ce{Fe^2+}})}{C_{\ce{Fe^2+}}} \approx \sqrt{ \left(\frac{u(C_{\ce{MnO4-}})}{C_{\ce{MnO4-}}}\right)^2 + \left(\frac{u(V_E)}{V_E}\right)^2 + \left(\frac{u(V_{\ce{Fe^2+}})}{V_{\ce{Fe^2+}}}\right)^2 } \]
\[ \approx \sqrt{ (0,005)^2 + \left(\frac{0,03}{10,04}\right)^2 + \left(\frac{0,02}{10,00}\right)^2 } \approx \sqrt{ 2,5\times10^{-5} + 8,9\times10^{-6} + 4\times10^{-6} } \approx \sqrt{3,8\times10^{-5}} \approx 0,0062 \]
$u(C_{\ce{Fe^2+}}) \approx 0,0062 \times \SI{0,1004}{M} \approx \SI{0,0006}{M}$.
\textbf{Résultat final :} $\boxed{C_{\ce{Fe^2+}} = (\SI{0,100 \pm 0,001}{mol.L^{-1}})}$ (arrondi de l'incertitude à 1 chiffre significatif).
\end{exoresolu}

\begin{exoresolu}[Question 7 : Variante iodométrique]
Dosage du diiode par le thiosulfate (iodométrie).
\tcblower
Si pas de \ce{KMnO4} : On dose une solution de diiode \ce{I2} (teinture d'iode diluée \SI{0,01}{M}) par une solution étalon de thiosulfate \ce{Na2S2O3} \SI{0,050}{M}.
\begin{itemize}
    \item Demi-équations : $\ce{I2 + 2e- -> 2I-}$ ($E^\circ = \SI{+0,54}{V}$) ; $\ce{S4O6^2- + 2e- -> 2S2O3^2-}$ ($E^\circ = \SI{+0,08}{V}$).
    \item Équation-bilan : $\ce{I2 + 2 S2O3^2- -> 2 I- + S4O6^2-}$.
    \item Indicateur : \textbf{Amidon} (solution de maïzena bouillie \SI{1}{g} dans \SI{100}{mL} eau, filtrée). L'amidon forme un complexe bleu intense avec \ce{I2} (complexe charge-transfert).
    \item Protocole : Verser \ce{S2O3^2-} dans \ce{I2} acidifié (ou neutre) jusqu'à décoloration brutale du bleu (ajout amidon vers \SI{1}{mL} de l'équivalence pour éviter l'adsorption d'iode sur l'amidon).
    \item Calcul : $C_{\ce{I2}} = \frac{C_{\ce{S2O3^2-}} \times V_E}{2 \times V_{\ce{I2}}}$.
\end{itemize}
\end{exoresolu}

\courssub{Conclusion}
Ce TP a permis de réaliser un \textbf{dosage manganimétrique} : la détermination de la concentration d'une solution d'ions fer(II) par titration avec une solution étalon de permanganate de potassium en milieu acide sulfurique.
\begin{itemize}
    \item La réaction est \textbf{totale, rapide et unique} ($E^\circ_{\ce{MnO4-}/\ce{Mn^2+}} - E^\circ_{\ce{Fe^3+}/\ce{Fe^2+}} = \SI{0,74}{V} > \SI{0,3}{V}$), ce qui assure un saut d'équivalence net.
    \item Le \textbf{point d'équivalence} est détecté \textbf{visuellement sans indicateur} grâce à la couleur propre du \ce{MnO4-} (auto-indication) : passage de « incolore » à « rose pâle persistant ».
    \item L'exploitation rigoureuse via le \textbf{tableau d'avancement} (ou la relation de dosage $n_{\ce{Fe^2+}} = 5 n_{\ce{MnO4-}}$) a conduit à une concentration expérimentale de \textbf{\SI{0,100}{mol.L^{-1}}}, en bon accord avec la valeur théorique de préparation (\SI{0,108}{M}), la légère différence s'expliquant par l'oxydation atmosphérique du \ce{Fe^2+} en \ce{Fe^3+} lors de la préparation.
    \item La concordance des trois essais ($\Delta V_E < \SI{0,1}{mL}$) valide la reproductibilité de la manipulation.
    \item La \textbf{variante iodométrique} (\ce{I2}/\ce{S2O3^2-}) avec indicateur amidon est une alternative classique (iodométrie) utile en l'absence de permanganate, reposant sur le même principe redox mais avec un indicateur exogène (complexe \ce{I2}-amidon bleu).
\end{itemize}
\emph{Compétence validée :} « Réaliser un dosage d'oxydoréduction, repérer l'équivalence, exploiter les résultats pour déterminer une concentration. » — Prêt pour l'épreuve pratique du Probatoire !

\newpage

\coursec{Méthodes et exercices résolus}

\courssub{Fiches méthodes et exercices résolus}

\begin{methode}[Écrire une équation-bilan d'oxydoréduction avec les couples au programme]
Pour équilibrer une réaction d'oxydoréduction en solution aqueuse acide, procède toujours dans l'ordre suivant :
\begin{enumerate}
\item \textbf{Identifier les couples} mis en jeu et écrire chaque demi-équation électronique dans le sens où elle se produit (oxydation pour le réducteur, réduction pour l'oxydant).
\item \textbf{Équilibrer les demi-équations} : d'abord les atomes autres que O et H, puis l'oxygène avec des molécules \ce{H2O}, puis l'hydrogène avec des ions \ce{H+}, enfin la charge électrique avec des électrons \ce{e-}.
\item \textbf{Multiplier} chaque demi-équation par le plus petit coefficient entier qui égalise le nombre d'électrons échangés.
\item \textbf{Additionner} membre à membre : les électrons doivent disparaître totalement du bilan.
\item \textbf{Vérifier} la conservation des éléments et des charges dans l'équation finale.
\end{enumerate}
Réflexe : les couples oxygénés au programme s'écrivent en milieu acide :
\begin{align*}
\ce{MnO4- + 8H+ + 5e- &-> Mn^2+ + 4H2O}\\
\ce{Cr2O7^2- + 14H+ + 6e- &-> 2Cr^3+ + 7H2O}\\
\ce{NO3- + 4H+ + 3e- &-> NO + 2H2O}\\
\ce{SO4^2- + 4H+ + 2e- &-> SO2 + 2H2O}
\end{align*}
et les couples simples : \ce{Fe^3+ + e- -> Fe^2+}, \ce{Cl2 + 2e- -> 2Cl-}, \ce{I2 + 2e- -> 2I-}, \ce{S4O6^2- + 2e- -> 2S2O3^2-}.
\end{methode}

\begin{exoresolu}[Action du dichromate sur les ions iodure]
En milieu acide, une solution orangée de dichromate de potassium oxyde les ions iodure \ce{I-} en diiode \ce{I2}. Écrire les demi-équations électroniques et l'équation-bilan de la réaction.
\tcblower
\textbf{Étape 1 : identification des couples.} Les couples mis en jeu sont \ce{Cr2O7^2-}/\ce{Cr^3+} (oxydant \ce{Cr2O7^2-}) et \ce{I2}/\ce{I-} (réducteur \ce{I-}).

\textbf{Étape 2 : demi-équations équilibrées.}
\[
\ce{Cr2O7^2- + 14H+ + 6e- -> 2Cr^3+ + 7H2O} \qquad \ce{2I- -> I2 + 2e-}
\]

\textbf{Étape 3 : égalisation des électrons.} La première demi-équation consomme $6$ électrons, la seconde en libère $2$ : on multiplie la seconde par $3$.

\textbf{Étape 4 : addition.}
\[
\ce{Cr2O7^2- + 14H+ + 6I- -> 2Cr^3+ + 3I2 + 7H2O}
\]

\textbf{Étape 5 : vérification.} Charges : à gauche $-2+14-6=+6$, à droite $2\times(+3)=+6$. Éléments : $2\,\mathrm{Cr}$, $7\,\mathrm{O}$, $14\,\mathrm{H}$, $6\,\mathrm{I}$ de chaque côté. L'équation est correcte. \textbf{Conclusion :} \ce{Cr2O7^2- + 6I- + 14H+ -> 2Cr^3+ + 3I2 + 7H2O}.
\end{exoresolu}

\begin{methode}[Prévoir une réaction à l'aide de la classification électrochimique]
\begin{enumerate}
\item \textbf{Repérer les deux couples} présents dans le mélange et les placer sur l'axe des potentiels normaux $E^{\circ}$.
\item \textbf{Appliquer la règle du gamma} : l'oxydant du couple de potentiel le plus élevé réagit spontanément avec le réducteur du couple de potentiel le plus faible.
\item \textbf{Écrire les demi-équations} dans les bons sens, puis l'équation-bilan.
\item \textbf{Prévoir les observations} : changement de couleur, dégagement gazeux, dépôt solide.
\end{enumerate}
Rappel des potentiels usuels des couples du programme (ordre croissant) :
\[
\ce{S4O6^2-}/\ce{S2O3^2-}\ (\SI{+0,08}{V}) < \ce{SO4^2-}/\ce{SO2}\ (\SI{+0,17}{V}) < \ce{I2}/\ce{I-}\ (\SI{+0,54}{V}) < \ce{Fe^3+}/\ce{Fe^2+}\ (\SI{+0,77}{V}) < \ce{NO3-}/\ce{NO}\ (\SI{+0,96}{V}) < \ce{Cr2O7^2-}/\ce{Cr^3+}\ (\SI{+1,33}{V}) < \ce{Cl2}/\ce{Cl-}\ (\SI{+1,36}{V}) < \ce{MnO4-}/\ce{Mn^2+}\ (\SI{+1,51}{V})
\]
\emph{Repère : le couple \ce{H+}/\ce{H2} (\SI{0,00}{V}) sert de référence.}
\end{methode}

\begin{exoresolu}[L'ion permanganate oxyde-t-il l'ion fer(II) ?]
On mélange en milieu acide une solution de permanganate de potassium et une solution de sulfate de fer(II). Données : $E^{\circ}(\ce{MnO4-}/\ce{Mn^2+})=\SI{+1,51}{V}$ et $E^{\circ}(\ce{Fe^3+}/\ce{Fe^2+})=\SI{+0,77}{V}$. La réaction est-elle spontanée ? Écrire son équation-bilan et prévoir l'observation.
\tcblower
\textbf{Étape 1 : positionnement des couples.} $E^{\circ}(\ce{MnO4-}/\ce{Mn^2+}) > E^{\circ}(\ce{Fe^3+}/\ce{Fe^2+})$ : le permanganate est l'oxydant le plus fort, \ce{Fe^2+} est le réducteur le plus fort.

\textbf{Étape 2 : règle du gamma.} La réaction spontanée est celle entre \ce{MnO4-} (oxydant du couple le plus haut) et \ce{Fe^2+} (réducteur du couple le plus bas) : elle est donc \textbf{spontanée}.

\textbf{Étape 3 : équation-bilan.}
\[
\ce{MnO4- + 8H+ + 5e- -> Mn^2+ + 4H2O} \qquad \ce{Fe^2+ -> Fe^3+ + e-} \ (\times 5)
\]
\[
\boxed{\ce{MnO4- + 5Fe^2+ + 8H+ -> Mn^2+ + 5Fe^3+ + 4H2O}}
\]

\textbf{Étape 4 : observation.} La solution violette de permanganate se \textbf{décolore} progressivement : les ions \ce{Mn^2+} formés sont quasi incolores et les ions \ce{Fe^3+} donnent une teinte jaunâtre pâle. Cette décoloration est la base du \cle{dosage manganimétrique}.
\end{exoresolu}

\begin{methode}[Exploiter l'équivalence d'un dosage d'oxydoréduction]
\begin{enumerate}
\item \textbf{Écrire l'équation-bilan} du dosage et relever les coefficients stœchiométriques de l'oxydant et du réducteur.
\item \textbf{Identifier le rôle des solutions} : solution \cle{titrée} (réactif à doser, volume $V$ connu avec précision, prélevé à la \cle{pipette jaugée}) et solution \cle{titrante} (concentration $C$ connue, placée dans la \cle{burette}).
\item \textbf{Repérer l'équivalence} : changement de couleur persistant (\cle{auto-indication} avec \ce{MnO4-}, empois d'amidon avec \ce{I2}).
\item \textbf{Écrire la relation d'équivalence} : si l'équation est $a\,\mathrm{Ox} + b\,\mathrm{Red} \rightarrow$ produits, alors
\[
\frac{n_{\mathrm{Ox}}}{a}=\frac{n_{\mathrm{Red}}}{b} \quad \text{soit} \quad \frac{C_{\mathrm{Ox}}V_{E}}{a}=\frac{C_{\mathrm{Red}}V_{\mathrm{Red}}}{b}.
\]
\item \textbf{Calculer} la concentration inconnue en respectant les unités (volumes dans la même unité) et les \cle{chiffres significatifs}.
\end{enumerate}
\end{methode}

\begin{exoresolu}[Dosage manganimétrique des ions fer(II)]
On dose $V_{\mathrm{Red}}=\SI{20,0}{mL}$ d'une solution de sulfate de fer(II) acidifiée à l'acide sulfurique par une solution de permanganate de potassium de concentration $C_{\mathrm{Ox}}=\SI{2,00e-2}{mol.L^{-1}}$. L'équivalence est atteinte pour $V_{E}=\SI{16,4}{mL}$ : la teinte rose persiste alors.
\begin{enumerate}
\item Écrire l'équation-bilan du dosage.
\item Définir l'équivalence et expliquer comment elle est repérée.
\item Calculer la concentration $C_{\mathrm{Red}}$ de la solution en ions \ce{Fe^2+}, puis la concentration massique en fer(II) ($M(\mathrm{Fe})=\SI{55,8}{g.mol^{-1}}$).
\end{enumerate}
\tcblower
\textbf{1. Équation-bilan.} Couples : \ce{MnO4-}/\ce{Mn^2+} et \ce{Fe^3+}/\ce{Fe^2+}.
\[
\ce{MnO4- + 5Fe^2+ + 8H+ -> Mn^2+ + 5Fe^3+ + 4H2O}
\]

\textbf{2. Équivalence.} C'est l'instant où les réactifs ont été mélangés dans les proportions stœchiométriques : tout le \ce{Fe^2+} est consommé. Avant l'équivalence, chaque goutte de \ce{MnO4-} est décolorée par le \ce{Fe^2+} en excès ; à l'équivalence, la première goutte en excès de permanganate donne à la solution une \textbf{teinte rose persistante} : \ce{MnO4-} joue le rôle d'indicateur de fin de réaction (\cle{auto-indication}).

\textbf{3. Calcul de la concentration.} Relation d'équivalence avec $a=1$ et $b=5$ :
\[
\frac{C_{\mathrm{Ox}}V_{E}}{1}=\frac{C_{\mathrm{Red}}V_{\mathrm{Red}}}{5}
\quad\Longrightarrow\quad
C_{\mathrm{Red}}=\frac{5\,C_{\mathrm{Ox}}V_{E}}{V_{\mathrm{Red}}}
=\frac{5\times\SI{2,00e-2}{}\times 16,4}{20,0}
\]
\[
C_{\mathrm{Red}}=\SI{8,20e-2}{mol.L^{-1}}.
\]
Concentration massique :
\[
C_{m}=C_{\mathrm{Red}}\times M(\mathrm{Fe})=\SI{8,20e-2}{}\times 55,8=\SI{4,58}{g.L^{-1}}.
\]
\textbf{Conclusion :} la solution contient $\SI{8,20e-2}{mol.L^{-1}}$ d'ions \ce{Fe^2+}, soit $\SI{4,58}{g.L^{-1}}$ de fer(II).
\end{exoresolu}

\begin{methode}[Réaliser et exploiter un dosage iodométrique (\ce{I2} par \ce{S2O3^2-})]
\begin{enumerate}
\item \textbf{Équation du dosage} : \ce{I2 + 2S2O3^2- -> 2I- + S4O6^2-}. Relation d'équivalence : $n(\ce{I2})=\dfrac{n(\ce{S2O3^2-})}{2}$.
\item \textbf{Montage} : solution de diiode dans l'erlenmeyer (volume prélevé à la pipette jaugée), solution de thiosulfate de sodium dans la burette.
\item \textbf{Indicateur} : l'empois d'amidon, ajouté quand la solution n'est plus que jaune pâle ; l'équivalence est marquée par la \textbf{décoloration} du bleu foncé.
\item \textbf{Calcul} : $C(\ce{I2})=\dfrac{C(\ce{S2O3^2-})\,V_{E}}{2\,V(\ce{I2})}$.
\item \textbf{Dosage en retour possible} : on peut d'abord faire réagir un oxydant (ex. \ce{MnO4-}) sur un excès de \ce{I-}, puis doser le \ce{I2} formé.
\end{enumerate}
\end{methode}

\begin{exoresolu}[Dosage d'une eau iodée]
Une solution de diiode (eau iodée utilisée comme antiseptique) est dosée par une solution de thiosulfate de sodium de concentration $C_{T}=\SI{5,00e-2}{mol.L^{-1}}$. On prélève $V_{I}=\SI{10,0}{mL}$ d'eau iodée ; la décoloration de l'empois d'amidon survient pour $V_{E}=\SI{12,8}{mL}$.
\begin{enumerate}
\item Écrire l'équation-bilan et schématiser le dispositif.
\item Calculer la concentration molaire puis massique du diiode ($M(\ce{I2})=\SI{253,8}{g.mol^{-1}}$).
\end{enumerate}
\tcblower
\textbf{1. Équation-bilan.} Couples \ce{I2}/\ce{I-} et \ce{S4O6^2-}/\ce{S2O3^2-} :
\[
\ce{I2 + 2e- -> 2I-} \qquad \ce{2S2O3^2- -> S4O6^2- + 2e-}
\]
\[
\boxed{\ce{I2 + 2S2O3^2- -> 2I- + S4O6^2-}}
\]
Dispositif : burette contenant la solution de thiosulfate au-dessus d'un erlenmeyer contenant le prélèvement d'eau iodée et quelques gouttes d'empois d'amidon, le tout sur agitation magnétique.

\textbf{2. Calculs.} À l'équivalence : $n(\ce{I2})=\dfrac{n(\ce{S2O3^2-})}{2}$, donc
\[
C(\ce{I2})=\frac{C_{T}\,V_{E}}{2\,V_{I}}
=\frac{\SI{5,00e-2}{}\times 12,8}{2\times 10,0}
=\SI{3,20e-2}{mol.L^{-1}}.
\]
Concentration massique :
\[
C_{m}=\SI{3,20e-2}{}\times 253,8=\SI{8,12}{g.L^{-1}}.
\]
\textbf{Conclusion :} l'eau iodée contient $\SI{3,20e-2}{mol.L^{-1}}$ de diiode, soit $\SI{8,12}{g.L^{-1}}$.
\end{exoresolu}

\begin{methode}[Exploiter un tableau d'avancement à l'équivalence]
Pour les mélanges non stœchiométriques ou les questions « que reste-t-il après réaction ? » :
\begin{enumerate}
\item \textbf{Calculer les quantités initiales} $n=C\times V$ de chaque réactif (attention : dilution éventuelle dans le mélange, mais pour l'avancement on utilise les quantités, pas les concentrations).
\item \textbf{Dresser le tableau d'avancement} avec l'équation-bilan.
\item \textbf{Déterminer l'avancement maximal} $x_{\max}$ : le plus petit rapport $n_{0}/\text{coefficient}$ désigne le \cle{réactif limitant}.
\item \textbf{Conclure} sur les quantités finales et, si besoin, sur les concentrations dans le volume total.
\end{enumerate}
\end{methode}

\begin{exoresolu}[Mélange permanganate -- fer(II) hors équivalence]
On verse $V_{1}=\SI{25,0}{mL}$ d'une solution de \ce{KMnO4} à $C_{1}=\SI{1,00e-2}{mol.L^{-1}}$ dans $V_{2}=\SI{25,0}{mL}$ d'une solution acidifiée de \ce{Fe^2+} à $C_{2}=\SI{8,00e-2}{mol.L^{-1}}$. Déterminer le réactif limitant, l'avancement maximal et la couleur finale du mélange.
\tcblower
\textbf{Quantités initiales :}
\[
n_{0}(\ce{MnO4-})=C_{1}V_{1}=\SI{1,00e-2}{}\times\SI{25,0e-3}{}=\SI{2,50e-4}{mol}
\]
\[
n_{0}(\ce{Fe^2+})=C_{2}V_{2}=\SI{8,00e-2}{}\times\SI{25,0e-3}{}=\SI{2,00e-3}{mol}
\]

\textbf{Tableau d'avancement} pour \ce{MnO4- + 5Fe^2+ + 8H+ -> Mn^2+ + 5Fe^3+ + 4H2O} (ions \ce{H+} en excès) :

\begin{center}
\begin{tabularx}{\linewidth}{|l|X|X|X|X|}
\hline
\rowcolor{popblueL} État & \ce{MnO4-} & \ce{Fe^2+} & \ce{Mn^2+} & \ce{Fe^3+} \\
\hline
Initial (mol) & $2,50\times10^{-4}$ & $2,00\times10^{-3}$ & $0$ & $0$ \\
\hline
Final (mol) & $2,50\times10^{-4}-x_{\max}$ & $2,00\times10^{-3}-5x_{\max}$ & $x_{\max}$ & $5x_{\max}$ \\
\hline
\end{tabularx}
\end{center}

\textbf{Réactif limitant :} $\dfrac{2,50\times10^{-4}}{1}=\SI{2,50e-4}{mol}$ et $\dfrac{2,00\times10^{-3}}{5}=\SI{4,00e-4}{mol}$. Le plus petit rapport correspond à \ce{MnO4-} : c'est le réactif limitant et $x_{\max}=\SI{2,50e-4}{mol}$.

\textbf{Bilan final :} $n(\ce{MnO4-})=0$ ; $n(\ce{Fe^2+})=2,00\times10^{-3}-5\times2,50\times10^{-4}=\SI{7,50e-4}{mol}$ ; $n(\ce{Mn^2+})=\SI{2,50e-4}{mol}$ ; $n(\ce{Fe^3+})=\SI{1,25e-3}{mol}$.

\textbf{Conclusion :} le permanganate est entièrement consommé ; la solution finale est \textbf{incolore à jaune pâle} (présence de \ce{Fe^3+} et de \ce{Fe^2+} restant), la teinte violette ayant disparu.
\end{exoresolu}

\begin{methode}[Schématiser et décrire le dispositif d'un dosage d'oxydoréduction]
\begin{enumerate}
\item \textbf{Burette graduée} contenant la solution titrante (concentration connue), fixée sur un support.
\item \textbf{Erlenmeyer} contenant le prélèvement exact de solution titrée (pipette jaugée), l'acidification éventuelle (acide sulfurique, \textbf{jamais \ce{HCl} avec \ce{MnO4-}} !) et l'indicateur si nécessaire.
\item \textbf{Agitation} constante (agitateur magnétique) ; repérage de l'équivalence goutte à goutte près de la fin.
\item \textbf{Légender} le schéma : nom et concentration de chaque solution, volumes.
\end{enumerate}
\end{methode}

\begin{exoresolu}[Schéma du dosage manganimétrique]
Schématiser et légender le dispositif du dosage d'une solution de \ce{Fe^2+} par une solution acidifiée de \ce{KMnO4}.
\tcblower
\begin{popfigure}
\begin{tikzpicture}[scale=0.9]
% Burette
\draw[thick,popdark] (0,4.2) -- (0,0.8);
\draw[thick,popdark] (0.5,4.2) -- (0.5,0.8);
\draw[thick,popdark] (0,0.8) -- (0.15,0.5) -- (0.15,0.15);
\draw[thick,popdark] (0.5,0.8) -- (0.35,0.5) -- (0.35,0.15);
\draw[thick,popdark] (0.15,0.15) -- (0.35,0.15);
\draw[fill=poppinkL,draw=none] (0.03,4.1) rectangle (0.47,1.6);
\draw[popdark] (0,4.2) -- (0.5,4.2);
\foreach \y in {1.8,2.2,2.6,3.0,3.4,3.8} \draw[popdark] (0.5,\y) -- (0.62,\y);
\node[right,popdark] at (0.7,2.9) {\small \ce{KMnO4} : $C$ connue};
% Erlenmeyer
\draw[thick,popdark] (-1.1,-2.6) -- (-0.35,-1.1) -- (-0.35,-0.6);
\draw[thick,popdark] (1.6,-2.6) -- (0.85,-1.1) -- (0.85,-0.6);
\draw[thick,popdark] (-1.1,-2.6) -- (1.6,-2.6);
\draw[thick,popdark] (-0.35,-0.6) -- (0.85,-0.6);
\draw[fill=popblueL,draw=none] (-0.95,-2.55) -- (1.45,-2.55) -- (0.95,-1.6) -- (-0.45,-1.6) -- cycle;
\node[right,popdark] at (1.8,-2.1) {\small \ce{Fe^2+} : $V=\SI{20,0}{mL}$ (pipette jaugée)};
\node[right,popdark] at (1.8,-2.6) {\small + acide sulfurique};
% Goutte
\draw[fill=poppink,draw=poppink] (0.25,-0.15) circle (0.05);
% Support
\draw[thick,popdark] (-1.6,4.2) -- (-1.6,-2.8);
\draw[thick,popdark] (-1.6,3.6) -- (0,3.6);
\draw[thick,popdark] (-2.2,-2.8) -- (2.2,-2.8);
\end{tikzpicture}
\legende{Dispositif du dosage manganimétrique : la solution titrante de permanganate (violette) est dans la burette ; la solution titrée d'ions fer(II), acidifiée à l'acide sulfurique, est dans l'erlenmeyer.}
\end{popfigure}
La solution violette de \ce{KMnO4} est versée goutte à goutte ; elle se décolore tant qu'il reste du \ce{Fe^2+}. L'équivalence est repérée à la \textbf{première teinte rose persistante}. On note alors le volume $V_{E}$ lu sur la burette et on applique la relation d'équivalence $C_{\mathrm{Red}}=\dfrac{5\,C_{\mathrm{Ox}}V_{E}}{V_{\mathrm{Red}}}$.
\end{exoresolu}

\begin{attention}[Les 5 erreurs qui coûtent des points au Probatoire]
\begin{enumerate}
\item \textbf{Électrons dans l'équation-bilan finale} : après multiplication et addition des demi-équations, aucun \ce{e-} ne doit subsister. Vérifie aussi la conservation des charges : un bilan non équilibré en charge est faux.
\item \textbf{Mauvais sens des demi-équations} : le réducteur s'oxyde (il \emph{perd} des électrons, écrits à droite), l'oxydant se réduit (il \emph{gagne} des électrons, écrits à gauche). Ne jamais écrire \ce{Fe^3+ + e- -> Fe^2+} si \ce{Fe^2+} est le réducteur du mélange.
\item \textbf{Oubli du coefficient stœchiométrique dans la relation d'équivalence} : pour \ce{MnO4- + 5Fe^2+ + ...}, c'est $n(\ce{Fe^2+})=5\,n(\ce{MnO4-})$, et non l'inverse. Pour l'iodométrie : $n(\ce{S2O3^2-})=2\,n(\ce{I2})$.
\item \textbf{Acidification avec l'acide chlorhydrique} : les ions \ce{Cl-} sont oxydés par \ce{MnO4-} (dégagement de \ce{Cl2}), ce qui fausse le dosage. On acidifie toujours à l'\textbf{acide sulfurique}.
\item \textbf{Confusion des rôles et des unités} : la solution titrante est dans la \textbf{burette}, la solution titrée est prélevée à la \textbf{pipette jaugée}. Convertis les volumes en litres pour les quantités de matière (ou garde les deux volumes dans la même unité dans la relation d'équivalence), et respecte les chiffres significatifs dans le résultat final.
\end{enumerate}
\end{attention}

\newpage

\coursec{Exercices}

\coursec{Leçon 08 : Généralisation et dosages d'oxydoréduction}

\courssub{1. Rappels et généralisation de la classification électrochimique}

\begin{savaistu}[Les couples au programme de Première]
Le programme officiel de Première C, D, E impose la connaissance des potentiels standards des couples suivants en solution aqueuse (valeurs à \SI{25}{\celsius}, pression \SI{1}{bar}, concentrations \SI{1}{mol.L^{-1}} pour les espèces dissoutes) :
\begin{center}
\begin{tabularx}{\linewidth}{|l|c|X|}\hline
\textbf{Couple redox} & $E^\circ$ (\si{V}) & \textbf{Rôle usuel} \\\hline
\ce{MnO4-}/\ce{Mn^2+} & $+1,51$ & Oxydant très fort (milieu acide) \\\hline
\ce{Cr2O7^2-}/\ce{Cr^3+} & $+1,33$ & Oxydant fort (milieu acide) \\\hline
\ce{Cl2}/\ce{Cl-} & $+1,36$ & Oxydant fort (gaz) \\\hline
\ce{NO3-}/\ce{NO} & $+0,96$ & Oxydant (milieu acide) \\\hline
\ce{Fe^3+}/\ce{Fe^2+} & $+0,77$ & Oxydant / Réducteur modéré \\\hline
\ce{I2}/\ce{I-} & $+0,54$ & Oxydant doux \\\hline
\ce{SO4^2-}/\ce{SO2} & $+0,17$ & Oxydant très doux (milieu acide) \\\hline
\ce{S4O6^2-}/\ce{S2O3^2-} & $+0,09$ & Réducteur (iodométrie) \\\hline
\end{tabularx}
\end{center}
\end{savaistu}

\begin{definition}[Potentiel standard d'un couple redox]
Le \cle{potentiel standard} $E^\circ$ d'un couple $\ce{Ox}/\ce{Red}$ est la tension (en volt) mesurée à \SI{25}{\celsius} pour une demi-pile où l'activité de l'oxydant et du réducteur vaut 1 (concentration \SI{1}{mol.L^{-1}} pour les espèces en solution, pression \SI{1}{bar} pour les gaz), reliée à une électrode de référence (électrode hydrogène standard, EHS, définie à \SI{0}{V}).
\end{definition}

\begin{propriete}[Échelle des potentiels standards et règle du gamma]
On représente les couples sur une \cle{échelle verticale} de potentiels croissants de bas en haut.
\begin{itemize}
\item Plus le potentiel $E^\circ$ est \textbf{élevé}, plus l'oxydant est \textbf{fort} (il a une forte tendance à capter des électrons).
\item Plus le potentiel $E^\circ$ est \textbf{faible}, plus le réducteur est \textbf{fort} (il a une forte tendance à donner des électrons).
\end{itemize}
\cle{Règle du gamma (prévisibilité)} : Une réaction d'oxydoréduction est spontanée dans le sens où l'\textbf{oxydant du couple de potentiel le plus élevé} réagit avec le \textbf{réducteur du couple de potentiel le plus faible}.
Graphiquement, la flèche $\gamma$ part de l'oxydant fort (en haut) vers le réducteur fort (en bas).
\end{propriete}

\begin{exemplebox}[Positionnement des couples et prévision]
On place les couples du programme sur l'échelle (extrait) :
\[
\begin{array}{c|c}
E^\circ (\si{V}) & \text{Couple} \\ \hline
+1,51 & \ce{MnO4-}/\ce{Mn^2+} \\
+1,36 & \ce{Cl2}/\ce{Cl-} \\
+1,33 & \ce{Cr2O7^2-}/\ce{Cr^3+} \\
+0,96 & \ce{NO3-}/\ce{NO} \\
+0,77 & \ce{Fe^3+}/\ce{Fe^2+} \\
+0,54 & \ce{I2}/\ce{I-} \\
+0,17 & \ce{SO4^2-}/\ce{SO2} \\
+0,09 & \ce{S4O6^2-}/\ce{S2O3^2-} \\
\end{array}
\]
\emph{Exemple :} $\ce{MnO4-}$ ($E^\circ=1,51$) peut oxyder $\ce{Fe^2+}$ ($E^\circ=0,77$) car $1,51 > 0,77$. $\ce{I2}$ ($E^\circ=0,54$) ne peut pas oxyder $\ce{Fe^2+}$.
\end{exemplebox}

\courssub{2. Équations-bilan d'oxydoréduction avec les couples au programme}

\begin{methode}[Écriture de l'équation-bilan d'une réaction redox]
\begin{enumerate}
\item \textbf{Identifier} les deux couples mis en jeu (Ox1/Red1 et Ox2/Red2).
\item \textbf{Écrire} les deux demi-équations électroniques (réduction de l'oxydant, oxydation du réducteur) en équilibrant : atomes $\neq$ O, H ; atomes O (par $\ce{H2O}$) ; atomes H (par $\ce{H+}$ en milieu acide) ; charges (par $\ce{e-}$).
\item \textbf{Égaliser} le nombre d'électrons en multipliant les demi-équations par les coefficients adéquats (PPCM des nombres d'électrons).
\item \textbf{Additionner} les deux demi-équations et \textbf{simplifier} les espèces communes ($\ce{e-}$, $\ce{H+}$, $\ce{H2O}$).
\item \textbf{Vérifier} : bilan des atomes et bilan des charges.
\end{enumerate}
\end{methode}

\begin{aretenir}[Demi-équations de référence (milieu acide)]
\begin{align*}
\ce{MnO4- + 8H+ + 5e- &-> Mn^2+ + 4H2O} \quad (E^\circ = \SI{1,51}{V}) \\
\ce{Cr2O7^2- + 14H+ + 6e- &-> 2Cr^3+ + 7H2O} \quad (E^\circ = \SI{1,33}{V}) \\
\ce{NO3- + 4H+ + 3e- &-> NO + 2H2O} \quad (E^\circ = \SI{0,96}{V}) \\
\ce{Fe^3+ + e- &-> Fe^2+} \quad (E^\circ = \SI{0,77}{V}) \\
\ce{I2 + 2e- &-> 2I-} \quad (E^\circ = \SI{0,54}{V}) \\
\ce{SO4^2- + 4H+ + 2e- &-> SO2 + 2H2O} \quad (E^\circ = \SI{0,17}{V}) \\
\ce{S4O6^2- + 2e- &-> 2S2O3^2-} \quad (E^\circ = \SI{0,09}{V}) \\
\ce{Cl2 + 2e- &-> 2Cl-} \quad (E^\circ = \SI{1,36}{V})
\end{align*}
\end{aretenir}

\begin{exoresolu}[Équation-bilan : \ce{MnO4-} / \ce{Fe^2+}]
\emph{Demi-équations :}
\begin{align*}
\text{Réduction : } & \ce{MnO4- + 8H+ + 5e- -> Mn^2+ + 4H2O} \\
\text{Oxydation : } & \ce{Fe^2+ -> Fe^3+ + e-} \quad (\times 5)
\end{align*}
\emph{Somme :} $\ce{MnO4- + 5Fe^2+ + 8H+ -> Mn^2+ + 5Fe^3+ + 4H2O}$.
\end{exoresolu}

\begin{exoresolu}[Équation-bilan : \ce{Cr2O7^2-} / \ce{Fe^2+}]
\emph{Demi-équations :}
\begin{align*}
\text{Réduction : } & \ce{Cr2O7^2- + 14H+ + 6e- -> 2Cr^3+ + 7H2O} \\
\text{Oxydation : } & \ce{Fe^2+ -> Fe^3+ + e-} \quad (\times 6)
\end{align*}
\emph{Somme :} $\ce{Cr2O7^2- + 6Fe^2+ + 14H+ -> 2Cr^3+ + 6Fe^3+ + 7H2O}$.
\end{exoresolu}

\begin{exoresolu}[Équation-bilan : \ce{I2} / \ce{S2O3^2-} (Iodométrie)]
\emph{Demi-équations :}
\begin{align*}
\text{Réduction : } & \ce{I2 + 2e- -> 2I-} \\
\text{Oxydation : } & \ce{2S2O3^2- -> S4O6^2- + 2e-}
\end{align*}
\emph{Somme :} $\ce{I2 + 2S2O3^2- -> 2I- + S4O6^2-}$.
\end{exoresolu}

\courssub{3. Principe général d'un dosage d'oxydoréduction}

\begin{definition}[Dosage d'oxydoréduction]
Un \cle{dosage d'oxydoréduction} est une opération qui consiste à déterminer la concentration d'une espèce chimique (le \cle{dosé}) en la faisant réagir complètement avec une solution étalonnée d'une autre espèce (le \cle{doseur}) au cours d'une réaction d'oxydoréduction rapide, totale et de stœchiométrie connue.
\end{definition}

\begin{propriete}[Point d'équivalence]
Au \cle{point d'équivalence} (volume $V_E$), le doseur et le dosé ont été introduits dans les proportions \textbf{stœchiométriques} de l'équation-bilan de la réaction de dosage.
Si l'équation-bilan est : $a\,\ce{Ox} + b\,\ce{Red} \rightarrow \text{produits}$, alors à l'équivalence :
\[
\frac{n_{\text{doseur}}}{a} = \frac{n_{\text{dosé}}}{b} \quad \text{soit} \quad \frac{C_{\text{doseur}} V_E}{a} = \frac{C_{\text{dosé}} V_{\text{dosé}}}{b}
\]
\end{propriete}

\begin{aretenir}[Repérage de l'équivalence]
Trois modes principaux :
\begin{enumerate}
\item \textbf{Dosage colorimétrique à auto-indicateur :} Le doseur est intensément coloré (ex. $\ce{MnO4-}$ violet). L'équivalence est marquée par l'apparition persistante de la couleur du doseur en excès (teinte \textbf{violette pâle} pour le permanganate).
\item \textbf{Dosage avec indicateur redox ajouté :} L'indicateur change de couleur lorsque le potentiel de la solution franchit son potentiel de virage (ex. diphénylamine sulfonée pour $\ce{Cr2O7^2-}/\ce{Fe^2+}$, virage incolore $\to$ violet).
\item \textbf{Dosage avec indicateur spécifique (complexation) :} En iodométrie, l'\textbf{empois d'amidon} forme un complexe bleu intense avec le diiode $\ce{I2}$. L'équivalence correspond à la \textbf{disparition brutale de la coloration bleue} (passage au jaune pâle puis incolore).
\end{enumerate}
\end{aretenir}

\begin{experience}[Dispositif expérimental type (Dosage au bécher)]
\begin{popfigure}
\begin{tikzpicture}[scale=0.8, >=Stealth]
% Support
\draw[thick, popmuted] (-1,0) -- (7,0);
% Burette
\draw[thick] (0.5,0.5) -- (0.5,6) -- (1.5,6) -- (1.5,0.5);
\draw[thick] (1.5,1.5) -- (2,1.5) -- (2,1) -- (2.5,1) -- (2.5,0.5) -- (2,0.5) -- (2,1) -- (1.5,1); % robinet
\fill[popblue!30] (0.5,0.8) rectangle (1.5,5.5);
\node[popblue, align=center] at (1,5.8) {Burette\\(Doseur)};
\node[popblue] at (1,0.3) {\ce{MnO4-} ou \ce{S2O3^2-}};
% Graduations
\foreach \y in {1,2,3,4,5} \draw[thin] (0.3,\y) -- (1.7,\y) node[right, font=\tiny] {$\y$};
% Bécher
\draw[thick] (4,0.5) -- (4,4) -- (6,4) -- (6,0.5);
\draw[thick] (3.8,0.5) -- (6.2,0.5); % base
\fill[poporange!20] (4,0.5) rectangle (6,3);
\node[poporange, align=center] at (5,4.3) {Bécher\\(Dosé)};
\node[poporange] at (5,0) {\ce{Fe^2+ + H2SO4} ou \ce{I2}};
% Agitateur magnétique
\draw[thick, popdark] (5,4) -- (5,5);
\draw[thick, decorate, decoration={coil, aspect=0.3, segment length=2mm, amplitude=1.5mm}] (5,3.5) -- (5,2);
\node[popdark, right] at (5,4.5) {Agitateur};
% Plaque blanche
\fill[white, draw=popmuted, thick] (3.5,-0.5) rectangle (6.5,-0.1);
\node[font=\small, popmuted] at (5,-0.8) {Plaque blanche (repérage couleur)};
% Pipette (symbolique)
\draw[->, thick, popgreen] (3,3.5) -- (4,3.5) node[midway, above, font=\small] {Prélèvement $V_{\text{dosé}}$};
\end{tikzpicture}
\legende{Dispositif de dosage d'oxydoréduction : burette graduée (doseur), bécher sur agitateur magnétique (dosé + milieu acide), plaque blanche pour la détection visuelle.}
\end{popfigure}
\end{experience}

\begin{attention}[Pièges classiques à éviter]
\begin{itemize}
\item \textbf{Acidification :} Pour $\ce{MnO4-}$ et $\ce{Cr2O7^2-}$, utiliser \textbf{uniquement $\ce{H2SO4}$ dilué}. $\ce{HCl}$ est oxydé par $\ce{MnO4-}$ ($\ce{Cl2}$ gazeux toxique) et par $\ce{Cr2O7^2-}$. $\ce{HNO3}$ est un oxydant parasite.
\item \textbf{Vitesse :} Les dosages $\ce{Fe^2+}$ doivent être rapides car $\ce{Fe^2+}$ s'oxyde au contact de l'air ($\ce{O2}$).
\item \textbf{Indicateur amidon :} L'ajouter \textbf{uniquement près de l'équivalence} (solution jaune paille) pour éviter la formation d'un complexe amidon-iode trop stable qui fausserait $V_E$.
\item \textbf{Stœchiométrie :} Ne pas confondre le coefficient du couple (ex: 6 $\ce{Fe^2+}$ pour 1 $\ce{Cr2O7^2-}$) et le nombre d'électrons du couple (6 e- pour $\ce{Cr2O7^2-}$).
\end{itemize}
\end{attention}

\courssub{4. Dosage des ions \texorpdfstring{$\ce{Fe^2+}$}{Fe2+} par les ions \texorpdfstring{$\ce{MnO4-}$}{MnO4-} (Dosage manganimétrique)}

\begin{propriete}[Réaction de dosage manganimétrique]
En milieu acide sulfurique, l'ion permanganate $\ce{MnO4-}$ (violet) oxyde l'ion fer(II) $\ce{Fe^2+}$ (vert pâle) en ion fer(III) $\ce{Fe^3+}$ (jaune/brun) et se réduit en ion manganèse(II) $\ce{Mn^2+}$ (incolore).
\[
\ce{MnO4- + 5Fe^2+ + 8H+ -> Mn^2+ + 5Fe^3+ + 4H2O}
\]
\cle{Auto-indicateur :} Tant que $\ce{Fe^2+}$ est en excès, le $\ce{MnO4-}$ ajouté est instantanément décoloré. À l'équivalence, la première goutte de $\ce{MnO4-}$ en excès donne une coloration \textbf{violette pâle persistante} (environ \SI{30}{\second}).
\end{propriete}

\begin{methode}[Protocole opératoire du dosage manganimétrique]
\begin{enumerate}
\item \textbf{Rinçage :} Rincer la burette à l'eau distillée, puis avec \textbf{quelques mL de la solution doseuse ($\ce{KMnO4}$)}. Évacuer par le robinet.
\item \textbf{Remplissage :} Remplir la burette avec la solution de $\ce{KMnO4}$, chasser les bulles d'air au niveau du robinet. Noter le volume initial $V_i$ (bas du ménisque, lecture en haut pour liquide foncé).
\item \textbf{Préparation du dosé :} Pipeter le volume $V_1$ de solution de $\ce{Fe^2+}$ dans le bécher. Ajouter \textbf{$\sim \SI{20}{mL}$ d'acide sulfurique $\ce{H2SO4}$ \SI{1}{mol.L^{-1}}} (milieu acide + volume suffisant). Placer sur agitateur magnétique avec plaque blanche.
\item \textbf{Titrage :} Verser le $\ce{KMnO4}$ goutte à goutte sous agitation. Au début, la coloration violette disparaît instantanément. Ralentir près de l'équivalence (coloration persistante quelques secondes).
\item \textbf{Équivalence :} Arrêter dès l'apparition d'une \textbf{coloration violette pâle persistante} (\SI{30}{\second}). Noter le volume final $V_f$. $V_E = V_f - V_i$.
\item \textbf{Calcul :} Utiliser la relation à l'équivalence pour trouver $C_{\ce{Fe^2+}}$.
\end{enumerate}
\end{methode}

\courssub{5. Dosage du diiode par les ions thiosulfate (Iodométrie)}

\begin{propriete}[Réaction de l'iodométrie]
Le diiode $\ce{I2}$ (brun/jaune en solution aqueuse, violet en phase organique) est réduit en iodure $\ce{I-}$ (incolore) par l'ion thiosulfate $\ce{S2O3^2-}$ qui s'oxyde en tétrathionate $\ce{S4O6^2-}$.
\[
\ce{I2 + 2S2O3^2- -> 2I- + S4O6^2-}
\]
\cle{Repérage :} La coloration jaune/brune de $\ce{I2}$ disparaît. Pour plus de précision, on ajoute \textbf{l'empois d'amidon} (indicateur spécifique) lorsque la solution devient jaune paille : le complexe $\ce{I2}$-amidon (bleu intense) disparaît brutalement à l'équivalence.
\end{propriete}

\begin{savaistu}[Pourquoi l'empois d'amidon en fin de dosage ?]
L'amidon forme un complexe insoluble et très stable avec l'iode ($\ce{I2}$-amidon bleu). Si on l'ajoute au début, le complexe "piège" l'iode et ralentit la réaction avec le thiosulfate, retardant l'équivalence apparente ($V_E$ mesuré $>$ $V_E$ réel). On l'ajoute quand $[\ce{I2}]$ est faible (solution jaune pâle) pour un virage net et précis.
\end{savaistu}

\begin{attention}[Particularités du thiosulfate]
\begin{itemize}
\item La solution de $\ce{Na2S2O3}$ n'est \textbf{pas une solution étalon primaire} (elle se dégrade par oxydation air, décomposition acide, bactéries). Elle doit être \textbf{étalonnée} avant usage (souvent par un $\ce{KIO3}/\ce{KI}$ standard).
\item Le dosage se fait \textbf{sans acide ajouté} (le milieu doit rester neutre ou légèrement basique). En milieu acide, $\ce{S2O3^2-}$ se décompose : $\ce{S2O3^2- + 2H+ -> SO2 ^ + S v + H2O}$.
\end{itemize}
\end{attention}

\courssub{6. Méthodologie, exploitation des résultats et entraînement type Probatoire}

\courssub{Je vérifie mes connaissances}

\begin{exercice}[QCM : classification et dosages]
Pour chaque question, choisis la (ou les) bonne(s) réponse(s) en justifiant brièvement.
\begin{enumerate}
\item Le potentiel standard du couple $\ce{MnO4-}/\ce{Mn^2+}$ vaut $E^\circ = \SI{1,51}{V}$ et celui du couple $\ce{Fe^3+}/\ce{Fe^2+}$ vaut $E^\circ = \SI{0,77}{V}$. L'ion $\ce{MnO4-}$ en milieu acide peut donc oxyder :
\begin{enumerate}
\item les ions $\ce{Fe^3+}$ ;
\item les ions $\ce{Fe^2+}$ ;
\item aucun des deux.
\end{enumerate}
\item Dans un dosage manganimétrique, l'équivalence est repérée grâce à :
\begin{enumerate}
\item un indicateur coloré ajouté (phénolphtaléine) ;
\item la persistance de la coloration violette de $\ce{MnO4-}$ ;
\item l'apparition d'un précipité.
\end{enumerate}
\item Lors du dosage du diiode par les ions thiosulfate $\ce{S2O3^2-}$, le diiode $\ce{I2}$ joue le rôle de :
\begin{enumerate}
\item réducteur ;
\item oxydant ;
\item indicateur.
\end{enumerate}
\item Pour acidifier le milieu avant un dosage par $\ce{KMnO4}$, on utilise :
\begin{enumerate}
\item l'acide chlorhydrique ;
\item l'acide sulfurique ;
\item l'acide nitrique.
\end{enumerate}
\end{enumerate}
\end{exercice}

\begin{exercice}[Vrai ou faux, justifié]
Réponds par VRAI ou FAUX en justifiant chaque réponse par une phrase ou une demi-équation.
\begin{enumerate}
\item Dans le couple $\ce{S4O6^2-}/\ce{S2O3^2-}$, l'ion tétrathionate $\ce{S4O6^2-}$ est le réducteur.
\item Une réaction d'oxydoréduction est spontanée lorsque l'oxydant du couple de potentiel le plus élevé réagit avec le réducteur du couple de potentiel le plus faible (règle du gamma).
\item À l'équivalence d'un dosage d'oxydoréduction, les réactifs ont été mélangés dans les proportions stœchiométriques de l'équation de dosage.
\item Dans l'éthylotest, le passage de la couleur orange au vert prouve que les ions $\ce{Cr2O7^2-}$ ont été réduits en ions $\ce{Cr^3+}$.
\item On peut doser une solution de diiode par une solution de thiosulfate en versant le diiode dans la burette et le thiosulfate dans le bécher, sans conséquence sur le résultat.
\end{enumerate}
\end{exercice}

\begin{exercice}[Texte à trous]
Complète le texte suivant avec les mots ou expressions de la liste : \emph{oxydant ; réducteur ; équivalence ; violette ; décolorée ; stœchiométriques ; empois d'amidon ; burette ; oxydation ; réduction}.

\medskip
Dans un dosage manganimétrique, la solution de permanganate de potassium est placée dans la \ldots\ldots\ldots\ldots\ldots ; l'ion $\ce{MnO4-}$ joue le rôle d'\ldots\ldots\ldots\ldots\ldots et subit une \ldots\ldots\ldots\ldots\ldots. Tant que des ions $\ce{Fe^2+}$ sont présents dans le bécher, la coloration violette est \ldots\ldots\ldots\ldots\ldots. À l'\ldots\ldots\ldots\ldots\ldots, les réactifs ont été mélangés dans les proportions \ldots\ldots\ldots\ldots\ldots : la première goutte de $\ce{MnO4-}$ en excès donne au mélange une teinte \ldots\ldots\ldots\ldots\ldots persistante. En iodométrie, on peut ajouter de l'\ldots\ldots\ldots\ldots\ldots pour repérer plus précisément la disparition du diiode.
\end{exercice}

\begin{exercice}[Définitions]
Définis en une ou deux phrases chacun des termes suivants, en donnant un exemple tiré de la leçon :
\begin{enumerate}
\item potentiel standard d'un couple redox ;
\item dosage d'oxydoréduction ;
\item point d'équivalence ;
\item dosage colorimétrique (auto-indicateur).
\end{enumerate}
\end{exercice}

\courssub{Je m'entraîne}

\begin{exercice}[Demi-équations et prévision de réactions]
On donne les potentiels standard : $E^\circ(\ce{Cr2O7^2-}/\ce{Cr^3+}) = \SI{1,33}{V}$ ; $E^\circ(\ce{NO3-}/\ce{NO}) = \SI{0,96}{V}$ ; $E^\circ(\ce{Fe^3+}/\ce{Fe^2+}) = \SI{0,77}{V}$ ; $E^\circ(\ce{I2}/\ce{I-}) = \SI{0,54}{V}$ ; $E^\circ(\ce{S4O6^2-}/\ce{S2O3^2-}) = \SI{0,09}{V}$.
\begin{enumerate}
\item Écris les demi-équations électroniques des couples $\ce{Cr2O7^2-}/\ce{Cr^3+}$, $\ce{NO3-}/\ce{NO}$ et $\ce{S4O6^2-}/\ce{S2O3^2-}$ en milieu acide.
\item Place ces cinq couples sur une échelle verticale de potentiels standard.
\item À l'aide de la règle du gamma, prévois deux réactions spontanées possibles entre les espèces citées et écris leurs équations-bilan.
\end{enumerate}
\end{exercice}

\begin{exercice}[Dosage manganimétrique simple]
On dose $V_1 = \SI{20,0}{mL}$ d'une solution acidifiée d'ions $\ce{Fe^2+}$ par une solution de permanganate de potassium de concentration $C_2 = \SI{0,015}{mol.L^{-1}}$. L'équivalence est atteinte pour $V_E = \SI{13,3}{mL}$.
\begin{enumerate}
\item Écris l'équation-bilan de la réaction de dosage.
\item Définis l'équivalence et établis la relation entre les quantités de matière à l'équivalence.
\item Calcule la concentration molaire $C_1$ de la solution en ions $\ce{Fe^2+}$.
\item Calcule la concentration massique en fer de cette solution. On donne $M(\ce{Fe}) = \SI{55,8}{g.mol^{-1}}$.
\end{enumerate}
\end{exercice}

\begin{exercice}[Dosage d'une teinture d'iode]
Une teinture d'iode commerciale (solution de diiode dans l'éthanol ioduré, vendue en pharmacie à Yaoundé) est diluée 20 fois. On dose $V_1 = \SI{10,0}{mL}$ de la solution diluée par une solution de thiosulfate de sodium de concentration $C_2 = \SI{0,010}{mol.L^{-1}}$. L'équivalence est obtenue pour $V_E = \SI{9,2}{mL}$.
\begin{enumerate}
\item Écris les demi-équations des couples $\ce{I2}/\ce{I-}$ et $\ce{S4O6^2-}/\ce{S2O3^2-}$, puis l'équation-bilan du dosage.
\item Comment repère-t-on l'équivalence ? Cite un moyen de rendre ce repérage plus précis.
\item Calcule la concentration en diiode de la solution diluée, puis celle de la teinture commerciale.
\end{enumerate}
\end{exercice}

\begin{exercice}[Dispositif et protocole du dosage manganimétrique]
\begin{enumerate}
\item Schématise et légende le dispositif expérimental du dosage d'une solution de $\ce{Fe^2+}$ par une solution de $\ce{KMnO4}$ (burette, bécher, agitateur, solutions).
\item Décris le protocole en cinq étapes numérotées, du rinçage de la burette jusqu'à la lecture du volume équivalent.
\item Explique, en écrivant les demi-équations utiles, pourquoi on acidifie le milieu avec l'acide sulfurique $\ce{H2SO4}$ et non avec l'acide chlorhydrique $\ce{HCl}$. On donne $E^\circ(\ce{Cl2}/\ce{Cl-}) = \SI{1,36}{V}$ et $E^\circ(\ce{MnO4-}/\ce{Mn^2+}) = \SI{1,51}{V}$.
\end{enumerate}
\end{exercice}

\courssub{Je me prépare au Probatoire}

\begin{exercice}[Contrôle du fer dans un antianémique]
Pour lutter contre l'anémie, très fréquente chez les femmes enceintes suivies au Centre Hospitalier d'Essos, on administre des solutions de sulfate de fer(II). Un technicien de laboratoire veut vérifier la teneur en ions $\ce{Fe^2+}$ d'une solution pharmaceutique $S$. Il prélève $V_1 = \SI{25,0}{mL}$ de $S$, les acidifie avec de l'acide sulfurique dilué, puis les dose par une solution de permanganate de potassium de concentration $C_2 = \SI{0,020}{mol.L^{-1}}$. La coloration violette persiste à partir de $V_E = \SI{11,8}{mL}$.

On donne : $M(\ce{Fe}) = \SI{55,8}{g.mol^{-1}}$ ; $E^\circ(\ce{MnO4-}/\ce{Mn^2+}) = \SI{1,51}{V}$ ; $E^\circ(\ce{Fe^3+}/\ce{Fe^2+}) = \SI{0,77}{V}$.
\begin{enumerate}
\item Écris les deux demi-équations électroniques, puis montre que l'équation-bilan du dosage s'écrit :
\[ \ce{MnO4- + 5Fe^2+ + 8H+ -> Mn^2+ + 5Fe^3+ + 4H2O} \]
\item Justifie, à l'aide des potentiels standard, que cette réaction est spontanée.
\item Explique comment le technicien repère l'équivalence sans ajouter d'indicateur coloré.
\item Établis la relation à l'équivalence et calcule la concentration molaire $C_1$ en ions $\ce{Fe^2+}$ de la solution $S$.
\item Calcule la concentration massique en fer de $S$. L'étiquette annonce \SI{5,5}{g.L^{-1}} de fer : le flacon est-il conforme (écart toléré : $5\,\%$) ?
\item Pourquoi le dosage doit-il être réalisé rapidement après l'ouverture du flacon ? (Indice : penser au dioxygène de l'air, $E^\circ(\ce{O2}/\ce{H2O}) = \SI{1,23}{V}$.)
\end{enumerate}
\end{exercice}

\begin{exercice}[L'éthylotest et la sécurité routière]
Lors d'un contrôle sur l'axe lourd Yaoundé--Douala, un conducteur souffle dans un éthylotest contenant des cristaux de dichromate de potassium $\ce{K2Cr2O7}$ acidifiés, de couleur orange. En présence d'éthanol dans l'air expiré, les cristaux verdissent.

On donne les couples : $\ce{Cr2O7^2-}/\ce{Cr^3+}$ ($E^\circ = \SI{1,33}{V}$) et $\ce{CH3COOH}/\ce{CH3CH2OH}$ ($E^\circ = \SI{0,04}{V}$) ; $M(\ce{C2H5OH}) = \SI{46,0}{g.mol^{-1}}$.
\begin{enumerate}
\item Écris la demi-équation de réduction des ions dichromate $\ce{Cr2O7^2-}$ en ions chrome(III) $\ce{Cr^3+}$ en milieu acide.
\item Écris la demi-équation d'oxydation de l'éthanol $\ce{CH3CH2OH}$ en acide éthanoïque $\ce{CH3COOH}$.
\item Montre que l'équation-bilan de la réaction s'écrit :
\[ \ce{2Cr2O7^2- + 3CH3CH2OH + 16H+ -> 4Cr^3+ + 3CH3COOH + 11H2O} \]
\item Identifie l'oxydant et le réducteur, puis explique l'origine du changement de couleur observé.
\item Un éthylotest contient $n_0 = \SI{2,0e-4}{mol}$ d'ions $\ce{Cr2O7^2-}$. Calcule la quantité de matière, puis la masse d'éthanol que ce tube peut déceler au maximum.
\end{enumerate}
\end{exercice}

\begin{exercice}[Dosage des ions fer(II) par le dichromate de potassium]
Dans les eaux usées d'une unité de traitement de la SONARA à Limbé, on détermine la concentration en ions $\ce{Fe^2+}$ par dosage avec une solution de dichromate de potassium de concentration $C_2 = \SI{0,010}{mol.L^{-1}}$, en milieu acide sulfurique. On prélève $V_1 = \SI{25,0}{mL}$ de la solution à analyser ; l'équivalence est atteinte pour $V_E = \SI{15,0}{mL}$, repérée à l'aide d'un indicateur redox approprié.

On donne : $M(\ce{Fe}) = \SI{55,8}{g.mol^{-1}}$ ; $E^\circ(\ce{Cr2O7^2-}/\ce{Cr^3+}) = \SI{1,33}{V}$ ; $E^\circ(\ce{Fe^3+}/\ce{Fe^2+}) = \SI{0,77}{V}$.
\begin{enumerate}
\item Écris les demi-équations électroniques des deux couples mis en jeu.
\item Montre que l'équation-bilan du dosage s'écrit :
\[ \ce{Cr2O7^2- + 6Fe^2+ + 14H+ -> 2Cr^3+ + 6Fe^3+ + 7H2O} \]
\item Schématise et légende le dispositif de dosage.
\item Établis la relation à l'équivalence et calcule la concentration molaire $C_1$ en ions $\ce{Fe^2+}$.
\item Calcule la concentration massique en fer de l'effluent. La norme camerounaise de rejet fixe la limite à \SI{2,0}{g.L^{-1}} : l'effluent peut-il être rejeté sans traitement ?
\item Pourquoi ne peut-on pas repérer l'équivalence aussi simplement que dans un dosage manganimétrique ? Propose une solution.
\end{enumerate}
\end{exercice}

\newpage

\coursec{Corrigés des exercices}

\begin{corrige}[Exercice 1 — QCM : classification et dosages]
\textbf{Question 1 : réponse (b).} L'ion \ce{MnO4-} est l'oxydant du couple de potentiel le plus élevé ($\SI{1,51}{V} > \SI{0,77}{V}$). D'après la règle du gamma, il réagit avec le réducteur du couple de potentiel le plus faible, c'est-à-dire \ce{Fe^2+} : $\ce{MnO4- + 5Fe^2+ + 8H+ -> Mn^2+ + 5Fe^3+ + 4H2O}$. Il ne peut pas oxyder \ce{Fe^3+}, qui est déjà la forme oxydée de son couple.

\textbf{Question 2 : réponse (b).} Le permanganate est un \cle{auto-indicateur} : tant que \ce{Fe^2+} est en excès, chaque goutte de \ce{MnO4-} est décolorée ; à l'équivalence, la première goutte en excès donne une teinte violette pâle persistante. La phénolphtaléine est un indicateur acide-base, inutilisable ici.

\textbf{Question 3 : réponse (b).} Dans $\ce{I2 + 2S2O3^2- -> 2I- + S4O6^2-}$, le diiode capte des électrons (il est réduit en \ce{I-}) : c'est l'\textbf{oxydant}.

\textbf{Question 4 : réponse (b).} On utilise l'acide sulfurique. L'acide chlorhydrique serait oxydé en \ce{Cl2} toxique ($E^\circ(\ce{MnO4-}/\ce{Mn^2+}) = \SI{1,51}{V} > E^\circ(\ce{Cl2}/\ce{Cl-}) = \SI{1,36}{V}$) et l'acide nitrique est lui-même un oxydant parasite.
\end{corrige}

\begin{corrige}[Exercice 2 — Vrai ou faux, justifié]
\begin{enumerate}
\item \textbf{FAUX.} Dans le couple \ce{S4O6^2-}/\ce{S2O3^2-}, l'oxydant est \ce{S4O6^2-} (il capte des électrons : $\ce{S4O6^2- + 2e- -> 2S2O3^2-}$). Le réducteur est l'ion thiosulfate \ce{S2O3^2-}.
\item \textbf{VRAI.} C'est l'énoncé de la règle du gamma : la réaction spontanée met en jeu l'oxydant le plus fort (potentiel le plus élevé) et le réducteur le plus fort (potentiel le plus faible).
\item \textbf{VRAI.} C'est la définition même du point d'équivalence : $\dfrac{n_{\text{doseur}}}{a} = \dfrac{n_{\text{dosé}}}{b}$, où $a$ et $b$ sont les coefficients stœchiométriques de l'équation de dosage.
\item \textbf{VRAI.} L'ion dichromate \ce{Cr2O7^2-} est orange ; réduit par l'éthanol, il donne l'ion chrome(III) \ce{Cr^3+}, vert : $\ce{Cr2O7^2- + 14H+ + 6e- -> 2Cr^3+ + 7H2O}$.
\item \textbf{FAUX.} Le diiode est peu soluble et volatil ; le placer en burette exposerait le doseur à des pertes par évaporation et fausserait $V_E$. C'est le thiosulfate, solution aqueuse stable à court terme, qui va dans la burette, et le diiode dans le bécher.
\end{enumerate}
\end{corrige}

\begin{corrige}[Exercice 3 — Texte à trous]
Dans un dosage manganimétrique, la solution de permanganate de potassium est placée dans la \textbf{burette} ; l'ion \ce{MnO4-} joue le rôle d'\textbf{oxydant} et subit une \textbf{réduction}. Tant que des ions \ce{Fe^2+} sont présents dans le bécher, la coloration violette est \textbf{décolorée}. À l'\textbf{équivalence}, les réactifs ont été mélangés dans les proportions \textbf{stœchiométriques} : la première goutte de \ce{MnO4-} en excès donne au mélange une teinte \textbf{violette} persistante. En iodométrie, on peut ajouter de l'\textbf{empois d'amidon} pour repérer plus précisément la disparition du diiode.
\end{corrige}

\begin{corrige}[Exercice 4 — Définitions]
\begin{enumerate}
\item \textbf{Potentiel standard} d'un couple redox : tension, mesurée à \SI{25}{\celsius} dans les conditions standard (concentrations \SI{1}{mol.L^{-1}}, pression \SI{1}{bar}), entre une demi-pile constituée du couple et l'électrode standard à hydrogène. Exemple : $E^\circ(\ce{MnO4-}/\ce{Mn^2+}) = \SI{1,51}{V}$.
\item \textbf{Dosage d'oxydoréduction} : opération consistant à déterminer la concentration d'une espèce en la faisant réagir totalement et rapidement avec une solution étalonnée d'un partenaire redox, selon une réaction de stœchiométrie connue. Exemple : dosage de \ce{Fe^2+} par \ce{KMnO4}.
\item \textbf{Point d'équivalence} : état du dosage où le doseur et le dosé ont été introduits dans les proportions stœchiométriques de l'équation-bilan ; les deux réactifs sont alors totalement consommés.
\item \textbf{Dosage colorimétrique à auto-indicateur} : dosage où l'un des réactifs, intensément coloré, signale lui-même l'équivalence par l'apparition ou la disparition de sa couleur, sans indicateur ajouté. Exemple : persistance de la teinte violette de \ce{MnO4-} en manganimétrie.
\end{enumerate}
\end{corrige}

\begin{corrige}[Exercice 5 — Demi-équations et prévision de réactions]
\textbf{1.} Demi-équations en milieu acide :
\begin{align*}
\ce{Cr2O7^2- + 14H+ + 6e- &-> 2Cr^3+ + 7H2O} \\
\ce{NO3- + 4H+ + 3e- &-> NO + 2H2O} \\
\ce{S4O6^2- + 2e- &-> 2S2O3^2-}
\end{align*}

\textbf{2.} Échelle verticale (potentiels croissants de bas en haut) :
\begin{center}
\begin{tabular}{|c|l|}
\hline
\rowcolor{popblueL}
$E^\circ$ (\si{V}) & Couple \\ \hline
$+1,33$ & \ce{Cr2O7^2-}/\ce{Cr^3+} \\
$+0,96$ & \ce{NO3-}/\ce{NO} \\
$+0,77$ & \ce{Fe^3+}/\ce{Fe^2+} \\
$+0,54$ & \ce{I2}/\ce{I-} \\
$+0,09$ & \ce{S4O6^2-}/\ce{S2O3^2-} \\
\hline
\end{tabular}
\end{center}

\textbf{3.} Deux réactions spontanées (règle du gamma) :
\begin{itemize}
\item \ce{Cr2O7^2-} ($1,33$ V) oxyde \ce{Fe^2+} ($0,77$ V) : réduction $\ce{Cr2O7^2- + 14H+ + 6e- -> 2Cr^3+ + 7H2O}$ ; oxydation $\ce{Fe^2+ -> Fe^3+ + e-}$ ($\times 6$) :
\[ \ce{Cr2O7^2- + 6Fe^2+ + 14H+ -> 2Cr^3+ + 6Fe^3+ + 7H2O} \]
\item \ce{I2} ($0,54$ V) oxyde \ce{S2O3^2-} ($0,09$ V) : réduction $\ce{I2 + 2e- -> 2I-}$ ; oxydation $\ce{2S2O3^2- -> S4O6^2- + 2e-}$ :
\[ \ce{I2 + 2S2O3^2- -> 2I- + S4O6^2-} \]
\end{itemize}
(On pouvait aussi proposer \ce{NO3-} oxydant \ce{Fe^2+}, \ce{Cr2O7^2-} oxydant \ce{I-}, etc.)
\end{corrige}

\begin{corrige}[Exercice 6 — Dosage manganimétrique simple]
\textbf{1.} Équation-bilan :
\[ \ce{MnO4- + 5Fe^2+ + 8H+ -> Mn^2+ + 5Fe^3+ + 4H2O} \]

\textbf{2.} L'équivalence est atteinte lorsque les réactifs ont été mélangés dans les proportions stœchiométriques : 1 mol de \ce{MnO4-} pour 5 mol de \ce{Fe^2+}. D'où :
\[ \frac{n(\ce{Fe^2+})}{5} = \frac{n(\ce{MnO4-})}{1} \quad \Longrightarrow \quad \frac{C_1 V_1}{5} = C_2 V_E \]

\textbf{3.} Calcul de $C_1$ :
\[ C_1 = \frac{5\, C_2 V_E}{V_1} = \frac{5 \times \SI{0,015}{mol.L^{-1}} \times \SI{13,3}{mL}}{\SI{20,0}{mL}} = \SI{4,99e-2}{mol.L^{-1}} \approx \SI{0,050}{mol.L^{-1}} \]

\textbf{4.} Concentration massique en fer :
\[ t = C_1 \times M(\ce{Fe}) = \SI{4,99e-2}{mol.L^{-1}} \times \SI{55,8}{g.mol^{-1}} \approx \SI{2,8}{g.L^{-1}} \]
\end{corrige}

\begin{corrige}[Exercice 7 — Dosage d'une teinture d'iode]
\textbf{1.} Demi-équations et équation-bilan :
\begin{align*}
\ce{I2 + 2e- &-> 2I-} \\
\ce{2S2O3^2- &-> S4O6^2- + 2e-} \\ \hline
\ce{I2 + 2S2O3^2- &-> 2I- + S4O6^2-}
\end{align*}

\textbf{2.} L'équivalence est repérée par la \textbf{décoloration} de la solution (disparition de la teinte jaune-brune du diiode). Pour plus de précision, on ajoute de \textbf{l'empois d'amidon} lorsque la solution devient jaune paille : le complexe bleu \ce{I2}-amidon disparaît brutalement à l'équivalence.

\textbf{3.} À l'équivalence : $n(\ce{I2}) = \dfrac{n(\ce{S2O3^2-})}{2}$, soit $C_1 V_1 = \dfrac{C_2 V_E}{2}$.
\[ C_1 = \frac{C_2 V_E}{2 V_1} = \frac{\SI{0,010}{mol.L^{-1}} \times \SI{9,2}{mL}}{2 \times \SI{10,0}{mL}} = \SI{4,6e-3}{mol.L^{-1}} \]
La teinture commerciale est 20 fois plus concentrée :
\[ C_{\text{teinture}} = 20 \times C_1 = 20 \times \SI{4,6e-3}{mol.L^{-1}} = \SI{9,2e-2}{mol.L^{-1}} \]
\end{corrige}

\begin{corrige}[Exercice 8 — Dispositif et protocole du dosage manganimétrique]
\textbf{1.} Schéma légendé :
\begin{popfigure}
\begin{tikzpicture}[scale=0.85, >=Stealth]
% Table
\draw[thick, popmuted] (-0.5,0) -- (7.5,0);
% Burette
\draw[thick, popdark] (0.8,0.5) -- (0.8,6.5) -- (1.8,6.5) -- (1.8,0.5);
\draw[thick, popdark] (1.8,1.8) -- (2.3,1.8) -- (2.3,1.2) -- (1.8,1.2); % robinet
\fill[poppurple!40] (0.8,0.8) rectangle (1.8,5.8);
\node[poppurple!80!black, align=center, font=\small] at (1.3,6.9) {Burette : \ce{KMnO4} $C_2$};
\foreach \y in {1,2,3,4,5} \draw[thin, popdark] (0.6,\y) -- (0.8,\y) node[left, font=\tiny] {\y0};
% Becher
\draw[thick, popdark] (4.2,0.5) -- (4.2,4.2) -- (6.2,4.2) -- (6.2,0.5);
\draw[thick, popdark] (4.0,0.5) -- (6.4,0.5); % base
\fill[popgreen!15] (4.2,0.5) rectangle (6.2,3.2);
\node[popgreen!60!black, align=center, font=\small] at (5.2,4.5) {Bécher : \ce{Fe^2+} $V_1$ + \ce{H2SO4}};
% Stirrer
\draw[thick, popdark] (5.2,4.2) -- (5.2,5.2);
\draw[thick, decorate, decoration={coil, aspect=0.3, segment length=2mm, amplitude=1.5mm}] (5.2,3.8) -- (5.2,2.0);
\node[popdark, right, font=\small] at (5.2,4.8) {Agitateur magnétique};
% White tile
\fill[white, draw=popmuted, thick] (3.7,-0.4) rectangle (6.7,-0.1);
\node[font=\small, popmuted] at (5.2,-0.7) {Plaque blanche};
\end{tikzpicture}
\legende{Dispositif du dosage manganimétrique : burette contenant le doseur \ce{KMnO4}, bécher contenant le dosé \ce{Fe^2+} acidifié, agitation magnétique et plaque blanche pour la lecture.}
\end{popfigure}

\textbf{2.} Protocole :
\begin{enumerate}
\item Rincer la burette à l'eau distillée puis avec un peu de solution de \ce{KMnO4}, évacuer par le robinet.
\item Remplir la burette, chasser les bulles d'air, ajuster le ménisque sur le zéro (ou noter le volume initial $V_i$).
\item Pipeter $V_1$ de solution de \ce{Fe^2+} dans le bécher ; ajouter environ \SI{20}{mL} de \ce{H2SO4} dilué (\SI{1}{mol.L^{-1}}) ; placer sur l'agitateur avec la plaque blanche.
\item Verser le \ce{KMnO4} goutte à goutte sous agitation vigoureuse ; ralentir le débit lorsque la décoloration devient lente.
\item Arrêter dès la persistance d'une teinte violette pâle pendant \SI{30}{\second} ; noter $V_f$ et calculer $V_E = V_f - V_i$.
\end{enumerate}

\textbf{3.} Pourquoi \ce{H2SO4} et pas \ce{HCl} : les ions \ce{H+} sont indispensables à la réduction du permanganate ($\ce{MnO4- + 8H+ + 5e- -> Mn^2+ + 4H2O}$). Avec \ce{HCl}, les ions chlorure seraient oxydés car $E^\circ(\ce{MnO4-}/\ce{Mn^2+}) = \SI{1,51}{V} > E^\circ(\ce{Cl2}/\ce{Cl-}) = \SI{1,36}{V}$ :
\begin{align*}
\ce{2Cl- &-> Cl2 + 2e-} \quad (\times 5)\\
\ce{2MnO4- + 10Cl- + 16H+ &-> 2Mn^2+ + 5Cl2 ^ + 8H2O}
\end{align*}
Cette réaction parasite consomme du \ce{MnO4-} (le $V_E$ serait faussé) et dégage du dichlore toxique. L'acide sulfurique, lui, n'est pas oxydable par \ce{MnO4-} dans ces conditions.
\end{corrige}

\begin{corrige}[Exercice 9 — Contrôle du fer dans un antianémique]
\textbf{1.} Demi-équations :
\begin{align*}
\text{Réduction : } & \ce{MnO4- + 8H+ + 5e- -> Mn^2+ + 4H2O} \\
\text{Oxydation : } & \ce{Fe^2+ -> Fe^3+ + e-} \quad (\times 5)
\end{align*}
En additionnant : $\ce{MnO4- + 5Fe^2+ + 8H+ -> Mn^2+ + 5Fe^3+ + 4H2O}$. \hfill \emph{(2 pts)}

\textbf{2.} $E^\circ(\ce{MnO4-}/\ce{Mn^2+}) = \SI{1,51}{V} > E^\circ(\ce{Fe^3+}/\ce{Fe^2+}) = \SI{0,77}{V}$ : d'après la règle du gamma, l'oxydant du couple le plus élevé (\ce{MnO4-}) réagit spontanément avec le réducteur du couple le plus faible (\ce{Fe^2+}). \hfill \emph{(1 pt)}

\textbf{3.} Le permanganate est auto-indicateur : tant que \ce{Fe^2+} est en excès, chaque goutte violette est décolorée ; à l'équivalence, la première goutte de \ce{MnO4-} en excès n'est plus consommée et donne une teinte violette pâle persistante. \hfill \emph{(1 pt)}

\textbf{4.} Relation à l'équivalence : $\dfrac{n(\ce{Fe^2+})}{5} = n(\ce{MnO4-})$, soit $C_1 V_1 = 5\, C_2 V_E$.
\[ C_1 = \frac{5 \times \SI{0,020}{mol.L^{-1}} \times \SI{11,8}{mL}}{\SI{25,0}{mL}} = \SI{4,72e-2}{mol.L^{-1}} \]
\hfill \emph{(2 pts)}

\textbf{5.} Concentration massique :
\[ t = C_1 \times M(\ce{Fe}) = \SI{4,72e-2}{mol.L^{-1}} \times \SI{55,8}{g.mol^{-1}} = \SI{2,63}{g.L^{-1}} \]
Écart relatif avec l'étiquette (\SI{5,5}{g.L^{-1}}) : $\dfrac{5,5 - 2,63}{5,5} \times 100 \approx 52\,\%$, très supérieur aux $5\,\%$ tolérés : \textbf{le flacon n'est pas conforme} (probable oxydation d'une partie du fer(II) en fer(III)). \hfill \emph{(2 pts)}

\textbf{6.} Le dioxygène de l'air ($E^\circ(\ce{O2}/\ce{H2O}) = \SI{1,23}{V} > \SI{0,77}{V}$) oxyde lentement \ce{Fe^2+} en \ce{Fe^3+} : $\ce{O2 + 4Fe^2+ + 4H+ -> 4Fe^3+ + 2H2O}$. Un dosage tardif donnerait une concentration en \ce{Fe^2+} sous-estimée ; il faut donc doser rapidement après ouverture. \hfill \emph{(2 pts)}

\textbf{Points de vigilance :} ne pas oublier le facteur 5 dans la relation à l'équivalence ; convertir correctement les volumes (les mL se simplifient ici) ; comparer l'écart relatif au seuil de $5\,\%$ et non l'écart absolu.
\end{corrige}

\begin{corrige}[Exercice 10 — L'éthylotest et la sécurité routière]
\textbf{1.} Réduction du dichromate en milieu acide :
\[ \ce{Cr2O7^2- + 14H+ + 6e- -> 2Cr^3+ + 7H2O} \]
\hfill \emph{(2 pts)}

\textbf{2.} Oxydation de l'éthanol en acide éthanoïque (on équilibre O par \ce{H2O}, H par \ce{H+}, charges par \ce{e-}) :
\[ \ce{CH3CH2OH + H2O -> CH3COOH + 4H+ + 4e-} \]
\hfill \emph{(2 pts)}

\textbf{3.} PPCM de 6 et 4 électrons : on multiplie la réduction par 2 et l'oxydation par 3 :
\begin{align*}
\ce{2Cr2O7^2- + 28H+ + 12e- &-> 4Cr^3+ + 14H2O} \\
\ce{3CH3CH2OH + 3H2O &-> 3CH3COOH + 12H+ + 12e-}
\end{align*}
Somme et simplification ($28\ce{H+} - 12\ce{H+} = 16\ce{H+}$ ; $14\ce{H2O} - 3\ce{H2O} = 11\ce{H2O}$) :
\[ \ce{2Cr2O7^2- + 3CH3CH2OH + 16H+ -> 4Cr^3+ + 3CH3COOH + 11H2O} \]
\hfill \emph{(2 pts)}

\textbf{4.} L'oxydant est \ce{Cr2O7^2-} (il capte les électrons) ; le réducteur est l'éthanol \ce{CH3CH2OH}. Le changement de couleur vient du remplacement des ions dichromate \textbf{orange} par les ions chrome(III) \ce{Cr^3+} \textbf{verts} : le verdissement des cristaux prouve la présence d'éthanol dans l'air expiré. \hfill \emph{(2 pts)}

\textbf{5.} D'après la stœchiométrie, 2 mol de \ce{Cr2O7^2-} réagissent avec 3 mol d'éthanol :
\[ n(\ce{C2H5OH}) = \frac{3}{2}\, n_0 = \frac{3}{2} \times \SI{2,0e-4}{mol} = \SI{3,0e-4}{mol} \]
\[ m = n \times M = \SI{3,0e-4}{mol} \times \SI{46,0}{g.mol^{-1}} = \SI{1,38e-2}{g} = \SI{13,8}{mg} \]
Le tube peut déceler au maximum \SI{13,8}{mg} d'éthanol. \hfill \emph{(2 pts)}

\textbf{Points de vigilance :} bien équilibrer la demi-équation de l'éthanol (4 électrons, ajouter \ce{H2O} à gauche) ; simplifier les \ce{H+} et \ce{H2O} après addition ; respecter le rapport $3/2$ et non $2/3$.
\end{corrige}

\begin{corrige}[Exercice 11 — Dosage des ions fer(II) par le dichromate de potassium]
\textbf{1.} Demi-équations :
\begin{align*}
\ce{Cr2O7^2- + 14H+ + 6e- &-> 2Cr^3+ + 7H2O} \\
\ce{Fe^2+ &-> Fe^3+ + e-}
\end{align*}
\hfill \emph{(2 pts)}

\textbf{2.} On multiplie l'oxydation du fer par 6 pour égaliser les électrons, puis on additionne :
\[ \ce{Cr2O7^2- + 6Fe^2+ + 14H+ -> 2Cr^3+ + 6Fe^3+ + 7H2O} \]
\hfill \emph{(1 pt)}

\textbf{3.} Schéma légendé :
\begin{popfigure}
\begin{tikzpicture}[scale=0.85, >=Stealth]
% Table
\draw[thick, popmuted] (-0.5,0) -- (7.5,0);
% Burette
\draw[thick, popdark] (0.8,0.5) -- (0.8,6.5) -- (1.8,6.5) -- (1.8,0.5);
\draw[thick, popdark] (1.8,1.8) -- (2.3,1.8) -- (2.3,1.2) -- (1.8,1.2);
\fill[poporange!50] (0.8,0.8) rectangle (1.8,5.8);
\node[poporange!80!black, align=center, font=\small] at (1.3,6.9) {Burette : \ce{K2Cr2O7} $C_2$};
\foreach \y in {1,2,3,4,5} \draw[thin, popdark] (0.6,\y) -- (0.8,\y) node[left, font=\tiny] {\y0};
% Becher
\draw[thick, popdark] (4.2,0.5) -- (4.2,4.2) -- (6.2,4.2) -- (6.2,0.5);
\draw[thick, popdark] (4.0,0.5) -- (6.4,0.5);
\fill[popgreen!15] (4.2,0.5) rectangle (6.2,3.2);
\node[popgreen!60!black, align=center, font=\small] at (5.2,4.5) {Bécher : \ce{Fe^2+} $V_1$ + \ce{H2SO4} + indicateur};
% Stirrer
\draw[thick, popdark] (5.2,4.2) -- (5.2,5.2);
\draw[thick, decorate, decoration={coil, aspect=0.3, segment length=2mm, amplitude=1.5mm}] (5.2,3.8) -- (5.2,2.0);
\node[popdark, right, font=\small] at (5.2,4.8) {Agitateur magnétique};
% White tile
\fill[white, draw=popmuted, thick] (3.7,-0.4) rectangle (6.7,-0.1);
\node[font=\small, popmuted] at (5.2,-0.7) {Plaque blanche};
\end{tikzpicture}
\legende{Dosage de \ce{Fe^2+} par \ce{K2Cr2O7} : burette (doseur), bécher (dosé acidifié + indicateur redox), agitation.}
\end{popfigure}
\hfill \emph{(2 pts)}

\textbf{4.} À l'équivalence : $\dfrac{n(\ce{Fe^2+})}{6} = n(\ce{Cr2O7^2-})$, soit $C_1 V_1 = 6\, C_2 V_E$.
\[ C_1 = \frac{6 \times \SI{0,010}{mol.L^{-1}} \times \SI{15,0}{mL}}{\SI{25,0}{mL}} = \SI{3,6e-2}{mol.L^{-1}} \]
\hfill \emph{(2 pts)}

\textbf{5.} Concentration massique :
\[ t = C_1 \times M(\ce{Fe}) = \SI{3,6e-2}{mol.L^{-1}} \times \SI{55,8}{g.mol^{-1}} = \SI{2,01}{g.L^{-1}} \]
Cette valeur est \textbf{légèrement supérieure} à la norme de \SI{2,0}{g.L^{-1}} : l'effluent \textbf{ne peut pas être rejeté sans traitement}. \hfill \emph{(2 pts)}

\textbf{6.} Contrairement à \ce{MnO4-} (violet intense, détectable à très faible excès), le dichromate orange est moins coloré et les produits (\ce{Cr^3+} vert, \ce{Fe^3+} jaune) masquent le léger excès de doseur : le virage visuel direct est imprécis. Solution : ajouter un \textbf{indicateur redox} approprié (ex. diphénylamine sulfonée de baryum, virage incolore $\to$ violet) dont le potentiel de virage se situe dans le saut de potentiel à l'équivalence. \hfill \emph{(1 pt)}

\textbf{Points de vigilance :} le facteur stœchiométrique est 6 (et non 5 comme pour le permanganate) ; la comparaison à la norme doit se faire sur la concentration \emph{massique} ; justifier le choix de l'indicateur par le saut de potentiel.
\end{corrige}

\newpage

\coursec{Fiche bilan}

\begin{popfigure}
\centering
\begin{tikzpicture}[
  font=\small,
  central/.style={circle, draw=popdark, fill=popblue, text=white, font=\bfseries\small, minimum size=2.6cm, align=center},
  branche/.style={rounded corners=3pt, draw=#1, fill=#1!18, text width=3.1cm, align=center, inner sep=3pt, font=\scriptsize},
  lien/.style={thick, #1}
]
\node[central] (C) {Dosages\\d'oxydoréduction};
\node[branche=popblue]   (A) at (-5.2, 2.6) {\textbf{Classification}\\$E^\circ$ : pouvoir oxydant croissant\\$\ce{MnO4-/Mn^2+}$, $\ce{Cr2O7^2-/Cr^3+}$, $\ce{Fe^3+/Fe^2+}$, $\ce{I2/I-}$, $\ce{S4O6^2-/S2O3^2-}$};
\node[branche=poporange] (B) at (0, 3.4) {\textbf{Équations-bilan}\\2 demi-équations\\électrons équilibrés\\réactif limitant};
\node[branche=popgreen]  (Cc) at (5.2, 2.6) {\textbf{Principe du dosage}\\réaction rapide, totale, unique\\équivalence : réactifs en proportions stœchiométriques};
\node[branche=poppurple] (D) at (-5.2,-2.6) {\textbf{Manganimétrie}\\$\ce{MnO4- + 5Fe^2+ + 8H+ -> Mn^2+ + 5Fe^3+ + 4H2O}$\\$\ce{MnO4-}$ violet = indicateur};
\node[branche=popgold]   (E) at (0,-3.4) {\textbf{Équivalence}\\$n_{\text{ox}} = \dfrac{n_{\text{red}}}{k}$\\$C_1V_1$ et $C_2V_2$ liés par la stœchiométrie};
\node[branche=poppink]   (F) at (5.2,-2.6) {\textbf{Iodométrie}\\$\ce{I2 + 2S2O3^2- -> 2I- + S4O6^2-}$\\empois d'amidon : décoloration};
\draw[lien=popblue]   (C) -- (A);
\draw[lien=poporange] (C) -- (B);
\draw[lien=popgreen]  (C) -- (Cc);
\draw[lien=poppurple] (C) -- (D);
\draw[lien=popgold]   (C) -- (E);
\draw[lien=poppink]   (C) -- (F);
\end{tikzpicture}
\legende{Carte mentale de la leçon : généralisation et dosages d'oxydoréduction.}
\end{popfigure}

\begin{bilan}[L'essentiel de la leçon]
\textbf{1. Définitions clés}
\begin{itemize}
  \item \cle{Potentiel normal} $E^\circ$ (en volts) : grandeur qui classe les couples oxydant/réducteur ; plus $E^\circ$ est grand, plus l'\cle{oxydant} est fort (et plus le réducteur conjugué est faible).
  \item \cle{Dosage d'oxydoréduction} : détermination de la concentration inconnue d'une espèce grâce à une réaction d'oxydoréduction \textbf{rapide, totale et unique} avec un réactif titrant de concentration connue.
  \item \cle{Équivalence} : état où les réactifs ont été mélangés dans les \textbf{proportions stœchiométriques} ; elle se repère par un changement de couleur (réactif coloré ou indicateur).
\end{itemize}

\textbf{2. Couples et demi-équations à connaître par cœur (milieu acide)}
\begin{center}
\begin{tabularx}{\linewidth}{|l|X|c|}
\hline
\rowcolor{popblueL} \textbf{Couple} & \textbf{Demi-équation électronique} & $E^\circ$ (V) \\
\hline
$\ce{MnO4-}/\ce{Mn^2+}$ & $\ce{MnO4- + 8H+ + 5e- <=> Mn^2+ + 4H2O}$ & $+1,51$ \\
\hline
$\ce{Cr2O7^2-}/\ce{Cr^3+}$ & $\ce{Cr2O7^2- + 14H+ + 6e- <=> 2Cr^3+ + 7H2O}$ & $+1,33$ \\
\hline
$\ce{Cl2}/\ce{Cl-}$ & $\ce{Cl2 + 2e- <=> 2Cl-}$ & $+1,36$ \\
\hline
$\ce{NO3-}/\ce{NO}$ & $\ce{NO3- + 4H+ + 3e- <=> NO + 2H2O}$ & $+0,96$ \\
\hline
$\ce{Fe^3+}/\ce{Fe^2+}$ & $\ce{Fe^3+ + e- <=> Fe^2+}$ & $+0,77$ \\
\hline
$\ce{I2}/\ce{I-}$ & $\ce{I2 + 2e- <=> 2I-}$ & $+0,54$ \\
\hline
$\ce{S4O6^2-}/\ce{S2O3^2-}$ & $\ce{S4O6^2- + 2e- <=> 2S2O3^2-}$ & $+0,09$ \\
\hline
$\ce{SO4^2-}/\ce{SO2}$ & $\ce{SO4^2- + 4H+ + 2e- <=> SO2 + 2H2O}$ & $+0,17$ \\
\hline
\end{tabularx}
\end{center}

\textbf{3. Équations-types des dosages au programme}
\[ \ce{MnO4- + 5Fe^2+ + 8H+ -> Mn^2+ + 5Fe^3+ + 4H2O} \qquad \text{(manganimétrie)} \]
\[ \ce{I2 + 2S2O3^2- -> 2I- + S4O6^2-} \qquad \text{(iodométrie)} \]

\textbf{4. Formules et relations d'équivalence}
\begin{center}
\begin{tabularx}{\linewidth}{|l|X|}
\hline
\rowcolor{poporangeL} \textbf{Grandeur} & \textbf{Relation} \\
\hline
Quantité de matière & $n = C \times V$ \quad (mol $=$ mol.L$^{-1}\times$ L) \\
\hline
Équivalence (général) & $\dfrac{n_{\text{ox}}}{a} = \dfrac{n_{\text{red}}}{b}$ \quad ($a$, $b$ : coefficients stœchiométriques) \\
\hline
Dosage $\ce{Fe^2+}$ par $\ce{MnO4-}$ & $n(\ce{Fe^2+}) = 5\,n(\ce{MnO4-})$ \; soit \; $C_{\text{red}}V_{\text{red}} = 5\,C_{\text{ox}}V_{E}$ \\
\hline
Dosage $\ce{I2}$ par $\ce{S2O3^2-}$ & $n(\ce{S2O3^2-}) = 2\,n(\ce{I2})$ \; soit \; $C_{\text{thio}}V_{E} = 2\,C_{\ce{I2}}V_{\ce{I2}}$ \\
\hline
Concentration massique & $C_m = C \times M$ \quad (g.L$^{-1}$) \\
\hline
\end{tabularx}
\end{center}

\textbf{5. Méthodes express}
\begin{itemize}
  \item \textbf{Écrire une équation-bilan} : équilibrer chaque demi-équation (O avec \ce{H2O}, H avec \ce{H+}, charges avec $e^-$) $\rightarrow$ multiplier pour égaliser les électrons $\rightarrow$ additionner et simplifier.
  \item \textbf{Prévoir une réaction} : l'oxydant du couple de plus grand $E^\circ$ réagit avec le réducteur du couple de plus petit $E^\circ$ (règle du gamma).
  \item \textbf{Repérer l'équivalence} : manganimétrie = première teinte rose persistante ($\ce{MnO4-}$ en léger excès) ; iodométrie = décoloration du bleu de l'empois d'amidon.
  \item \textbf{Exploiter} : tableau d'avancement $\rightarrow$ relation stœchiométrique à l'équivalence $\rightarrow$ $C$ inconnue $\rightarrow$ éventuellement $C_m = C \times M$.
\end{itemize}
\end{bilan}

\begin{aretenir}[Checklist avant le Probatoire]
\begin{itemize}
  \item[$\square$] Je sais classer les couples au programme sur un axe de $E^\circ$ croissant.
  \item[$\square$] Je sais écrire sans hésiter les demi-équations de $\ce{MnO4-}/\ce{Mn^2+}$, $\ce{Cr2O7^2-}/\ce{Cr^3+}$, $\ce{Fe^3+}/\ce{Fe^2+}$, $\ce{I2}/\ce{I-}$ et $\ce{S4O6^2-}/\ce{S2O3^2-}$.
  \item[$\square$] Je sais équilibrer une équation-bilan en égalisant les électrons échangés.
  \item[$\square$] Je sais prédire le sens d'une réaction spontanée avec la règle du gamma.
  \item[$\square$] Je connais les trois qualités d'une réaction de dosage : rapide, totale, unique.
  \item[$\square$] Je sais schématiser le montage : burette (titrant) + erlenmeyer (solution à doser) + agitateur.
  \item[$\square$] Je sais que $\ce{MnO4-}$ violet sert d'indicateur : équivalence = teinte rose persistante.
  \item[$\square$] Je sais que l'empois d'amidon donne un bleu foncé avec $\ce{I2}$ : équivalence = décoloration.
  \item[$\square$] Je sais écrire la relation d'équivalence $n(\ce{Fe^2+}) = 5\,n(\ce{MnO4-})$ et $n(\ce{S2O3^2-}) = 2\,n(\ce{I2})$.
  \item[$\square$] Je convertis toujours les volumes en litres avant d'appliquer $n = CV$.
  \item[$\square$] Je sais passer de la concentration molaire à la concentration massique ($C_m = C \times M$).
  \item[$\square$] Je vérifie l'homogénéité des unités et je donne le résultat avec le bon nombre de chiffres significatifs.
\end{itemize}
\end{aretenir}

\findecours

\end{document}
